\documentclass[11pt,a4paper]{article}
\usepackage[margin=27mm]{geometry}
\usepackage{fontspec}
\setmainfont{Latin Modern Roman}
\usepackage{amsmath,amssymb,amsthm,mathtools}
\usepackage{mathrsfs}
\usepackage{longtable,booktabs,array,enumitem,microtype}
\usepackage{xurl}
\usepackage[unicode,colorlinks=true,linkcolor=blue,citecolor=blue,urlcolor=blue]{hyperref}
\usepackage{bookmark}
\newtheorem{theorem}{Statement}
\newtheorem{proposition}[theorem]{Statement}
\newtheorem{lemma}[theorem]{Statement}
\newtheorem{corollary}[theorem]{Statement}
\newtheorem{definition}[theorem]{Statement}
\newtheorem{remark}[theorem]{Statement}
\newtheorem{obstruction}[theorem]{Statement}
\newcommand{\Gal}{\operatorname{Gal}}
\newcommand{\Cent}{\operatorname{Cent}}
\newcommand{\Norm}{\operatorname{Norm}}
\newcommand{\Aut}{\operatorname{Aut}}
\newcommand{\Hom}{\operatorname{Hom}}
\newcommand{\im}{\operatorname{im}}
\newcommand{\Q}{\mathbb Q}
\newcommand{\Z}{\mathbb Z}
\newcommand{\F}{\mathbb F}
\newcommand{\Gm}{\mathbb G_m}
\newcommand{\Zhat}{\widehat{\mathbb Z}}
\newcommand{\Galois}{\Gamma}
\newcommand{\Gzero}{\Gamma^{\circ}}
\newcommand{\Cyc}{K(\mu_\infty)}
\newcommand{\sep}{\mathrm{sep}}
\newcommand{\ab}{\mathrm{ab}}
\newcommand{\tors}{\mathrm{tors}}
\newcommand{\et}{\mathrm{\acute et}}
\newcommand{\status}[1]{\textbf{#1.}}
\newcommand{\passage}[1]{\par\smallskip\noindent\textit{Manuscript: #1}\par\smallskip}
\newcommand{\resref}[1]{\hyperref[res:#1]{#1}}
\newcommand{\Rtwo}{\resref{A6}}
\newcommand{\Arith}{\resref{A1}}
\makeatletter
\renewcommand\paragraph{\@startsection{paragraph}{4}{\z@}%
  {2.7ex \@plus 0.7ex \@minus 0.2ex}{0.5ex \@plus 0.2ex}%
  {\normalfont\normalsize\bfseries}}
\makeatother
\setlength{\emergencystretch}{2em}
\title{The cyclotomic kernel and cuspidal centralizers\\\large Proof audit and calculations for cote 149}
\author{AI-produced by GPT-6 Astra\\OpenAI, ChatGPT Pro mode}
\date{12 November 2026}
\begin{document}
\maketitle
\begin{abstract}
We prove the cyclotomic-kernel centralizer theorem with two cusp-uniqueness arguments, audit the arithmetic and birational comparison maps, and give explicit covering and nilpotent-quotient calculations. Corrections to the manuscript are distinguished from transcription evidence. The coverage ledger records the remaining obligations, including the full local-field calculation and the unrestricted action question. This AI-produced proof audit makes no claim of priority or of a proof of a section conjecture.
\end{abstract}

% BEGIN sections/conventions.tex
We fix an algebraic closure $\overline K$ of a characteristic-zero field $K$. All fundamental groups and homomorphisms of them are profinite and continuous. A geometric base point is fixed; changing it conjugates the displayed subgroups by the geometric group. A section is an actual continuous right inverse of the arithmetic projection; its class means its orbit under conjugation by the geometric group.

\small
\begin{longtable}{>{\raggedright\arraybackslash}p{.16\textwidth}>{\raggedright\arraybackslash}p{.28\textwidth}>{\raggedright\arraybackslash}p{.45\textwidth}}
\toprule Our symbol & Manuscript symbol & Meaning and principal scans\\\midrule\endhead
$\Gamma_K$ & $\Gamma$, $\Gamma_{\overline K/K}$ & Absolute Galois group; scans 45, 53.\\
$\Gamma_K^\circ$ & $\Gamma^\circ$, $\Gamma^0$ & Kernel of $\chi:\Gamma_K\to\Zhat^\times$; scans 45, 55, 61.\\
$E_U$ & $E_U$, $\mathcal E_U$ & Arithmetic fundamental group; scans 53 and 126--127.\\
$\pi_U$ & $\Pi_U$, $\pi$, $\pi_{U/K}$ & Geometric fundamental group; scans 49, 53, 127.\\
$I_c$ & $L_i$, $I_V$ & Inertia at a selected lift of a cusp or place; scans 53--55, 127.\\
$D_c$ & $Z(L_i)$, $N_V$ & Decomposition group; the centralizer/normalizer reading at scan 54 is audited below.\\
$\Zhat(1)$ & $T$, $T_\infty(\overline K^*)$ & Full Tate module of roots of unity; scans 53--54.\\
$C(s)$ & $\pi^{\Gamma^\circ}$, $\mathrm{Zentr}(\Gamma^\circ)$ & $\Cent_{\pi_U}(s(\Gamma_K^\circ))$; scans 55--57.\\
$F$ & corps cyclotomique & $K(\mu_\infty)$; scans 61--62.\\
$X'$ & $X'=\widehat{\widetilde U}$ & Compatible geometric points of completed cover levels; scans 59--61.\\
$E_L$, $P_L$ & $\mathcal E_L$, $\pi_{L/K}$ & Function-field Galois group and geometric kernel; scans 126--127.\\
$P$, $N_r$ & Quotients tested in LM3b D.5 & Pro-$5$ group and its class-$r$ quotient, $P/\gamma_{r+1}P$; Part E.\\
\bottomrule
\end{longtable}
\normalsize

% END sections/conventions.tex

\clearpage
\tableofcontents
\section{Scope and claim ledger}
Statement identifiers A1--E4 and AE1--AE4 are stable identifiers for the accompanying results table. An Open statement gives its proved portion and an explicit remaining obligation. The report distinguishes these from requested work still marked ``unfinished''. None of the claims labelled as unresolved is used as a hypothesis of a proof. The supplied research prompt, prior AI drafts and handwritten scans are evidence to audit, rather than mathematical authorities.
\section{A. Arithmetic and recognition of inertia}

% BEGIN sections/A_agent.tex
% Integrated in Part A of the master paper.

We use genuine continuous actions after fixing a geometric base point and an actual continuous section $s:\Gamma_K\to E_U$. Changing the section by conjugation with $a\in\pi$ conjugates every centralizer and every lifted cusp below by $a$. Every subgroup is closed, every homomorphism between Tate modules is continuous and $\widehat{\mathbb Z}$-linear, every abelianization is profinite, and $\langle z\rangle$ denotes the closed subgroup generated by $z$. Put $F=K(\mu_\infty)$ and $\Gamma_K^\circ=\operatorname{Gal}(\overline K/F)$. For an abelian variety $A$, write $T(A)=\prod_\ell T_\ell(A)$, with its usual continuous Galois action. ``Divisible element'' means an element of $\bigcap_{n\geq1}nA(L)$, or of $\bigcap_{n\geq1}(L^*)^n$ in multiplicative notation.

\subsection{The arithmetic input}
\paragraph{Statement A1 (Proved; Part A(a)).}\label{res:A1}
\passage{Section 3, manuscript pp. 26–27, scans 54–56}

This is the arithmetic input for cote 149, section 3, manuscript pp.~26--27, scans 54--56.
Let $K$ be a number field and $F'/K(\mu_\infty)$ finite inside $\overline K$. For every abelian variety $A/F'$ and every torus $S/F'$,
\[
 A(F')_{\rm tors}\text{ is finite},\qquad T(A)^{\Gamma_{F'}}=0,
 \qquad\operatorname{Hom}_{\Gamma_{F'}}(T(S),T(A))=0,
 \qquad\bigcap_{n\geq1}nA(F')=0.
\]
Moreover $\bigcap_{n\geq1}(F'^*)^n=\mu_\infty$.

\emph{Source input.} Ribet's appendix to Katz--Lang, Theorem 1, has the following data: $k$ is a number field, $A/k$ is an abelian variety, and $k^{\rm cyc}=k(\mu_\infty)$. Its assertion is that $A(k^{\rm cyc})_{\rm tors}$ is finite \cite[Appendix, Theorem 1, p.~315]{Ribet1981}. We use precisely this assertion, including finiteness across all primes.

\emph{Proof.} A finite set of generators of $F'/F$ and a finite set of coefficients defining $A$ lie in a number field $k_1\subset F'$ containing $K$. Enlarging $k_1$ inside $F'$ as needed gives the equality $F'=k_1(\mu_\infty)$ and an abelian variety $A_1/k_1$ with $A=A_1\otimes_{k_1}F'$. Thus Ribet applies to $A_1$ and proves the first assertion. A nonzero element of $T_\ell(A)^{\Gamma_{F'}}$ provides compatible $F'$-rational torsion points of unbounded $\ell$-power order. This contradicts that finiteness. Taking the product over $\ell$ proves the second assertion. Choose a finite extension $F''/F'$ splitting $S$. The same argument applies to $F''$. On $\Gamma_{F''}$ the module $T(S)$ is a direct sum of copies of $\widehat{\mathbb Z}(1)$ with trivial action, because $F''$ contains all roots of unity. A Galois-equivariant homomorphism $T(S)\to T(A)$ therefore has image in $T(A)^{\Gamma_{F''}}=0$. Restriction changes neither the underlying homomorphism nor its image, proving the third assertion.

For the assertion about divisible points, take $P\in\bigcap_n nA(F')$ and enlarge $k_1$ so that $P$ is defined over $k_1$. The finite nonempty sets of $n!$-division points of $P$ in $A(F')$ have a nonempty inverse limit: any finite collection of compatibility requirements is satisfied by a sufficiently high factorial division point. Choose a compatible system and let
\[
 c:\Gamma_{k_1}\longrightarrow T(A),\qquad
 c(\sigma)=(\sigma(P_{n!})-P_{n!})_n
\]
be its continuous Kummer cocycle. Its restriction to the normal subgroup $H=\Gamma_{F'}$ is zero. For $h\in H$ and $g\in\Gamma_{k_1}$, comparing $c(hg)$ and $c(g(g^{-1}hg))$ gives $h c(g)=c(g)$. Consequently $c(g)\in T(A)^H=0$ for every $g$. All the chosen division points are defined over $k_1$, so $P\in\bigcap_n nA_1(k_1)$.

Here this last intersection can be checked without a finite-generation assumption. Spread $A_1$ to an abelian scheme at one finite place $v$ of good reduction. Properness puts $A_1(k_1)$ inside $A_1(\mathcal O_{k_{1,v}})$. Reduction modulo $\mathfrak m_v^r$ is a homomorphism to a finite group for every $r$. An element divisible by every positive integer has zero image in each of these finite groups: use multiplication by the group's order. The reductions separate points, by separatedness and completeness. Hence $P=0$. This also proves directly the needed injectivity of the Kummer map $A_1(k_1)\to H^1(\Gamma_{k_1},T(A))$. In cohomological terms, the restriction map used above has kernel $H^1(\operatorname{Gal}(F'/k_1),T(A)^H)=0$; the cocycle calculation supplies its proof with profinite coefficients.

For the multiplicative assertion, the inclusion $\mu_\infty\subset\bigcap_n(F'^*)^n$ follows by taking roots of roots of unity. Let $a$ belong to the intersection and enlarge $k_1$ so that $a\in k_1^*$. A compatible root system in $F'$ gives a Kummer cocycle which is inflated from
\[
 Q=\operatorname{Gal}(k_1(\mu_\infty)/k_1)
 \subset\widehat{\mathbb Z}^{\,*}
\]
with coefficients $T=\widehat{\mathbb Z}(1)$. This subgroup is open: its index is the finite degree of $k_1\cap\mathbb Q(\mu_\infty)$ over $\mathbb Q$. Choose $m$ such that $Q$ contains the kernel of reduction modulo $m$. There is $t\in Q$ such that $t-1=du$, where $d$ is a positive integer and $u\in\widehat{\mathbb Z}^{\,*}$. Explicitly, at primes dividing $m$ take $t_\ell=1+\ell^{v_\ell(m)}$; at odd primes not dividing $m$ take $t_\ell=2$; at $2$ not dividing $m$ take $t_2=3$. All nonunit factors of $t-1$ then have finite support and finite valuation.

For a continuous cocycle $f:Q\to T$, commutativity of $Q$ gives
\[
 (t-1)f(q)=(q-1)f(t).
\]
The right side is the coboundary of $f(t)$; therefore $d[f]=0$ in $H^1(Q,T)$. The Kummer class of $a^d$ over $k_1$ is zero. The map $k_1^*\to H^1(\Gamma_{k_1},T)$ is injective. Indeed, an element with roots of every order in $k_1$ has valuation zero at every finite prime. All its roots in $k_1$ are algebraic units whose complex absolute values are uniformly bounded. There are only finitely many algebraic integers of bounded degree with all conjugates bounded, since their monic minimal polynomials have bounded integral coefficients. Repetition of roots of unbounded order shows that the element is a root of unity. Finally use an order divisible by the exponent of the finite group $\mu(k_1)$ to show that the element is $1$. Thus $a^d=1$, as required.\hfill$\square$

\paragraph{Statement A2 (Proved; the geometric input).}\label{res:A2}
\passage{Section 3, manuscript pp. 26–27, scans 53–56}

This supplies the local groups used in cote 149, section 3, manuscript pp.~26--27, scans 53--56.
Let $U$ be an affine hyperbolic curve over an algebraically closed field of characteristic zero, and let $\mathcal C$ be the profinite set of cusps of its universal cover. Every $I_c$ is isomorphic to $\widehat{\mathbb Z}(1)$, and
\[
 I_c\cap I_{c'}\ne1\Longrightarrow c=c',\qquad
 N_\pi(I_c)=I_c,\qquad C_\pi(z)=I_c\quad(1\ne z\in I_c).
\]
For a curve over a characteristic-zero field, $N_{E_U}(I_c)=D_c$, its cusp decomposition group. The same normalizer identity holds with any nontrivial closed subgroup of $I_c$ in place of $I_c$. For a connected finite geometric covering with group $H\subset\pi$, compactification $X_H$, and cusp set $S_H$, the continuous quotient map induced by $U_H\hookrightarrow X_H$ fits into
\[
 \widehat{\mathbb Z}(1)\xrightarrow{\rm diagonal}
 \bigoplus_{d\in S_H}\widehat{\mathbb Z}(1)e_d
 \longrightarrow H^{\rm ab}\longrightarrow T(J_{X_H})\longrightarrow0.\tag{A.1}
\]
The image of the $d$-summand is the image of its inertia group.

\emph{Sources and proof.} The assertions with their characteristic-zero, full-profinite, hyperbolic hypotheses are Stix 2008, Lemma 10, Proposition 11, Corollaries 12--13 and (3.1), and Stix 2012, Lemma 14, Proposition 16(1)--(3) and (3.1) \cite{Stix2008,Stix2012}. We spell out the local argument that is used repeatedly. The boundary map in the localization sequence identifies a positive loop with the corresponding summand in (A.1); its only relation is the diagonal one. Once there are at least two cusps the map from any individual summand is injective, proving the required local inertia injection. Arbitrarily many cusps can be obtained in finite covers: a cusp inertia subgroup has infinite index, since an open procyclic subgroup would be the fundamental group of a finite hyperbolic cover, whose abelianization has rank at least two. The number of points over that cusp in the finite normal cover associated to $N$ is $[\pi:NI_c]$, which is unbounded as $N$ shrinks.

Two different lifted cusps can be separated at a finite level. Refine that level until there are at least three cusps. A nonzero common inertia element has a nonzero positive integral power in the subgroup defining this level, and (A.1) would put that nonzero element in two different summand lines. Those lines intersect trivially when there are at least three summands. This proves the intersection assertion. An element normalizing $I_c$, or centralizing a nonidentity element of it, conjugates $I_c$ to an inertia group meeting $I_c$ nontrivially. It therefore fixes $c$. Its geometric stabilizer is $I_c$; its arithmetic stabilizer is $D_c$. For a nontrivial subgroup $J\subset I_c$ the same argument gives $N_{E_U}(J)\subset D_c$. Conversely the cyclotomic action of $D_c$ preserves every closed subgroup of the procyclic group $I_c$, so equality holds. These arguments also prove that every abelian subgroup meeting $I_c$ nontrivially is contained in $I_c$.\hfill$\square$

\subsection{A criterion for an inertia element}
\paragraph{Statement A3 (Proved; Part A(b)).}\label{res:A3}
\passage{Section 3, manuscript pp. 27, scans 55–56}

This supplies the inertia recognition needed for cote 149, section 3, manuscript p.~27, scans 55--56.
With the geometric data of A2, let $z\in\pi$. For every open normal subgroup $N\triangleleft\pi$ put $H_N=N\langle z\rangle$. Suppose that the image of $z$ in $H_N^{\rm ab}$ belongs to the subgroup generated by the images of the cusp inertia groups of $U_{H_N}$. Then $z\in I_c$ for some lifted cusp $c$.

\emph{Proof.} The assertion includes $z=1$, for which any cusp works. Let $z\ne1$. The quotient $Q_N=H_N/N$ is a finite cyclic group generated by $zN$. The assumed abelianized assertion, composed with the continuous surjection $H_N^{\rm ab}\to Q_N$, says that the cusp inertia subgroups generate $Q_N$. For each prime $\ell$, the $\ell$-Sylow subgroup of a finite cyclic group generated by subgroups is contained in one of those subgroups: the largest $\ell$-adic order among the finitely many generators equals the order of the Sylow subgroup. Thus the $\ell$-primary component $z_\ell$ of $z$ fixes a cusp in the finite normal cover $U_N$. The component is defined using the product decomposition of the procyclic group $\langle z\rangle$ into its Sylow subgroups.

For all $N$ the finite sets of cusps fixed by $z_\ell$ are nonempty. For any finitely many levels choose a common normal refinement; a fixed cusp there projects to fixed cusps at all those levels. Compactness consequently gives a cusp of the universal cover fixed by $z_\ell$, so $z_\ell\in I_{c_\ell}$. Select one nontrivial component $z_{\ell_0}$. Every other component commutes with it, and A2 gives $z_\ell\in I_{c_{\ell_0}}$. Finite products of the primary components are dense in $\langle z\rangle$. Since $I_{c_{\ell_0}}$ is closed, it contains $z$.\hfill$\square$

\emph{The exact comparison with Nakamura and Stix.} In Nakamura 1994, \S2.1, the data preceding (2.1.1) are a number field $k$, a smooth geometrically connected noncomplete curve of type $(g,n)$ with $n\ge1$, $2-2g-n<0$, and a full class $\mathcal D$ of finite groups, closed under subgroups, quotients and extensions. The geometric group is the pro-$\mathcal D$ completion $G$ of the displayed surface-group presentation. His (2.1.1) says that a nonidentity $z\in G$ is an inertia element exactly when there is a closed subgroup of its arithmetic normalizer mapping openly to $\Gamma_k$ and acting on $\langle z\rangle$ through the cyclotomic character. Lemma (2.1.2) says that the centralizer of a nonidentity element in the subgroup of a free generator of a free pro-$\mathcal D$ group is that generator subgroup. Corollary (2.1.3) excludes an abelian subgroup meeting two distinct free-generator subgroups nontrivially. For $Z$ the union of all conjugates of the cusp subgroups, Lemma (2.1.4) gives
\[
 z\in Z\quad\Longleftrightarrow\quad
 z\in\overline{[H,H]H^\ell\langle Z\cap H\rangle}
 \quad\text{for all open }H\ni z\text{ and all }\ell\in|\mathcal D|.
\]
These are precise restatements of (2.1.1)--(2.1.4), pp.~94--97 \cite{Nakamura1994}.

For the full profinite group, A3 implies the right side of (2.1.4) directly: for a given $H$ choose open normal $N\subset H$ and use the inclusion $N\langle z\rangle\hookrightarrow H$; its induced map on abelianizations carries boundary summands to boundary summands. Conversely an inertia element satisfies A3 by the same inclusion of its inertia subgroup at each level. Thus A3 and Nakamura's full-profinite (2.1.4) give the same recognition of $z$, through the specified inclusion and quotient maps.

Stix 2008 Theorem 20 has a characteristic-zero field $k$ satisfying
$\operatorname{Hom}_{\Gamma_k}(T(S),T(A))=0$ for every $k$-torus $S$ and abelian variety $A$, a hyperbolic $U/k$, and a nontrivial closed $I\subset\pi$. It equates containment in inertia with procyclicity and essential cyclotomic normalization, and with procyclicity and cyclotomic normalization by the entire normalizer. The 2012 manuscript numbers its refined statement Theorem 30: fix a prime $\ell$ for which every $\Gamma_k$-map $\mathbb Z_\ell(1)\to T_\ell(A)$ is zero; for nontrivial closed pro-$\ell$ $I$, containment in inertia is equivalent to $I\simeq\mathbb Z_\ell$ with essential, or full, cyclotomic normalization. Lemma 28 proves stability of that hypothesis under finite extensions by Weil restriction. Applying this statement to a nonzero $\ell$-component of a procyclic $I$ and then applying A2 to the commuting remaining components proves the 2008 assertion whenever its hypothesis holds. For number fields and the full class of finite groups it gives Nakamura (2.1.1) with $I=\langle z\rangle$. Nakamura also treats the quotient $\pi\twoheadrightarrow\pi^{\mathcal D}$ for other full classes; that scope is stated in his theorem itself, rather than inferred by an unproved passage through this quotient.

The phrase ``number field'' in Stix 2008 Lemma 19 is left as printed; A1 supplies the necessary assertion for $F'$ directly. Stix 2012 Lemma 27 explicitly permits algebraic $k/\mathbb Q$ and has the stronger prime-specific hypothesis $[kK_\ell:K_\ell]<\infty$, with $K_\ell$ the fixed field of an $\ell$-Sylow subgroup of $\Gamma_{\mathbb Q^{\rm ab}}$. In particular it applies to finite extensions of $K(\mu_\infty)$.

\subsection{The reduction from centralizers to inertia}
\paragraph{Statement A4 (Proved; Part A(c), R2 steps (1)--(3)).}\label{res:A4}
\passage{Section 2, manuscript pp. 24, scans 49; Section 3, manuscript pp. 27, scans 55–56}

The manuscript passages are cote 149, section 2, manuscript p.~24, scan 49, and section 3, manuscript p.~27, scans 55--56.
For a number field $K$, a section $s$ and $C=C_\pi(s(\Gamma_K^\circ))$, the centralizer is trivial when $U$ is proper. When $U$ is affine hyperbolic, every nonidentity $z\in C$ belongs to an inertia subgroup $I_c$; the subgroup $s(\Gamma_K^\circ)$ fixes $c$, and $I_c\subset C$.

\emph{Proof.} Let $H\subset\pi$ be open and contain $z\in C$. A continuous action on a topologically finitely generated profinite group has open stabilizer of any open subgroup: choose an open characteristic $N\subset H$ and use the finite action on $\pi/N$. Such characteristic subgroups exist because there are finitely many subgroups of any given finite index, so the intersection of all subgroups of index at most a fixed bound is characteristic and open. Thus some open $\Gamma_{K'}\subset\Gamma_K$ stabilizes $H$ through $s$. The subgroup $H s(\Gamma_{K'})$ of $E_U$ defines a finite etale covering of $U_{K'}$ with geometric group $H$, and $s|_{\Gamma_{K'}}$ is a section for that covering. The inclusion $H\hookrightarrow\pi$ is the induced continuous fundamental-group homomorphism, with the chosen base points.

The image of $z$ in $H^{\rm ab}$ is fixed by
$\Gamma_{K'(\mu_\infty)}=\Gamma_{K'}\cap\Gamma_K^\circ$. The map in (A.1) to the Tate module of the compactified covering is Galois-equivariant. A1 therefore makes the image of $z$ there zero. If $U$ is proper, choose an open characteristic $N$ with $z\notin N$ and set $H=N\langle z\rangle$. Then $H^{\rm ab}=T(J_{X_H})$, so $z$ has zero abelianized image. But the homomorphism $H\to H/N$ factors through $H^{\rm ab}$ and takes $z$ to a nonidentity generator, a contradiction. For affine hyperbolic $U$, exactness of (A.1), applied to all $N\langle z\rangle$, is the hypothesis of A3. Hence $z\in I_c$.

For $\sigma\in\Gamma_K^\circ$, the conjugate $s(\sigma)I_cs(\sigma)^{-1}=I_{s(\sigma)c}$ contains the same nonidentity element $z$. A2 gives $s(\sigma)c=c$. The action of $D_c$ on $I_c=\widehat{\mathbb Z}(1)$ is the cyclotomic one, trivial on this subgroup; hence $I_c\subset C$. For proper genus zero, $\pi=1$. For proper genus one, $\pi=T(J_X)$ and inner conjugation is trivial, so A1 applies directly. This covers those two proper types as well.\hfill$\square$

\paragraph{Statement A5 (Proved; Part A(d), direct route).}\label{res:A5}
\passage{Section 3, manuscript pp. 27, scans 55–56}

This repairs the proof required by cote 149, section 3, manuscript p.~27, scans 55--56.
For an affine hyperbolic geometric curve, any closed subgroup $B\subset\pi$ all of whose nonidentity elements are inertia elements is contained in a single inertia subgroup. That subgroup is unique when $B\ne1$.

\emph{Proof.} Assume $B\ne1$ and write $\mathcal T$ for the cusps $c$ with $B\cap I_c\ne1$. It is a nonempty $B$-stable set. If $B$ is not contained in one inertia subgroup, every $J_c=B\cap I_c$ is proper. The set $\mathcal T$ cannot be one $B$-orbit. Otherwise $B$ would be the union of the conjugates of a single proper closed subgroup $J_c$. A proper closed subgroup has proper image in some finite quotient $Q$ of $B$. No finite group is the union of conjugates of a proper subgroup $J$: in its action on $Q/J$, counting pairs consisting of an element and a fixed coset gives average fixed-coset count $1$. The identity fixes $[Q:J]>1$ cosets, so some element fixes none. This contradicts the proposed union.

Choose $c,c'\in\mathcal T$ in distinct $B$-orbits. There is an open normal $N\subset\pi$ such that their images at level $N$ are in distinct $B$-orbits. Indeed, failure at every level would give by compactness one element of $B$ carrying $c$ to $c'$. Set $H=BN$. This is open, contains $B$, and has distinct images of $c,c'$ among its cusps. If it has at least three cusps, keep it. If it has exactly two, hyperbolicity forces the compactification $X_H$ to have genus at least one. Every element of $B$ maps to $1$ under the continuous compactification map $H\twoheadrightarrow\pi_1(X_H)$: every nonidentity element of $B$ lies in an inertia subgroup of $H$. Choose a degree-two finite etale cover of $X_H$, using the quotient $\pi_1(X_H)\twoheadrightarrow\mathbb Z/2\mathbb Z$ supplied by $T(J_{X_H})$. Pullback to $U_H$ replaces $H$ by an open subgroup still containing $B$ and gives four cusps. The two specified images remain distinct. This proves the required three-cusp assertion without assuming that arbitrary distinct cusps can be separated by a subgroup containing $B$.

Choose $1\ne z\in B\cap I_c$ and $1\ne z'\in B\cap I_{c'}$. In (A.1) for this $H$, let their classes be $a e_c$ and $b e_{c'}$, with $a,b\in\widehat{\mathbb Z}(1)$ nonzero. Each of $zz'$ and $z^2z'$ is a nonidentity element of $B$, because its boundary image is nonzero. It must therefore lie in an inertia group. If there are at least four cusps, neither $a e_c+b e_{c'}$ nor $2a e_c+b e_{c'}$ belongs to a single summand line modulo the diagonal: a fourth coordinate forces the diagonal coefficient to vanish. If there are exactly three cusps, membership of $n a e_c+b e_{c'}$ in a single line forces that line to be the third one and forces $na=b$. Applying this for $n=1,2$ gives $a=b$ and $2a=b$, whence $a=0$, a contradiction. This proof is valid over $\widehat{\mathbb Z}$ even though that ring has zero divisors: subtraction gives $a=0$ directly. Thus $B\subset I_c$; uniqueness follows from A2.\hfill$\square$

The program \texttt{code/A\_boundary\_checks.py} verifies the derangement count in $S_2,\ldots,S_6$ and the three-cusp calculation modulo $3,5,7,9,15,25$; its exact output is attached. Modulo $2$ both products can lie in inertia lines, which records explicitly why a single finite quotient cannot replace the integral argument.

The original R2 step (4) omits the distinction between cusps and their $C$-orbits, and does not justify the existence of its covering containing $C$. Its written proof is therefore incomplete. A5 repairs both omissions and uses three cusps with the two products $zz'$ and $z^2z'$; four cusps are unnecessary. Hyperbolicity enters through A2 and, in the two-cusp stage, through the positive genus of $X_H$.

\paragraph{Statement A6 (Proved; Part A(d), literature route; Part A(e); R2).}\label{res:A6}
\passage{Section 3, manuscript pp. 27, scans 55–56}

This answers the formula (3) of cote 149, section 3, manuscript p.~27, scan 55, and the marginal precisions on scan 56.
For a hyperbolic curve $U/K$ over a number field and a continuous section $s:\Gamma_K\to E_U$,
\[
 C_\pi(s(\Gamma_K^\circ))\ne1
 \quad\Longleftrightarrow\quad s\text{ is cuspidal}.
\]
In this event there is a unique lifted cusp $c$ fixed by $s(\Gamma_K)$ and
$C_\pi(s(\Gamma_K^\circ))=I_c$. For a proper hyperbolic curve the centralizer is $1$.

\emph{First proof.} A4 and A5 give $C\subset I_c$ and $I_c\subset C$ whenever $C\ne1$. Since $\Gamma_K^\circ$ is normal in $\Gamma_K$, its centralizer is preserved by $s(\Gamma_K)$. Thus $s(\Gamma_K)$ normalizes $I_c$, and A2 puts it in $D_c$. This is the definition of a cuspidal section, and it proves that the base cusp is $K$-rational. Conversely a cuspidal section acts on its nontrivial inertia group through $\chi$, so $I_c\subset C$. A5 and A2 force equality and uniqueness. The proper assertion is A4.\hfill$\square$

\emph{Second proof, using the opened literature.} Stix 2008 Theorem 14 assumes injectivity of points into section classes for proper genus-at-least-two curves over every finite extension of the base field. Under that hypothesis it proves cusp uniqueness, injection of the splitting classes of a cusp decomposition extension, and disjointness from finite-point sections and from other cusps. In the 2012 manuscript these are Theorem 17(1)--(3); Remark 18(1) supplies absence of divisible abelian-variety points as a sufficient hypothesis. We verify it over $F$ with A1. Explicitly, on a proper curve $Y/F'$ two points $y_1,y_2$ with equal section classes give equal classes after the morphism $y\mapsto[y-y_2]$ to $J_Y$. Its Kummer cocycle is the difference of the two point-section cocycles, as in Stix 2008 Proposition 44. Therefore $[y_1-y_2]$ is divisible by every integer and A1 makes it zero. A principal divisor $y_1-y_2$ with distinct points would define a degree-one map to $\mathbb P^1$, so this is impossible when the genus is positive.

Base change $U_F\to U$ induces a continuous homomorphism $\pi_1(U_F)\to E_U$ identifying $\pi_1(U_F)$ with the inverse image of $\Gamma_F$. This follows by descending every finite etale cover to a finite intermediate field and taking the inverse limit of the usual finite-base-change subgroup identifications. Now restrict $s$ to $\Gamma_F$. A4 makes each nonidentity element of $C$ an element of an inertia group of a cusp fixed by this restricted section. Theorem 14(1), together with the final argument of its proof applied to every finite covering and finite extension of $F$, gives uniqueness of the lifted cusp, rather than merely uniqueness of its image on $U$. Consequently all elements of $C$ lie in the same $I_c$. The previous normality argument finishes the proof. The finite-cover step here is exactly where hyperbolicity is needed in Stix's route: one takes neighborhoods whose compactifications have genus at least two. Neither proof uses surjectivity in a section conjecture.\hfill$\square$

\paragraph{Statement A7 (Proved; germs and larger subgroups, Part A(e)).}\label{res:A7}
\passage{Section 3, manuscript pp. 30, scans 61}

This addresses the subgroup clause of cote 149, section 3, manuscript p.~30, scan 61.
A6 holds for every germ of section, after replacing $K$ by its finite field of definition. For a section over $K$, let $\Omega$ be any open subgroup of $\Gamma_K^\circ$. Then its centralizer in $\pi$ equals $C_\pi(s(\Gamma_K^\circ))$. More generally, let $\Delta\subset\Gamma_K$ be closed and contain such an $\Omega$. For a noncuspidal $s$ one has $C_\pi(s(\Delta))=1$. For a cuspidal $s$ at its unique lifted cusp $c$,
\[
 C_\pi(s(\Delta))=I_c^{\chi(\Delta)}
 =\prod_{\ell:\,\chi_\ell|_\Delta=1}\mathbb Z_\ell(1)
 \ \subset\ I_c=\prod_\ell\mathbb Z_\ell(1).\tag{A.2}
\]
The factors at all other primes are zero.

\emph{Proof.} The fixed field of $\Omega$ is finite over $F$ and hence is $K_1(\mu_\infty)$ for some number field $K_1\supset K$. A6 applied over $K_1$ proves the assertion for the restricted section. Cuspidality descends to $K$: choose a finite Galois extension $K_2/K$ containing $K_1$. Its cyclotomic-kernel subgroup is normal in $\Gamma_K$, and its centralizer is the same inertia group, since it contains the nonzero inertia group obtained over $K_1$ and A6 gives uniqueness. Therefore $s(\Gamma_K)$ normalizes that inertia group. This proves both descent and equality of centralizers for $\Omega$. The containment $\Omega\subset\Delta$ gives a reversed containment of centralizers. In the cuspidal case it places $C_\pi(s(\Delta))$ inside $I_c$, where the action is exactly $\chi$. For $a_\ell\in\mathbb Z_\ell(1)$, invariance is the equation $(\chi_\ell(\delta)-1)a_\ell=0$ for every $\delta\in\Delta$. The ring $\mathbb Z_\ell$ is an integral domain. The fixed subgroup is all of the factor when the character is trivial and is zero otherwise. Taking products proves (A.2).\hfill$\square$

\paragraph{Statement A8 (Proved; exceptional types and the cyclotomic assertion).}\label{res:A8}
\passage{Section 3, manuscript pp. 26, scans 54}

These are the exceptional types in cote 149, section 3, manuscript p.~26, scan 54.
The cyclotomic centralizer for a proper genus-one curve is zero whenever a section exists. For geometric types $(0,0)$ and $(0,1)$ it is zero because $\pi=1$. For a form of type $(0,2)$, after a finite extension splitting the two cusps, the geometric group is $\widehat{\mathbb Z}(1)$ and its entire group centralizes every section restricted to the cyclotomic kernel. Both cusps have that same inertia group. For every hyperbolic type, including $(0,3)$ and $(1,1)$, A6 applies with the full geometric fundamental group.

\emph{Proof.} The first three assertions follow from the explicit geometric groups and A1. For $\mathbb G_m$ the Kummer covers $t\mapsto t^n$ identify the geometric group with $\widehat{\mathbb Z}(1)$; the genuine action of any section on this abelian group is the same cyclotomic action, since an inner automorphism of an abelian group is trivial. The maps from the punctured local rings at $0$ and $\infty$ induce respectively a generator and its inverse, so both inertia images are the entire group. The last assertion is A6 with its stated hyperbolicity hypothesis.\hfill$\square$

For the real base field the exact calculation is
$\Gamma_{\mathbb R}^{\circ}=1$ and therefore
$C(s)=\pi_1(U_{\mathbb C})$ for every actual section over $\mathbb R$.
Indeed $\Gamma_{\mathbb R}=\{1,c\}$ and $c$ acts on a primitive fourth
root of unity by inversion, so $\chi(c)=-1\ne1$; its cyclotomic kernel
is trivial. The centralizer of its image is the entire geometric group.
For example this gives the whole free profinite rank-two group for
the point $2$ of $U_{0,3}(\mathbb R)$, as well as for a real tangent at
$0$. The finite character calculation modulo $4$ and the nonabelian
$S_3$ quotient are recorded in \texttt{code/root\_math\_checks.py}.
At a real place of a number field, the decomposition homomorphism
$\Gamma_{\mathbb R}\to\Gamma_K$ sends $c$ to an element with the same
cyclotomic value $-1$, so that decomposition group intersects
$\Gamma_K^\circ$ trivially. This specifies the comparison with the
global calculation A6. Over $\mathbb R$, a positive-dimensional
abelian variety has $T(A)^{\Gamma_{\mathbb R}^\circ}=T(A)\ne0$;
the arithmetic vanishing input A1 thus does not hold for this base.

% END sections/A_agent.tex


% BEGIN sections/A_root.tex
\subsection{Explicit loops, Fermat covers, and the Ihara series}
\paragraph{Statement A9 (Proved; Part A(f)).}\label{res:A9}
\passage{Section 3, manuscript pp. 27, scans 55–56; La Longue Marche, section 34, manuscript p. 245, scan 483}

For $U=\mathbb P^1_{\mathbb Q}\setminus\{0,1,\infty\}$ choose the rational tangential base point at $0$ directed toward $1$, the corresponding section $s$, and geometric generators $x,y$ around $0,1$. Then
\[
 C_\pi(s(\Gamma_{\mathbb Q}^{\circ}))=\overline{\langle x\rangle},
 \qquad y\notin C_\pi(s(\Gamma_{\mathbb Q}^{\circ})),
 \qquad [x,y]\notin C_\pi(s(\Gamma_{\mathbb Q}^{\circ})).
\]
These assertions concern actual elements and the genuine action defined by $s$. They test the formula in cote 149, section 3, manuscript p.~27, scans 55--56, and the characterizations in Longue Marche, section 34, scan 483.

\begin{proof}
The equality is A6. The image of $y$ in $\pi^{\rm ab}=\widehat{\mathbb Z}(1)^2$ is $(0,1)$, whereas the image of $\overline{\langle x\rangle}$ is the first summand. This proves the assertion for $y$ with the quotient map specified. We give an independent arithmetic obstruction for the commutator.

Let $H_n$ be the kernel of $\pi\to(\mathbb Z/n)^2$ sending $x,y$ to the coordinate generators. Its geometric covering is
\[
 w^n+v^n=1,\qquad z=w^n.
\]
It is finite etale over $U$, and its compactification is the Fermat curve $X_n$. The maps
\[
 \overline{[\pi,\pi]}\longrightarrow H_n\longrightarrow H_n^{\rm ab}
 \longrightarrow T(J_{X_n})
 \tag{A.3}
\]
are continuous and are equivariant for the restricted section over $\mathbb Q(\mu_\infty)$. A finite field extension supplies any chosen compatible lift and does not affect the invariant-vanishing input A1. The first arrow is an inclusion, the second is abelianization, and the last is induced by the open immersion into the compactification. In particular the last two arrows kill the second derived subgroup of $\pi$.

Here is a direct certificate that $[x,y]$ has nonzero image for $n=3$. The abelian covering of the topological wedge of the two generators has vertices $(a,b)\in(\mathbb Z/3)^2$ and directed edges $e_x(a,b),e_y(a,b)$ incrementing one coordinate. Its commutator cycle is
\[
 e_x(0,0)+e_y(1,0)-e_x(0,1)-e_y(0,0).
\]
Compactification kills the cycles around $x^3$, $y^3$ and $(xy)^3$, at all their lifts. Fix a primitive cube root $\omega$ and give $e_x(a,b)$ weight $\omega^{a+b}$ and $e_y(a,b)$ weight zero. This functional vanishes on all the boundary cycles: the relevant sums are $1+\omega+\omega^2$ and $1+\omega^2+\omega^4$. Its value on the commutator is $1-\omega\ne0$. The calculation works over $\mathbb Z_3[\omega]$, so it proves nonvanishing in $T_3(J_{X_3})$, through the comparison of topological finite covers with etale finite covers. The exact source for that comparison is SGA~1, Expose XII, Theorem 5.1: finite etale covers of a complex scheme locally of finite type correspond to its finite analytic etale covers \cite{SGA1}. The map on first homology is precisely the pro-$3$ abelianization of the displayed cover.

An element fixed by $s(\Gamma_{\mathbb Q}^{\circ})$ would have fixed image in this Tate module, contrary to A1. This proves the commutator assertion without using its conclusion as an input to the Fermat calculation.

To identify the requested Ihara power series, use the free pro-$\ell$ quotient $P$ of $\pi$, $M=P'/P''$, and $\Lambda=\mathbb Z_\ell[[T_1,T_2]]$. The opened primary research source is Kodani's thesis, section 4.3, Lemma 4.3.1 and equation (4.3.2), which restates Ihara's Theorem 2(i) and proves the module assertion independently \cite{Kodani2017,Ihara1991}. Ihara's own opened ICM paper, section 6.4, formula (6.4.1), gives the corresponding multiplier $B'_\sigma$ and the Fermat-Jacobian connection on pp.~116--117. In the thesis's tangential convention, $M$ is free of rank one over $\Lambda$, generated by the commutator, and
\[
 \rho(g)([a,b])=F_g(T_1,T_2)[a,b]\quad\text{in }M.
 \tag{A.4}
\]
The action on $\Lambda$ sends $1+T_i$ to $(1+T_i)^{\chi_\ell(g)}$. The paragraph following Theorem 4.3.3 identifies the resulting module with the inverse limit of the Tate modules of the Fermat Jacobians; its transition maps are induced by the power maps between Fermat curves. Equation (A.3) is the finite-level compactification map in that construction. The exact hypotheses here are $\{0,1,\infty\}\subset\mathbb P^1_{\mathbb Q}$, a fixed prime $\ell$, the pro-$\ell$ geometric group, and the specified tangential Galois action.

We also check the change of tangential convention. Applying $z\mapsto1/z$ to the convention with tangent at infinity takes it to the present tangent at zero. In the old generators $a,b$, the new generators are $x=(ab)^{-1}$ and $y=b$. The identity
\[
 [x,y]=x[a,b]^{-1}x^{-1}
\]
shows that their commutator bases differ by the invertible scalar $-(ab)^{-1}$ in the group ring. On $\Gamma_{\mathbb Q}^{\circ}$ the scalar ring has trivial action, so the multiplier $F_g$ in (A.4) is unchanged. Taking $\ell=3$, the preceding nonzero Fermat image and A1 show that $F_g\ne1$ for at least one $g\in\Gamma_{\mathbb Q}^{\circ}$. Otherwise (A.4) would fix the commutator and hence its nonzero image under (A.3). This supplies the requested power-series obstruction with its actual module and comparison maps; no numerically guessed Galois matrix is used.
\end{proof}

The executed program \texttt{code/root\_math\_checks.py} computes the incidence and boundary matrices for $n=2,3,4,5$. Their compact first-homology ranks are $0,2,6,12$; the commutator raises the boundary rank exactly for $n=3,4,5$. The degree-three calculation above is the integral proof behind these checks. The CM elliptic test, including the image of the full centralizer, is given in E1--E3 below.

% END sections/A_root.tex

\section{B. Arithmetic centralizers}

% BEGIN sections/B_agent.tex
\subsection{An elementary proof of the arithmetic centralizer assertion}

\paragraph{Statement B1 (Proved: R3).}\label{res:B1}
\passage{section 2; p.22; scan 45}

Fix an algebraic closure $\overline{\mathbf Q}$, with its profinite absolute
Galois group $G_{\mathbf Q}$. For every number field
$k\subset\overline{\mathbf Q}$, put $F=k(\mu_\infty)$. Then
\[
 C_{G_{\mathbf Q}}(G_F)=1.
\]
Here $G_F=\operatorname{Gal}(\overline{\mathbf Q}/F)$ is a closed subgroup,
and the centralizer is taken in the full profinite group, without passing to
an abelian or pro-$\ell$ quotient. This proves the starred assertion in
cote 149, \S2, manuscript p.~22, scan 45.

\begin{proof}
Replacing $k$ by its finite Galois closure $k_0/\mathbf Q$ gives the inclusion
$G_{k_0(\mu_\infty)}\subseteq G_F$, and hence the inclusion of centralizers
in the opposite direction. It therefore suffices to consider $k/\mathbf Q$
Galois. Fix $\sigma\in C_{G_{\mathbf Q}}(G_F)$.

For $a\in F$ and $h\in G_F$ we have
$h(\sigma a)=\sigma(ha)=\sigma a$. Applying the same argument to
$\sigma^{-1}$ proves $\sigma(F)=F$. For arbitrary
$x\in\overline{\mathbf Q}$ the identities
\[
 h(x)=x\quad\Longleftrightarrow\quad h(\sigma x)=\sigma x
 \qquad(h\in G_F)
\]
identify the two closed stabilizers in $G_F$. By the Galois correspondence,
\begin{equation}\label{eq:B-stabilizer}
 F(x)=F(\sigma x).
\end{equation}

We record explicitly the choice of split prime needed below. A nonconstant
polynomial $f\in\mathbf Z[T]$ has values divisible by infinitely many
rational primes. Indeed, after translating $T$ we have $f(0)=c\ne0$.
For any finite set $P$ of primes, enlarge $P$ to contain the prime divisors
of $c$ and put $b=\prod_{p\in P}p$. Then
$f(cbt)/c$ is an integer congruent to $1$ modulo every prime in $P$;
for some integer $t$ its absolute value exceeds $1$. It has a prime divisor
outside $P$. Apply this to the minimal polynomial of an integral primitive
element of a finite Galois number field $k'/\mathbf Q$. Outside the finitely
many primes dividing its discriminant or the index of its generated order,
a root modulo $p$ gives a residue-degree-one prime of $k'$. All residue
degrees are equal in an unramified Galois extension; thus $p$ splits
completely. This proves the existence of completely split primes outside
any prescribed finite set, without an appeal to a density theorem.

Suppose that $\sigma$ acts nontrivially on a finite Galois subfield
$k'\subset F$. Choose a rational prime $p$ which splits completely in $k'$
and is unramified in $k$. The action of
$\operatorname{Gal}(k'/\mathbf Q)$ on the primes above $p$ is simply
transitive, so $\sigma$ moves each such prime. Choose one, $\mathfrak p$.
The Chinese remainder theorem, applied modulo $\mathfrak p^2$ and the
other primes above $p$, gives $a\in k'^*$ such that
\[
 v_{\mathfrak p}(a)=1,\qquad
 v_{\mathfrak q}(a)=0\quad(\mathfrak q\mid p,\ \mathfrak q\ne\mathfrak p).
\]
The normalization here is $v(p)=1$. Extend $v_{\mathfrak p}$ to $F$.
Its value group is exactly
\begin{equation}\label{eq:B-value-group}
 v(F^*)=\frac1{p-1}\mathbf Z[1/p].
\end{equation}
To check the equality, roots of unity of order prime to $p$ belong to the
maximal unramified extension of $\mathbf Q_p$: they lift the roots in finite
residue fields by the simple-root form of Hensel's lemma. The extensions
$\mathbf Q_p(\mu_{p^r})/\mathbf Q_p$ are totally ramified of degree
$(p-1)p^{r-1}$; this degree and total ramification follow from the
Eisenstein polynomial $\Phi_{p^r}(1+T)$.
Since $k/\mathbf Q$ is unramified at $p$, its completions do not add a
ramification factor. The value groups of the finite composita are therefore
contained in, and the $p$-power cyclotomic fields already realize, the union
in \eqref{eq:B-value-group}. In particular this argument requires
unramifiedness in $k$, not merely complete splitting in $k'$.

Choose a prime $\ell\nmid p(p-1)$. Equation \eqref{eq:B-value-group} gives
$1\notin\ell v(F^*)$, so $a\notin(F^*)^\ell$. Choose
$\alpha^\ell=a$. The extension $F(\alpha)/F$ is cyclic of degree $\ell$,
because $F$ contains $\mu_\ell$. Equation \eqref{eq:B-stabilizer} gives
$F(\alpha)=F(\sigma\alpha)$. If a generator $\tau$ sends
$\alpha$ to $\zeta_\ell\alpha$, then
$\tau(\sigma\alpha)/(\sigma\alpha)\in\mu_\ell$.
The $\zeta_\ell^j$-eigenspace for $\tau$ in
$F(\alpha)=\bigoplus_{i=0}^{\ell-1}F\alpha^i$ is $F\alpha^j$.
Consequently
\[
 \sigma\alpha=c\alpha^j,\qquad
 \sigma(a)=c^\ell a^j,
 \qquad c\in F^*,\quad 1\le j\le\ell-1.
\]
The condition $j\ne0\bmod\ell$ follows from $\sigma\alpha\notin F$.
At our chosen valuation, $v(\sigma a)=0$ whereas $v(a)=1$. Thus
$0=\ell v(c)+j$. Since
$\ell v(F^*)\cap\mathbf Z=\ell\mathbf Z$, this contradicts
$1\le j\le\ell-1$.

We have proved that $\sigma$ fixes every finite Galois subfield of $F$;
their union is $F$, since $F/\mathbf Q$ is Galois. Finally, for an arbitrary
finite Galois number field $k''$ containing $k$, the inclusion
$G_{k''(\mu_\infty)}\subset G_F$ shows that $\sigma$ centralizes that
subgroup too. The argument just given, with $k''$ in place of $k$, proves
that $\sigma$ fixes $k''(\mu_\infty)$ pointwise. These fields exhaust
$\overline{\mathbf Q}$, so $\sigma=1$.
\end{proof}



\paragraph{Audit of the supplied R3 argument.}
Every implication in its proposed chain is valid after the preceding two
additions: the prime $p$ is unramified in the original finite coefficient
field, and the chosen Kummer class is nonzero. In the final passage one may
use finite Galois number fields $k''$ containing $k$; it is unnecessary to
assert that an arbitrary finite extension of $F$ is Galois over
$\mathbf Q$. It is true that every finite $F'/F$ has the form
$k_1(\mu_\infty)$ for a number field $k_1$: choose finitely many algebraic
generators of $F'/F$ and put them together with $k$ into $k_1$.

\subsection{A general Hilbertian criterion and the fields in Part B(c)}

\paragraph{Statement B2 (Proved: a stronger arithmetic criterion).}\label{res:B2}
\passage{section 2; p.22; scan 45}

For a Hilbertian field $k$ and a separable closure $k^{\rm sep}$, every
nontrivial closed normal subgroup $H\triangleleft G_k$ satisfies
$C_{G_k}(H)=1$. Consequently, for every algebraic Galois extension
$M/k$ in $k^{\rm sep}$,
\[
 C_{G_k}(G_M)=
 \begin{cases}
 1,&M\ne k^{\rm sep},\\
 G_k,&M=k^{\rm sep}.
 \end{cases}
\]
The comparison is in the category of profinite groups with continuous
homomorphisms.

\begin{proof}
We use two precisely identified inputs from the opened primary texts.
Jarden \cite[Proposition 2.5]{B_Jarden2011} treats a Hilbertian field $k$
and an algebraic Galois extension $N/k$ that is not separably closed. Its
conclusion is that $G_N$ is almost locally free: for each integer $e\ge1$,
almost every $e$-tuple, with respect to Haar measure, generates a free
profinite group of rank $e$ (definition (2) in the introduction).
Haran--Jarden \cite[Corollary 2.5]{B_HaranJarden1991} states that the
separable closure of a Hilbertian field cannot be the compositum of two
proper Galois extensions of the field. Their Theorem 2.4 gives the stronger
input: the compositum of two Galois extensions, neither contained in the
other, is Hilbertian. These statements include infinite Galois extensions.

There is no nontrivial closed abelian normal subgroup $A$ of $G_k$.
Otherwise its fixed field is a proper Galois extension of $k$, and Jarden's
Proposition 2.5 supplies a pair in $A$ generating a free profinite group of
rank two. Such a group has the nonabelian quotient $S_3$, whereas every
subgroup and every quotient of the abelian group $A$ is abelian, a
contradiction.

Now put $C=C_{G_k}(H)$. Centralizers are closed, since each commutation
condition is closed, and $C$ is normal since $H$ is normal. The subgroup
$H\cap C=Z(H)$ is closed, abelian, and normal in $G_k$, so it is trivial.
If $C$ were nontrivial, the fixed fields
$(k^{\rm sep})^H$ and $(k^{\rm sep})^C$ would both be proper Galois
extensions of $k$. The identity
\[
 G_{(k^{\rm sep})^H(k^{\rm sep})^C}=H\cap C=1
\]
would make their compositum $k^{\rm sep}$, contrary to the cited
Corollary 2.5. Hence $C=1$. The second displayed formula follows by taking
$H=G_M$; when $H=1$, every element centralizes it.
\end{proof}

This is a derivation from the two opened theorems, not a claim of priority.
Our search included the phrases ``absolute Galois centralizer normal
subgroup'', ``Hilbertian prosolvable normal subgroup'', and ``centralizer
Galois abelian extension''. We did not find the exact formula of Statement
B1 stated separately in the opened papers. Statement B2 supplies a
literature-based proof of it after taking a Galois closure of $k$; Statement
B1 independently checks the proposed elementary proof.

\paragraph{Statement B3 (Proved: all the named fields).}\label{res:B3}
\passage{section 2; p.22; scan 45}

For each field in the following table the indicated centralizer is trivial.
Here $k$ is any number field; for the first two rows a finite Galois closure
of $k/\mathbf Q$ may be used when $k(\mu_{\ell^\infty})$ or
$k(\mu_\infty)$ itself is not Galois over $\mathbf Q$.
\begin{center}
\begin{tabular}{ll}
\hline
$M$ & Centralizer in $G_{\mathbf Q}$ of $G_M$\\\hline
$k(\mu_{\ell^\infty})$ & $1$\\
$k(\mu_\infty)$ & $1$\\
$\mathbf Q^{\rm ab}$ & $1$\\
the maximal nilpotent extension of $\mathbf Q$ & $1$\\
the maximal solvable extension of $\mathbf Q$ & $1$\\
the field $\mathbf Q^{\rm tr}$ of totally real algebraic numbers & $1$\\\hline
\end{tabular}
\end{center}

\begin{proof}
The first row follows directly from
$G_{k(\mu_\infty)}\subset G_{k(\mu_{\ell^\infty})}$ and Statement B1.
The second is Statement B1. For all remaining rows we use Statement B2
with $k=\mathbf Q$, which is Hilbertian by Hilbert's irreducibility theorem
(opened as \cite[\S4, Theorem 1]{B_Geyer2012}). We must check that each
Galois extension is proper.

The splitting field of $T^3-2$ has group $S_3$: Eisenstein's criterion at
$2$ gives irreducibility, and its discriminant $-108$ is not a square.
It cannot be contained in an abelian extension or in a pro-nilpotent
extension of $\mathbf Q$. This proves properness for $\mathbf Q^{\rm ab}$
and the maximal nilpotent extension, without using the Kronecker--Weber
identification. The splitting field of $T^5-T-1$ has group $S_5$.
Here is a finite certificate with the inference spelled out. Its reduction
modulo $3$ is irreducible, so its rational Galois group is transitive of
degree five. Its squarefree reduction modulo $2$ has irreducible factors
of degrees two and three, so the group has an element of cycle type
$(2)(3)$ and hence a transposition. A transitive group of prime degree is
primitive. The graph whose edges are the conjugates of that transposition
has connected components forming blocks; it is therefore connected.
Transpositions along a connected graph generate $S_5$.
Thus this Galois group is $S_5$, and is not solvable. It cannot be a quotient
of the prosolvable group of the maximal solvable extension. The field
$\mathbf Q^{\rm tr}$ is Galois, because being totally real is invariant
under conjugation, and it omits $i$. It too is proper.

The program \texttt{code/B\_field\_certificates.py} computes the indicated
factorizations, the discriminants, the derived series orders
$120,60,60$ of $S_5$, and the three Sylow-$2$ subgroups of $S_3$.
Its exact output is attached. The factorization-to-cycle inference follows
by reducing the roots at an unramified prime: the residue Frobenius permutes
the roots in cycles whose lengths are the irreducible factor degrees.
\end{proof}

\subsection{Absolute groups of finitely generated fields and inverse limits}

\paragraph{Statement B4 (Proved: the assertion for $E_K$).}\label{res:B4}
\passage{section 1; pp.6-7; scans 13-15; section 2; p.22; scan 45}

Let $K/\mathbf Q$ be a finitely generated field, fix an algebraic closure
$\overline K$, and write $E_K=G_K$. For every open subgroup $E'\subset E_K$
let $E'^\circ$ be the kernel of its action on all roots of unity in
$\overline K$. Then
\[
 C_{E_K}(E'^\circ)=1.
\]
This proves the assertion accompanying the starred clause of scan 45 in
its stated field-theoretic generality.

\begin{proof}
The field $K$ is Hilbertian: the opened source
\cite[\S4, Theorem 2(e)]{B_Geyer2012} states this for every infinite field
finitely generated over its prime field. Choose a finite Galois extension
$K_1/K$ whose absolute Galois group is contained in $E'$ (the normal core
of $E'$ provides it). The extension $K_1(\mu_\infty)/K$ is Galois and
$G_{K_1(\mu_\infty)}\subset E'^\circ$.
It is a proper subfield of $\overline K$. Indeed $K_1$ is Hilbertian and
$K_1(\mu_\infty)/K_1$ is abelian; the opened Kuyk statement in
\cite[\S4, Theorem 2(g)(2)]{B_Geyer2012} says that every algebraic abelian
extension of a Hilbertian field is Hilbertian. A separably closed field is
not Hilbertian: $Y^2-T$ is irreducible in its rational function field but
every nonzero specialization is reducible in characteristic zero.
Statement B2, applied to the nontrivial normal subgroup
$G_{K_1(\mu_\infty)}\triangleleft G_K$, now gives the result.
\end{proof}

\paragraph{Statement B5 (Proved, with the rank-one correction).}\label{res:B5}
\passage{section 1; pp.6-7; scans 13-15}

The property
\[
 \text{every open subgroup $H\subset G$ has $C_G(H)=1$}
 \tag{B-slim}
\]
is preserved by extensions of profinite groups and by projective limits
whose transition homomorphisms are surjective. A free profinite group of
finite rank at least two has this property. The trivial group has it.
The free profinite group of rank one does not: it is $\widehat{\mathbf Z}$,
and its centralizer of every subgroup is the whole group. Thus the
corollary as worded on scan 15 needs the nonabelian-rank qualification on
every nontrivial free factor.

\begin{proof}
For a free profinite group $F$ of rank $r\ge2$, every finite-index subgroup
$V$ is free profinite of rank $1+[F:V](r-1)$. One can see the comparison
explicitly: the preimage $V_0$ of $V$ in the dense discrete free group
$F_r$ corresponds to a finite covering of its graph with $r$ loops.
Choosing a spanning tree gives a free basis with
$1+[F:V](r-1)$ elements. Its full profinite topology is induced from $F_r$,
because the core in $F_r$ of any finite-index subgroup of $V_0$ still has
finite index. Therefore its profinite completion maps isomorphically onto
$V$, and $V^{\rm ab}\simeq\widehat{\mathbf Z}^{1+[F:V](r-1)}$.

This also proves torsion-freeness without assuming it. A nonidentity
finite-order $z\in F$ has a nonidentity image in a finite quotient $F/N$.
Put $V=N\overline{\langle z\rangle}$. The finite-order image of $z$ in
the torsion-free group $V^{\rm ab}$ is zero, but the map
$V\twoheadrightarrow V/N$, whose target is the nontrivial cyclic group
generated by $zN$, factors through $V^{\rm ab}$. This is a contradiction.

The center of every free profinite group of rank at least two is trivial
by the opened Nakamura \cite[Corollary (1.3.4)]{B_Nakamura1994}; its
hypotheses are the pro-$\mathcal C$ completion of a finite-rank free group
of rank $r\ge2$, with $\mathcal C$ a full class of finite groups. We take
the class of all finite groups. For an open $H\subset F$, choose an open
normal $V\subset H$. Its rank is at least two, so $Z(V)=1$.
The homomorphism $C_F(V)\to F/V$ is injective, because its kernel is
$Z(V)$. Thus $C_F(V)$ is finite; torsion-freeness makes it trivial.
The inclusion $C_F(H)\subset C_F(V)$ proves (B-slim).

For an exact sequence of profinite groups
$1\to A\xrightarrow{i}B\xrightarrow{q}D\to1$, where $A,D$ have
(B-slim), let $H$ be open in $B$ and $z\in C_B(H)$. The quotient map is
open, so $q(H)$ is open in $D$, and $q(z)$ centralizes it. Hence $q(z)=1$.
Then $z\in A$ centralizes $H\cap A$, open in $A$, and so $z=1$.

Finally let $G=\varprojlim_iG_i$, with continuous surjective transition
maps, each $G_i$ satisfying (B-slim). The projections $q_i:G\to G_i$
are surjective. An open $H\subset G$ contains $\ker q_{i_0}$ for some
index $i_0$: finitely many inverse-image neighborhoods contain an open
neighborhood of $1$ in $H$, and directedness combines their indices.
For each $j\ge i_0$ the subgroup $q_j(H)$ is open in $G_j$, since it is
the inverse image of $q_{i_0}(H)$ under $G_j\to G_{i_0}$.
For $z\in C_G(H)$, its coordinate $q_j(z)$ centralizes $q_j(H)$, and is
therefore $1$. These $j$ are cofinal, so $z=1$ in the inverse limit.
\end{proof}

Here is the passage relevant to the manuscript, with its maps specified.
For a regular integral model $Y$ of a function field $L$, and a dense open
$V\subset Y$, restriction of finite \emph{\'etale} covers gives a
continuous surjection $\pi_1(V)\twoheadrightarrow\pi_1(Y)$.
A connected finite \emph{\'etale} cover of the normal integral $Y$ stays
connected on a dense open: its function field is unchanged. This is the
connected-cover criterion for surjectivity of the fundamental-group map.
Every finite separable extension of $L$ is finite \emph{\'etale} on some
dense open of its normalization, after deleting the branch locus. Thus
\[
 G_L\ \xrightarrow{\sim}\ \varprojlim_{V\subset Y}\pi_1(V)
\]
is the continuous isomorphism induced by restriction to the corresponding
maximal extensions unramified on $V$. The same construction after extending
the constant field gives the geometric inverse limit. Whenever the cofinal
model groups are the successive nonabelian free extensions constructed in
\S1 of the manuscript, Statement B5 gives exactly the missing passage
to the limit; it requires surjectivity of the transition maps, and no
injectivity assertion.

There is also a direct verification of the needed property of the
geometric absolute group, independent of that choice of models. Let
$k$ be the algebraic closure of $\mathbf Q$ in $K$, and put
$\pi_K=G_{K\overline k}$. For positive transcendence degree,
$K\overline k$ is a finite extension of a rational function field
$\overline k(t_1,\ldots,t_d)$, so it is Hilbertian by
\cite[\S4, Theorem 2(a),(d)]{B_Geyer2012}. Statement B2 applied to the
normal core of an open subgroup gives (B-slim) for $\pi_K$. In
transcendence degree zero it is trivial. Projection of an element of
$C_{G_K}(E'^\circ)$ to $G_k$ centralizes an open cyclotomic kernel, and
is trivial by B1; the element then lies in $\pi_K$ and centralizes
$E'\cap\pi_K$. This provides the manuscript's original reduction as well.

The three marginal questions on scans 14--15 concern, respectively,
the center of the open arithmetic subgroup, the center of a free
profinite group, and its torsion. In this last proof the first is answered
by B1; the second is answered, with rank at least two, by Nakamura's cited
Corollary (1.3.4); and the third is proved above. All three enter the
manuscript's finite-centralizer argument. The alternative proof B4 uses
the Hilbertian criterion instead and does not require those three
auxiliary steps.

\subsection{Arithmetic and geometric fixed groups}

\paragraph{Statement B6 (Proved: the equivalence used at scan 127).}\label{res:B6}
\passage{section 4; p.63; scan 127}

Let $K/\mathbf Q$ be finitely generated, let
\[
 1\longrightarrow P\longrightarrow E\xrightarrow{p}G_K\longrightarrow1
\]
be an exact sequence of profinite groups, and let
$s:G_{K'}\to E\times_{G_K}G_{K'}$ be a continuous section over a finite
extension $K'/K$. We regard its image as a subgroup of $E$ through the
inclusion of the pullback. For every subgroup $\Omega\subset G_{K'}^\circ$
that is open in $G_{K'}^\circ$,
\[
 C_E(s(\Omega))=C_P(s(\Omega)).
\]
In particular the corresponding two assertions of triviality are
equivalent. This equality applies both to the absolute group of a
function field and to the arithmetic fundamental group of a model.

\begin{proof}
An open subgroup $\Omega$ of $G_{K'}^\circ$ contains
$G_{K''}^\circ$ for a finite extension $K''/K'$.
To construct it, choose an open subgroup $W\subset G_{K'}$ with
$W\cap G_{K'}^\circ\subset\Omega$, and let $K''$ be its fixed field.
Then $G_{K''}^\circ=W\cap G_{K'}^\circ$. For
$z\in C_E(s(\Omega))$, the element $p(z)$ centralizes
$G_{K''}^\circ$ in $G_K$, and is trivial by B4. Thus $z\in P$.
The reverse inclusion of the two centralizers is the inclusion $P\to E$.
Every map used is the displayed continuous quotient or subgroup inclusion;
no comparison of fixed groups under a quotient is inferred in the reverse
direction.
\end{proof}

\paragraph{Statement B12 (Proved: Theorem 1 over cyclotomic fields).}\label{res:B12}
\passage{section 1; pp.6-7; scans 13-15; section 3; pp.30-31; scans 61-63}

Let $k$ be a number field, $F=k(\mu_\infty)$, and let $M/F$ be a finitely
generated field extension. Every open subgroup $H\subset G_M$ has
$C_{G_M}(H)=1$. Since $\mu_\infty\subset M$, its cyclotomic kernel is
$G_M$ itself. For positive transcendence degree, the geometric absolute
group $G_{M\overline F}$ has the same open-subgroup centralizer property;
in transcendence degree zero it is trivial. This proves Theorem 1 of scan
13, including the geometric assertion, over $F$ and its finitely generated
extensions.

\begin{proof}
The abelian extension $F/k$ is Hilbertian by Kuyk's theorem as opened in
\cite[\S4, Theorem 2(g)(2)]{B_Geyer2012}. For $M/F$ algebraic and finitely
generated, $M/F$ is finite and preserves Hilbertianity by that source's
Theorem 2(d). For positive transcendence degree, a transcendence basis
expresses $M$ as a finite extension of a rational function field over
$F$, which is Hilbertian by Theorem 2(a), followed again by (d).
Choose an open normal subgroup $H_0\subset H$. Its fixed field is finite
over $M$, hence Hilbertian and not separably closed, so $H_0\ne1$.
Statement B2 gives $C_{G_M}(H_0)=1$, and hence $C_{G_M}(H)=1$.
For the geometric group, $M\overline F$ is a finite extension of
$\overline F(t_1,\ldots,t_d)$ when $d>0$, so the identical proof applies.
When $d=0$ it is the algebraically closed field $\overline F$, with trivial
absolute Galois group.
\end{proof}

% END sections/B_agent.tex

\section{C. Consequences for cote 149}

% BEGIN sections/C_root.tex
\subsection{The local extension and its packets}

\paragraph{Statement C1 (Proved; correction to the centralizer reading).}\label{res:C1}
Let $U=X\setminus S$ be a hyperbolic curve over a number field $K$, let $x\in S(K)$, and fix a lift $c$ of $x$ to the completed geometric universal cover. With the chosen base point, the local arithmetic group is
\[
 D_c=\Norm_{E_U}(I_c),\qquad
 1\longrightarrow I_c\simeq\Zhat(1)\longrightarrow D_c
   \longrightarrow\Gamma_K\longrightarrow1.
\]
Its conjugation action on $I_c$ is $\chi_K$. In particular
\[
 \Cent_{E_U}(I_c)=D_c\cap p^{-1}(\Gamma_K^\circ).
\]
A local parameter supplies a splitting. The set of $I_c$-conjugacy classes of splittings is a torsor under
\[
 H^1_{\mathrm{cont}}(\Gamma_K,\Zhat(1))
   \simeq\varprojlim_{n\mid m} K^*/K^{*n}.
\]
These identifications describe the torsor; identifying it with its acting group additionally chooses a splitting.
\passage{Section 3, manuscript pp.~26--27, scans 53--55; formula (2) and the first formula (3).}
\begin{proof}
We use the following exact source input. Stix~\cite[Proposition 11 and Corollary 12]{Stix2008}, for a hyperbolic curve over a characteristic-zero field, proves that distinct lifted cusps have inertia groups with trivial intersection and that the geometric and arithmetic normalizers of $I_c$ are $I_c$ and $D_c$, respectively. The identification of local inertia with $\Zhat(1)$ is the characteristic-zero tame description: over an algebraically closed residue field the covers of a punctured strictly henselian smooth disc are the compatible covers $t=u^n$. Conjugating $u\mapsto\zeta_nu$ by a coefficient automorphism sends $\zeta_n$ to $\zeta_n^{\chi_n(\sigma)}$. Passing to the inverse limit gives the displayed action and the formula for the centralizer. A centralizer element lies in the normalizer; within the normalizer its action is trivial exactly when every $\chi_n$ equals $1$.

Choose the compatible roots $t^{1/n}$ and let $\Gamma_K$ act on coefficients while fixing those roots. This constructs a splitting of the local extension; restriction to the finite extensions in the chosen algebraic closure gives the indicated embedding into $E_U$. Any second splitting has the unique form $s_a(\sigma)=a(\sigma)s_0(\sigma)$, with $a$ a continuous $1$-cocycle in $I_c$. Conjugation by an element of $I_c$ changes $a$ by a coboundary. This proves the torsor assertion without choosing an origin for it.

For each $n$, the exact sequence $1\to\mu_n\to\overline K^*\xrightarrow{n}\overline K^*\to1$ gives $H^1(K,\mu_n)=K^*/K^{*n}$. Here the vanishing used for $H^1(K,\overline K^*)$ has the following elementary proof. A continuous cocycle descends to a finite Galois extension $L/K$. Linear independence of the automorphisms of $L$ supplies $b\in L$ for which $\sum_\sigma a_\sigma\sigma(b)\ne0$; applying an automorphism to this sum expresses the cocycle as a coboundary. At the inverse-limit step, the finite groups $\mu_n(K)$ have the Mittag--Leffler property, so the possible $\varprojlim^1 H^0$ term is zero. Equivalently, the finite nonempty sets of compatible representatives of a fixed compatible family of cohomology classes have nonempty inverse limit. Thus continuous $H^1$ with coefficients in $\Zhat(1)$ is the asserted limit of Kummer groups.

For a number field, $\chi_K(\Gamma_K)$ is open in $\Zhat^\times$: $K\cap\Q(\mu_\infty)$ is finite over $\Q$, and restriction identifies this image with the subgroup fixing that intersection. It is nontrivial. Consequently the literal centralizer in scan 54 cannot be the middle group of its extension onto all of $\Gamma_K$. We retain the word visible in the scan and record the required mathematical replacement in the readings file.
\end{proof}

\paragraph{Statement C2 (Proved; the packets are embedded and disjoint).}\label{res:C2}
For hyperbolic $U/K$, the map from $I_c$-conjugacy classes of splittings of $D_c\to\Gamma_K$ to $\pi_U$-conjugacy classes of sections of $E_U\to\Gamma_K$ is injective. Its image is disjoint from the point sections of $U(K)$ and from the packet of any other rational cusp. These assertions include type $(0,3)$ and remain true after each finite extension of $K$.
\passage{Section 3, manuscript pp.~26--27, scans 54--55, including the exception left for type $(0,3)$.}
\begin{proof}
First use the proved centralizer theorem \Rtwo. A local splitting $s$ has $C(s)=I_c$. If two local splittings satisfy $s'=hsh^{-1}$ with $h\in\pi_U$, then $I_c=C(s')=hC(s)h^{-1}=hI_ch^{-1}$. The geometric normalizer formula in C1 puts $h$ in $I_c$. This proves injectivity. For splittings at different cusps, the same equality identifies their lifted inertia groups after conjugation, hence their images in $S$; the two cusps must coincide. A point section has zero cyclotomic centralizer by \Rtwo\ and the separation argument below, whereas a local splitting has the nonzero group $I_c$.

Here is a second proof that also checks the separation used in the preceding sentence. Stix~\cite[Theorem 14(1)--(2)]{Stix2008} assumes injectivity of the point-to-section map for proper curves of genus at least two over every finite extension of the base field. Under exactly that hypothesis it proves all the packet assertions in C2. Its proof first replaces $U$ by a finite étale neighbourhood of the section whose compactification has genus at least two. Two sections coming from different points of the compactification are then separated there. Repeating over all finite neighbourhoods separates different lifted cusps over the same cusp. An element conjugating two local splittings therefore stabilizes the selected lifted cusp, and lies in $I_c$.

For completeness the arithmetic hypothesis follows as follows. Given distinct rational points $a,b$ of a smooth proper curve $Y$ of positive genus, the Abel--Jacobi map based at $a$ sends $b$ to $[b-a]\in J_Y(K)$. This class is nonzero: a function with divisor $b-a$ would have degree one as a map to $\mathbf P^1$, forcing $Y$ to have genus zero. Equality of the two sections would make $[b-a]$ divisible by every positive integer: pull back the multiplication-by-$n$ covering of $J_Y$, and read the equality of the two fibre torsors. The no-divisible-elements statement \Arith\ excludes this. The argument works over every finite extension. The higher-genus neighbourhood can be constructed explicitly over the algebraic closure and then descended by a characteristic refinement and the section. In genus zero, carry three punctures to $0,1,\infty$ and pull back the Fermat quartic; its compactification has genus three. In genus one, first take an etale double cover of the compactification, so that its open part has at least two cusps $P,Q$. Choose a degree-one line bundle $L$ with $L^{\otimes2}\simeq\mathcal O(P+Q)$; it exists because multiplication by two on the elliptic curve's degree-zero Picard group is surjective. The double cover defined by the section with divisor $P+Q$ is branched at these two cusps and its compactification has genus two by Riemann--Hurwitz. In compact genus at least two the identity cover suffices. Each constructed cover is etale on the relevant open curve. Take its Galois closure and an open characteristic subgroup contained in its subgroup. The new compactification dominates the constructed one, so Riemann--Hurwitz preserves genus at least two. The subgroup together with the section descends this cover.

The published Esnault--Hai statement~\cite[Theorem 5.5]{EsnaultHai2008} concerns a smooth projective $X$ over a finitely generated extension of $\Q$, a nonempty open $U$, and a neutral fibre functor for finite connections on $U$. For $g(X)\ge1$, or $g(X)=0$ with $U$ hyperbolic, it bounds the number of rational points in the corresponding normalized universal cover by one. Its Lemma 5.6 supplies a positive-genus finite level in genus zero over any characteristic-zero field. The proof just given is the group-theoretic version of this positive-genus reduction. The attached older preprint has a different item 5.6; the published lemma has been opened and checked.
\end{proof}

\paragraph{Statement C3 (Refuted as written in type $(0,2)$; exact replacement).}\label{res:C3}
\passage{Section 3, manuscript pp. 26, scans 54}

For $U=\Gm/K$, both cuspidal packets are the whole set
$H^1(K,\Zhat(1))$ of section classes. In particular they contain every point section. For an affine curve with nontrivial geometric fundamental group, the disjointness conclusion of scan 54 holds for every hyperbolic type and fails for type $(0,2)$. In types $(0,0)$ and $(0,1)$ the geometric group is trivial. For a proper genus-one curve admitting a section, its cyclotomic centralizer is zero over a number field.
\begin{proof}
The covers $u^n=t$ of $\Gm$ identify $\pi_U$ with $\Zhat(1)$. The point $1$ splits the arithmetic extension, giving $E_U=\Zhat(1)\rtimes\Gamma_K$. The local group at $0$ is this entire group, and the coordinate $t^{-1}$ at infinity gives the same group with inertia generator inverted. Thus each map from local splittings to global splittings is surjective, indeed bijective on conjugacy classes. A point $a\in K^*$ gives its Kummer class in precisely the same group. Since the geometric kernel is abelian, the conjugation action of every splitting on it is $\chi_K$ and $C(s)=\Zhat(1)$ for every splitting.

A sphere and a once-punctured sphere are simply connected over $\mathbf C$; finite-cover comparison identifies their geometric étale groups with $1$. A torus has fundamental group $\Z^2$, whose profinite completion is the Tate module of its Jacobian. In the genus-one case changing a section conjugates by this abelian kernel and does not change its action; \Arith\ gives zero invariants. The finite-cover comparison used here is SGA~1, Exposé XII, Theorem 5.1: analytification induces an equivalence between finite étale covers of a complex scheme locally of finite type and finite étale analytic covers. The comparison with the complex fibre is valid in characteristic zero, as in the same exposé and the specialization statement of Exposé XIII; all our curves descend to finitely generated subfields embedded in $\mathbf C$.
\end{proof}

\paragraph{Statement C4 (Proved; the precisions without section surjectivity).}\label{res:C4}
Over a number field, the section of any rational point of a hyperbolic curve has $C(s)=1$. Every cuspidal section has $C(s)=I_c$ for exactly one lifted cusp $c$. The same formulas hold for germs over finite extensions and for restriction to an arbitrary open subgroup of the appropriate cyclotomic kernel. These formulas alone impose no existence statement for sections whose centralizer is trivial.
\passage{Section 3, manuscript p.~27, scans 55--56; the second formula (3) and the marginal precisions to Conjecture A.}
\begin{proof}
The local formula and its uniqueness are \Rtwo; C2 excludes a point section from every local packet. The open-subgroup statement is the germ version in Part A. More explicitly, an open subgroup of $\Gamma_K^\circ$ is $\Gamma_{F'}$ for a finite extension $F'/K(\mu_\infty)$, and the arithmetic vanishing and packet separation continue over that field. The normality of $\Gamma_K^\circ$ then makes its uniquely detected cusp invariant under the whole section. None of these deductions asserts that a section with zero centralizer is in the image of $U(K)$.
\end{proof}

\subsection{Dévissage and the quantifier over all open subgroups}

\paragraph{Statement C5 (Proved; exact sequences for elementary fibrations).}\label{res:C5}
\passage{Section 2, manuscript pp. 24, scans 49}

Let $Y\to B$ be an elementary fibration in hyperbolic curves in characteristic zero, with $B$ obtained by a finite tower of such fibrations from a point. The elementary hypothesis means a smooth proper curve compactification whose omitted divisor is finite étale over the base. At compatible geometric base points, the maps induced by the fibre inclusion and by the fibration give an exact sequence of profinite groups
\[
 1\longrightarrow\pi_1(Y_{\bar b})\longrightarrow\pi_1(Y_{\overline K})
 \longrightarrow\pi_1(B_{\overline K})\longrightarrow1.
\]
For $b\in B(K)$ there is also a natural isomorphism of extensions
\[
 E_{Y_b}\ \xrightarrow{\ \sim\ }\ E_Y\times_{E_B}s_b(\Gamma_K).
\]
\begin{proof}
We give the part of the proof which prevents a loss of injectivity on completion. Bresciani, Lemma A.22, supplies invariance of geometric fundamental groups under extension of algebraically closed characteristic-zero fields. After descending to a field of definition and embedding it in $\mathbf C$, the smooth proper family and its finite collection of disjoint marked sections are locally trivial as differentiable pairs. To see this locally, lift tangent vector fields from a relatively compact coordinate ball in the base to the total space, choosing the lifts tangent to the marked divisor; a partition of unity patches the lifts, and properness supplies the flows throughout a smaller ball. The flows trivialize the pair, and hence its open complement. Hyperbolic complex curves are aspherical. The homotopy lifting property and the long exact homotopy sequence therefore prove inductively that the base and total space are aspherical and that
\[
 1\to P\to G\to B_0\to1
\]
is exact for their discrete topological fundamental groups.

The goodness input can be stated precisely: Bresciani~\cite[Lemmas 10.9 and 10.11]{Bresciani2021} proves goodness and residual finiteness for elementary anabelian topological stacks, with goodness preserved under extensions whose kernel has finite cohomology with finite coefficients and is finitely generated. For open curve fibres the kernel is a finitely generated free group. Its discrete and completed free-group resolutions have length one; in degree one the comparison identifies derivations by their values on the same free generators, and in degrees at least two both cohomologies vanish. This proves goodness and the finite-cohomology property in this case. Lemma 10.10 of the same source proves that finite-index subgroups of a good group remain good, as needed below. Thus the identical induction includes the open elementary fibrations used here. We now justify the resulting profinite exactness rather than treating it as a formal property of completion.

Take a characteristic finite-index subgroup $N$ of $P$ and put $A=P/N$. The quotient $G/N$ has finite normal subgroup $A$. Its subgroup centralizing $A$ has finite index, meets $A$ in $Z(A)$, and is a central extension of a finite-index subgroup $B_1\subset B_0$ by the finite abelian group $Z(A)$. Goodness of $B_1$ identifies its $H^2$ with continuous $H^2$ of $\widehat B_1$. A continuous cocycle with finite coefficients factors through a finite quotient. The central-extension class thus restricts to zero on a finite-index subgroup $B_2$, and over $B_2$ a splitting supplies a subgroup disjoint from $Z(A)$. Its preimage in $G$ has finite index and intersects $P$ inside $N$. Characteristic finite-index subgroups are cofinal in $P$: intersect the finitely many subgroups of a fixed index bound in a finitely generated group. Hence the full profinite topology on $P$ is induced from $G$, proving $\widehat P\hookrightarrow\widehat G$. The quotient by the closure of $P$ is $\widehat B_0$ because its finite quotients are exactly the finite quotients of $G/P$. This proves the short exact sequence after completion. SGA~1 XII.5.1 identifies these completions with the geometric étale groups.

For the second assertion the morphism of schemes $Y_b\to Y$ induces the displayed map, with identity on the quotient $\Gamma_K$. Its restriction to geometric kernels is the identification of the kernel just proved, so injectivity and surjectivity follow by lifting an element to the quotient and then correcting it in the kernel. All maps are continuous; a continuous bijection of compact profinite groups is a homeomorphism.
\end{proof}

% Optional supporting paragraph for C5; root chooses its insertion point.
An independent exact reference, with no goodness or embedding assumption,
is Hoshi~\cite[Proposition 2.4(i),(ii), pp.~169--172]{B_Hoshi2014}.
His base $S$ is connected, noetherian, separated and normal over a field
of characteristic zero, and his relative hyperbolic curves have smooth
proper compactifications and finite \emph{\'etale} cusp divisors.
Part (i), with $n=1$, $m=0$, $S=B_{\overline K}$ and
$X=Y_{\overline K}$, gives precisely the displayed geometric exact
sequence, including injectivity on the left. Part (ii), with
$S=B$, $X=Y$ and $T=\operatorname{Spec}K\xrightarrow{b}B$, gives
$E_{Y_b}\simeq E_Y\times_{E_B}\Gamma_K$, with the two maps supplied
by the fibre inclusion and the structure map. Compatible geometric
base paths turn his outer isomorphisms into the representatives used
here.

\paragraph{Statement C6 (Proved; point sections of elementary anabelian varieties).}\label{res:C6}
Let $U/K$ be an elementary anabelian variety over a number field and $u\in U(K)$. For every open $\Omega\subset\Gamma_K^\circ$,
\[
 \Cent_{\pi_U}(s_u(\Omega))=1.
\]
Moreover every morphism $f:V\to U$ of nonempty connected finite-type schemes over $\Q$ induces a representative $f_*:E_V\to E_U$ such that
\[
 \Cent_{E_U}(f_*(H))=\Cent_{\pi_U}(f_*(H))=1
 \quad\text{for every open }H\subset E_V^\circ.
\]
Changing geometric base paths conjugates the maps and their centralizers and does not affect these equalities.
\passage{Section 2, manuscript pp.~23--24, scans 48--49, condition b) and its asserted reduction to a point.}
\begin{proof}
Induct along the fibration tower. An element centralizing $s_u(\Omega)$ maps to an element centralizing the point section at the image $b$ in the base. By induction its image is $1$. C5 places it in the geometric group of the fibre over $b$, and identifies the action on that kernel with the point section of the fibre at $u$. C4 for that hyperbolic fibre makes the element $1$. The zero-dimensional base case has trivial geometric group. All comparison maps are the fibre inclusions and fibration projections of C5, so the induction uses injectivity precisely where needed.

Now let $H$ be open in the closed subgroup $E_V^\circ$. There is an open normal subgroup $N\triangleleft E_V$ with $N\cap E_V^\circ\subset H$: an identity neighbourhood in the subspace topology contains the intersection with an open normal neighbourhood in $E_V$. Let $V_N\to V$ be the corresponding connected finite étale cover. It has a closed point $v$ whose residue field $k$ is a number field, since it is nonempty of finite type over $\Q$. Choose base paths so that the induced point homomorphism $s_v:\Gamma_k\to E_V$ has image in $N$. Its cyclotomic kernel has image in $N\cap E_V^\circ\subset H$. Under $f_*$ this is the restriction of the point section at $f(v)$, with arithmetic domain $\Gamma_k$. The first part therefore gives $\Cent_{\pi_U}(f_*(H))=1$.

An element of the arithmetic centralizer projects to an element of $\Gamma_\Q$ centralizing the image of $\Gamma_k^\circ$. Part B proves that centralizer trivial. Thus the arithmetic element lies in $\pi_U$, where it is already zero. This proves the other equality. The argument displays a point after choosing $H$, so it preserves the manuscript's universal quantifier over all open $H$. Regularity of $V$ is unnecessary for this direction.
\end{proof}

\paragraph{Statement C7 (Proved equivalence of predicates; surjectivity remains unresolved over number fields).}\label{res:C7}
\passage{Section 2, manuscript pp. 23–24, scans 48–49; Section 3, manuscript pp. 27, scans 56}

For a hyperbolic curve $U/K$, let $\mathcal S(U/K')$ denote section classes over a finite extension $K'/K$, and $\mathcal S_{\mathrm{nc}}(U/K')$ its noncuspidal part. Condition b) of scan 48 for $V=\operatorname{Spec}K'$ holds for exactly $\mathcal S_{\mathrm{nc}}(U/K')$. Consequently the unresolved map
\[
 U(K')\ \lhook\joinrel\longrightarrow\ \mathcal S_{\mathrm{nc}}(U/K')
\]
has the same surjectivity defect as both the point case of the fundamental conjecture and Conjecture A. This is an equality of the two requested assertions, without assuming either one.
\begin{proof}
C4 and \Rtwo\ identify trivial centralizer on every open subgroup of $\Gamma_{K'}^\circ$ with noncuspidality. C2 puts all point classes in this subset, and the Abel--Jacobi and finite-cover argument proves the displayed injectivity. Every other class is cuspidal exactly when the displayed map is onto the noncuspidal subset. The same statement, read as a claim that every class satisfying b) arises from a point, is the point case of scan 48. These implications apply separately over every finite $K'$. The smallest instance left by this argument is a noncuspidal section of $\mathbf P^1\setminus\{0,1,\infty\}$ over a number field that has not been identified with a point. The calculations in Part D separate a specific pair of point sections; they do not provide a point for an arbitrary class in this unresolved set.
\end{proof}

% Root splice: replaces Statement C8, not the surrounding subsections.
% Root splice: replaces Statement C8, not the surrounding subsections.
\paragraph{Statement C8 (Proved: boundary nonvanishing in every dimension;
local injectivity refuted).}\label{res:C8}
Let $U/K$ be an elementary anabelian variety, namely a finite tower of
relative hyperbolic curves over a field of characteristic zero, and let
$U\hookrightarrow X$ be a smooth proper compactification with
normal-crossings boundary. For a rational boundary point $x$, write
$I(x)$ for the nonempty set of geometric local branches at $x$, and put
$T=\widehat{\mathbf Z}(1)$. The geometric group of the punctured strict
henselian neighbourhood is $T^{I(x)}$, and its natural map
\[
 j_x:T^{I(x)}\longrightarrow\pi_1(U_{\overline K})
\]
has nonzero image. More precisely, its restriction to the diagonal
$d_x:T\to T^{I(x)}$ has nonzero image. Every section through the
punctured henselian neighbourhood therefore has nontrivial cyclotomic
centralizer. When all local branches are rational, that centralizer
contains the whole image of $j_x$; in general it contains the nonzero
image of $j_xd_x$. These assertions hold over every finite base extension
and hence for germs. The assertion that $j_x$ is injective for every
such compactification is false, already on a blown-up product surface.
\passage{Section 4, manuscript pp.~63--65, scans 128--131, especially
the withdrawn assertion $1^\circ$ and the margin at scan 130.}
\begin{proof}
The punctured henselian neighbourhood is
$U_{(x)}=U\times_X\operatorname{Spec}\mathcal O^h_{X,x}$.
Its local exact sequence and action are computed in
Carlson--Corwin~\cite[Proposition 3.4]{B_CarlsonCorwin2026}:
the geometric kernel is
\[
 \operatorname{Hom}\bigl(\mathbf Z[I(x)],T\bigr)=T^{I(x)},
\]
with Galois action obtained by permuting the branches and by the
cyclotomic action on $T$. The natural morphism $U_{(x)}\to U$ supplies
$j_x$ and makes it equivariant for every chosen local section. If the
branches are rational, the cyclotomic kernel fixes this entire local
geometric group. Without that restriction it still fixes the diagonal,
because every permutation fixes a constant tuple. It remains to prove
that the diagonal survives the displayed map.

We use the following precise valuative input. Hoshi
\cite[Lemma 2.8, pp.~175--176]{B_Hoshi2014} treats a relative hyperbolic
curve $Y\to S$, where $S$ is connected, noetherian, separated and
normal over a characteristic-zero field. For a strictly henselian
discrete valuation ring $R$ over $S$, with fraction field $F$, a given
$S$-map $\operatorname{Spec}F\to Y$ extends over
$\operatorname{Spec}R$ exactly when its homomorphism
$\pi_1(\operatorname{Spec}F)\to\pi_1(Y)$ has trivial image.

This gives the same extension criterion for a hyperbolic polycurve over
an algebraically closed characteristic-zero field. Here is the
induction, including the maps involved. For a tower
$U=U_n\to U_{n-1}\to\cdots\to U_0=\operatorname{Spec}\overline K$
and $a:\operatorname{Spec}F\to U$ with trivial fundamental-group
image, its composite to $U_{n-1}$ also has trivial image. By induction
that composite extends to $b:\operatorname{Spec}R\to U_{n-1}$.
Regard $R$ as a ring over the normal base $U_{n-1}$ by $b$.
Hoshi's lemma applied to the last relative curve and to the original
map $a$ extends $a$ to $U_n$. The base case is a point. The converse
follows from $\pi_1(\operatorname{Spec}R)=1$, and uniqueness follows
from separatedness.

Apply this criterion to an explicit arc at the chosen geometric point
$\bar x$ over $x$. Let $r=|I(x)|$ and $d=\dim U$. Choose branch
parameters $z_1,\ldots,z_r$ in the strict henselian local ring and
complete them to a regular system of parameters. Normal crossings and
smoothness identify its completion with
\[
 \widehat{\mathcal O}_{X_{\overline K},\bar x}
   \simeq\overline K[[z_1,\ldots,z_d]],
 \qquad X\setminus U:\ z_1\cdots z_r=0.
\]
The continuous local homomorphism to $R=\overline K[[t]]$ defined by
\[
 z_i\longmapsto t\quad(1\le i\le r),\qquad
 z_i\longmapsto0\quad(r<i\le d)
\]
defines an arc $a:\operatorname{Spec}R\to X_{\overline K}$ through
$\bar x$. Its generic point lands in $U_{\overline K}$ because the
boundary equation pulls back to the nonzero element $t^r$.
Since $R$ is strictly henselian, the arc factors through the strict
henselization at $\bar x$. On punctured neighbourhoods the covers
$z_i^{1/n}$ pull back to $t^{1/n}$ for every $i$ and every $n$.
Consequently the induced geometric homomorphism is precisely
\[
 \pi_1(\operatorname{Spec}\overline K((t)))
   =T\xrightarrow{d_x}T^{I(x)}\xrightarrow{j_x}\pi_1(U_{\overline K}).
\]
If $j_xd_x$ were zero, the extension criterion just proved would give
$\widetilde a:\operatorname{Spec}R\to U_{\overline K}$ with the same
generic-point map as $a$. Its composite with $U_{\overline K}\to
X_{\overline K}$ equals $a$: two maps from this integral trait to the
separated scheme $X_{\overline K}$ which agree at the generic point
agree everywhere. Their closed point would then lie both in
$U_{\overline K}$ and at $\bar x\in X_{\overline K}\setminus
U_{\overline K}$, a contradiction. Thus $\operatorname{im}(j_xd_x)$
is nonzero. The previously proved equivariance now gives the asserted
nontrivial subgroup of the cyclotomic centralizer.

For the failure of injectivity, take
$U=(\mathbf P^1\setminus\{0,1,\infty\})^2$, compactify in
$(\mathbf P^1)^2$, and blow up $(0,2)$. With old local coordinates
$t$ and $u=z_2-2$, the old boundary is $t=0$. In the chart
$t=ab$, $u=b$, the new boundary has the two branches $a=0$ and $b=0$.
Each branch loop maps to the loop around $t=0$, while the loop around
$u=0$ contracts in the original neighbourhood. Thus the actual map is
\[
 T^2\longrightarrow I_0\subset\pi_1(U_{1,\overline K}),
 \qquad(m,n)\longmapsto m+n.
\]
Its kernel is $\{(m,-m):m\in T\}$ and its image is $I_0\ne1$.
The diagonal maps by multiplication by $2$, so it is indeed nonzero.
The finite matrix reductions in \texttt{code/root\_math\_checks.py}
record this kernel and image modulo $2,3,5,25$.
\end{proof}

The converse requires a geometric boundary centre and a factorization
through its actual punctured henselian group for every section with
nonzero cyclotomic centralizer, with a specified compactification and
arithmetic domain. Merely recognizing a nontrivial abelian subgroup
does not construct that factorization. This obligation already contains
the curve section-surjectivity problem through explicit product maps.
Indeed, let $C/K$ be any hyperbolic curve, let $Q=\mathbf
P^1\setminus\{0,1,\infty\}$, and fix a rational cusp section $t$ of
$Q$. For every section $\sigma$ of $C$, the product section
$(\sigma,t)$ of $C\times Q$ has a centralizer containing
$\{1\}\times I_0$. A boundary-local factorization of that section in
the product compactification would, by the first projection, express
$\sigma$ either as the point section at an interior first coordinate
or as a cuspidal section at a boundary first coordinate. Thus a
universal boundary-factorization assertion for these products entails
the point-or-cusp assertion for all such $\sigma$ over the same field.
For a statement about germs, the same argument retains the corresponding
finite-extension domain. No such exhaustion of the section set is
asserted here.

For source comparison, Carlson--Corwin
\cite[Theorem 3.7 and Proposition 4.12]{B_CarlsonCorwin2026} prove
compactification independence of the \emph{closure} of the boundary
section set and an inclusion into the \emph{closure} of the inductively
defined cuspidal set of a tower. The closure bars are essential in both
statements. Their Proposition 3.4 supplies the precise local group and
action used above; the nonvanishing argument here additionally uses
Hoshi's valuative lemma. None of these comparisons supplies the
universal factorization in the preceding paragraph.

\subsection{Admissible germs and loop recognition}

\paragraph{Statement C9 (Proved; the precise object in scan 66).}\label{res:C9}
The admissible germs for $U_{03}=\mathbf P^1\setminus\{0,1,\infty\}$, defined by trivial cyclotomic centralizer, are exactly the noncuspidal germs. There is a natural injective $\Gamma_\Q$-equivariant map from $U_{03}(\overline\Q)$ to this germ set. Equality is precisely the section-surjectivity assertion for $U_{03}$ over all number fields. Classes of sections defined over all of $\Gamma_\Q$, by themselves, do not give the germ set appearing in this equivariant formulation.
\passage{Section 3, manuscript p.~32, scan 66.}
\begin{proof}
Take the filtered colimit, under restriction, over finite number fields in $\overline\Q$ of the sets of section classes. A point lies over a finite field, so it determines an element of this colimit. If two points have equal germs, their classes agree over some common finite field; the point injectivity in C2 proves equality of the points. Conjugating fields, points and sections by $\Gamma_\Q$ makes this construction equivariant. The identification of admissible and noncuspidal germs is \Rtwo\ applied after restriction. The exhaustion statement follows exactly as in C7, with the quantifier over all fields retained.

There is a concrete type check on the literal whole-group reading. For a section $s:\Gamma_\Q\to E_{03}$, transport by $\tau\in\Gamma_\Q$ using the lift $s(\tau)$ changes the section by
$s(\tau)s(\tau^{-1}\sigma\tau)s(\tau)^{-1}=s(\sigma)$.
Thus this transport on whole-group section classes is trivial. The action on $U_{03}(\overline\Q)$ is nontrivial, for example on the two roots of $z^2-z+1$. Germs carry the required action because their arithmetic domains change. The phrase ``à \ldots\ cas près'' is not used as a theorem: neither the suggested reading ``deux'' nor an unspecified exceptional-case reading repairs a mismatch of arithmetic domains. The exceptions $(0,1)$ and $(0,2)$ calculated in C3 concern different types; the scan here fixes type $(0,3)$, so those calculations do not identify its proposed exceptions. The available mathematics does not identify the omitted words.
\end{proof}

\paragraph{Statement C10 (Proved correction; the character-free strengthening remains unresolved).}\label{res:C10}
For a hyperbolic curve over a number field, its lifted inertia subgroups are exactly the nontrivial subgroups $C(s)$ as $s$ ranges over germs. They are also exactly the maximal procyclic subgroups of $\pi_U$ whose arithmetic normalizer has open image and acts on them through the cyclotomic character. The condition involving only open image of the normalizer, with arbitrary closed subgroups permitted, is false. The same condition restricted to maximal procyclic subgroups without specifying the character is not decided here in the full group.
\passage{La Longue Marche, \S34, cote 140-2, scan 483; cote 149, scans 55--57.}
\begin{proof}
Every geometric cusp and a nonzero tangent vector at its image are defined over a finite extension. C1 supplies the corresponding tangential section, and \Rtwo\ gives $C(s)=I_c$. The reverse inclusion is exactly \Rtwo. The second characterization is Stix~\cite[Theorem 20 and Corollary 21]{Stix2008}, with the arithmetic vanishing of Part A verifying its hypotheses; the maximality step is also explicit: a larger procyclic group would be abelian and intersect $I_c$ nontrivially, and Stix's Corollary 13 then puts that larger group inside $I_c$.

For the unqualified normalizer condition, $\pi_U$ itself and every characteristic subgroup of it have normalizer $E_U$. To make this an explicit counterexample, take $U_{03}$. Its geometric group maps onto the nonabelian finite group $S_3$ by sending its two free generators to $(12)$ and $(123)$. Therefore it is not procyclic, whereas it satisfies the full-image normalizer condition and contains every proper candidate. The closure of its commutator subgroup is characteristic as well. This supplies the comparison map witnessing why the original maximal-subgroup test cannot select the cusp groups. Part E verifies the CM class-two counterexample for the procyclic formulation. The maps from the full group onto these nilpotent quotients send normalizers into quotient normalizers; lifting a quotient normalizer element or a normalized cyclic subgroup back through every level is an additional requirement. The precise remaining full-group question is the existence of a maximal procyclic subgroup with open arithmetic normalizer and no nonzero potentially cyclotomic primary component, already for a once-punctured CM elliptic curve. The finite and Lie calculations below do not supply such a compatible lift.
\end{proof}

% END sections/C_root.tex


% BEGIN sections/B_birational.tex
\subsection{The function-field centralizer and the margin of scan 127}

\paragraph{Statement B7 (Proved: the birational centralizer criterion).}\label{res:B7}
\passage{section 2; p.22; scan 45; section 4; pp.62-63; scans 126-127}

Let $K$ be a field finitely generated over $\mathbf Q$, and let $X/K$ be a smooth, proper, geometrically connected
curve of any genus, and $L=K(X)$. Fix an algebraic closure $\overline L$
and the resulting copy of $\overline K\subset\overline L$. Put
\[
 E_L=\operatorname{Gal}(\overline L/L),\qquad
 P_L=\operatorname{Gal}(\overline L/\overline K(X)).
\]
Let $s:G_K\to E_L$ be a continuous section of the restriction
$E_L\twoheadrightarrow G_K$. Then
\[
 C_{P_L}(s(G_K^\circ))\ne1
\]
holds exactly for the sections whose image lies in the decomposition
group $D_v$ of a prolongation $v$ to $\overline L$ of a $K$-rational place
$x$ of $L$. In that case the prolongation is unique and
\[
 C_{P_L}(s(G_K^\circ))=I_v\simeq\widehat{\mathbf Z}(1),\qquad
 C_{E_L}(s(G_K^\circ))=I_v.
\]
The same assertion holds for a germ over $G_{K'}$, for every finite
extension $K'/K$, with a $K'$-rational place of $K'(X)$ and with unchanged
geometric group $P_L$. This proves the margin formula of cote 149,
\S4, manuscript p.~63, scan 127, in relative transcendence degree one.

\begin{proof}
For every nonempty $K$-open $U\subset X$, restriction gives a continuous
surjection
\[
 q_U:E_L\twoheadrightarrow E_U=\pi_1(U),
 \qquad q_U|_{P_L}:P_L\twoheadrightarrow\pi_1(U_{\overline K}).
\]
For $V\subset U$ the comparison map $q_{VU}:E_V\twoheadrightarrow E_U$
satisfies $q_U=q_{VU}q_V$, and all maps commute with their projections to
$G_K$. The function-field argument in B5 gives
$E_L=\varprojlim_U E_U$ and
$P_L=\varprojlim_U\pi_1(U_{\overline K})$.
Deleting at least three geometric points gives a cofinal family of
hyperbolic affine opens, in every genus.

Write $C=C_{P_L}(s(G_K^\circ))$ and choose $1\ne z\in C$.
Some hyperbolic open $U_0$ satisfies $q_{U_0}(z)\ne1$.
For every hyperbolic $V\subset U_0$, its image $z_V$ remains nontrivial,
because it maps to $z_{U_0}$. Apply the centralizer theorem of Part A and its proved finitely generated extension in Statement AE1 to
$s_V=q_Vs$. It gives a unique cusp $c_V$ of the completed universal
pro-cover of $V$, fixed by $s_V(G_K)$, with
\[
 C_{\pi_1(V_{\overline K})}(s_V(G_K^\circ))=I_{c_V}.
\]

We verify compatibility instead of assuming it. The map of normalized
completions associated to $q_{VU_0}$ sends $c_V$ either to a cusp over
$X\setminus U_0$ or to a point over $U_0$. In the latter case its inertia
maps trivially, since the universal cover of $U_0$ is \emph{\'etale}
there. That would kill $z_{U_0}$, which is impossible. Thus the image is
a cusp, and its inertia contains the nonidentity element $z_{U_0}$.
The distinct-inertia assertion of Stix
\cite[Proposition 11]{B_Stix2008}, for hyperbolic curves in characteristic
zero with full profinite geometric group, identifies that image cusp
with $c_{U_0}$. The same argument for any pair $V\subset U$ in this
cofinal family gives compatibility of all $c_U$.

These compatible points determine a valuation prolongation $v$ of
$\overline L$. Here is the field comparison. Every finite extension of
$L$ inside $\overline L$ is unramified on some $U$; the normalization of
$X$ in that extension is a finite level in the corresponding completed
universal tower. The points supplied by $c_U$ restrict compatibly on all
such finite normalizations. Their valuation rings have a directed union
inside $\overline L$, and for each $a\in\overline L^*$ either $a$ or
$a^{-1}$ belongs to that union, as this holds already at a finite level
containing $a$. This is the required valuation ring. Since all $c_U$ are
fixed by the image of the full section, $s(G_K)\subset D_v$. The
underlying geometric point of $X$ is fixed by $G_K$, so its residue field
is $K$: it is a $K$-rational place $x$.

The characteristic-zero local exact sequence is
\[
 1\longrightarrow I_v\simeq\widehat{\mathbf Z}(1)
 \longrightarrow D_v\longrightarrow G_K\longrightarrow1.
\]
For a hyperbolic $U$ deleting $x$, its map to $E_U$ is the map of the
punctured henselian neighborhood at $x$. Stix
\cite[Lemma 10]{B_Stix2008} proves that this map is injective for
$\chi(U)\le0$; its image is the cusp decomposition group, and its
restriction identifies $I_v$ with $I_{c_U}$. The arithmetic action on this
inertia is the cyclotomic action. Thus $I_v\subset C$.

For the reverse inclusion, let $w\in C$. For every $U$ in the cofinal
family, $q_U(w)\in I_{c_U}$. Choose $a\in I_v$ with
$q_{U_0}(a)=q_{U_0}(w)$. The comparison maps
$I_{c_V}\to I_{c_{U_0}}$ for $V\subset U_0$ are isomorphisms, because
both are identified with the same $I_v$ by the preceding local
injections. Therefore $q_V(a)=q_V(w)$ for every such $V$. The inverse-limit
isomorphism gives $w=a$, so $C=I_v$.

Conversely a section in a decomposition group $D_v$ centralizes $I_v$
on $G_K^\circ$, so $C\ne1$ and the equality just proved applies.
Uniqueness follows from the uniqueness of every $c_U$: any other
prolongation fixed by the section would supply the same compatible
system of points. The equality of the centralizers in $E_L$ and $P_L$
is B6. Passing to a finite $K'/K$ replaces $E_L$ by the inverse image of
$G_{K'}$, and does not change $P_L$; repeating the same argument gives
the assertion for germs.
\end{proof}

\paragraph{Statement B8 (Refuted in the written function-field form;
proved for the model quotient).}\label{res:B8}
\passage{section 4; p.63; scan 127}

For a section in $D_v$ the two groups written as trivial in (32)--(33)
of scan 127 are both $I_v$, and are nontrivial. For a hyperbolic model
$U/K$ containing the rational point $x$ underneath $v$, its point section
$s_x=q_Us$ satisfies
\[
 C_{E_U}(s_x(G_K^\circ))
 =C_{\pi_1(U_{\overline K})}(s_x(G_K^\circ))=1.
\]
The comparison responsible for the two formulas is the continuous
surjection $q_U:E_L\twoheadrightarrow E_U$; it kills $I_v$ when $x\in U$.

\begin{proof}
The function-field formula is B7. For the model, the point section and a
cuspidal section are disjoint by Stix
\cite[Theorem 14(2)]{B_Stix2008}, whose hypothesis is injectivity of point
sections on proper curves of genus at least two over all finite base
extensions. Over the present finitely generated field this hypothesis follows from the
no-divisible-points argument of Part A, extended in Statement AE1, and the Jacobian Kummer
comparison. The centralizer theorem of Part A therefore gives the
geometric equality $1$, and B6 gives the arithmetic equality.
The local map at an interior point is the map from a strict henselian
trait into an \emph{\'etale} neighborhood, so its ramification inertia is
in the kernel of $q_U$.

For a completely explicit counterexample take $K=\mathbf Q$,
$L=\mathbf Q(t)$, $x=0$ on $X=\mathbf P^1$, and
$U=\mathbf P^1\setminus\{1,2,\infty\}$. A parameter $t$ at $0$ splits
$D_v\to G_{\mathbf Q}$, by the local construction in Stix
\cite[Proposition 8]{B_Stix2008}; compose that splitting with
$D_v\hookrightarrow E_L$. The inertia $I_v$ is already seen to be
nontrivial on the quadratic extension $\overline{\mathbf Q}(t^{1/2})$
of $\overline{\mathbf Q}(t)$, where it maps onto $\mathbf Z/2\mathbf Z$.
The program \texttt{B\_field\_certificates.py} checks the discriminant
$\operatorname{disc}_T(T^2-t)=4t$, with a simple zero at $t=0$.
All of $I_v$ is killed on passage to $E_U$, and the section there is the
point section of $0\in U(\mathbf Q)$.
\end{proof}

We inspected scan 127, including its lower half. The two written
subscripts read $L$. The proved comparison supports the replacement of
the \emph{mathematical assertion} by its $U$ version in the indicated
context. It does not authorize silently changing the transcription, and
does not establish what the author intended. The slash/comma variation
in $\pi_{L/K}$ has no effect on the argument once that group is defined
as the displayed geometric kernel. Notice also that the arithmetic and
geometric formulas with $L$ are still equivalent to each other by B6;
their common value, rather than that equivalence, is the erroneous
assertion of triviality for a decomposition section.

\paragraph{Statement B9 (Proved: the precise logical content of statement 2).}\label{res:B9}
\passage{section 2; p.22; scan 45; section 4; pp.62-63; scans 126-127}

For $L=K(X)$ as in B7, the following two predicates on the set of germs
of sections have the same extension:
\[
 \{[s]:C_{P_L}(s(G_{K'}^\circ))\ne1\}
 =
 \{[s]:s(G_{K'})\subset D_v\text{ for a rational place }v\}.
\]
Here the equivalence relation on sections is conjugation by $P_L$, and
$K'/K$ ranges through finite extensions. Consequently the universal
assertion in scan 45, restricted to relative transcendence degree one,
and the birational section assertion a) of scans 126--127 over all finite
base extensions are logically equivalent.

\begin{proof}
Statement B7 proves equality of the predicates for each actual section
over each $K'$. Both predicates are invariant under $P_L$-conjugacy:
conjugation sends the centralizer to its conjugate and sends a valuation
decomposition group to that of the conjugate prolongation. Restriction
of a section to an open subgroup gives the finite-extension case of the
same assertion. Universal quantification over those sections therefore
gives the stated equivalence. We make no assertion that the universal
predicate is satisfied by all sections.
\end{proof}

The exact remaining assertion is the existence, for every birational
section $s$, of a nonidentity element of
$C_{P_L}(s(G_K^\circ))$. Our proved criterion detects every section for
which such an element is present; it does not supply the element for an
arbitrary section. A basic unproved instance of that remaining assertion
is $K=\mathbf Q$, $L=\mathbf Q(t)$, with an arbitrary continuous section
$G_{\mathbf Q}\to G_{\mathbf Q(t)}$. The finite quotient calculations
in this work concern specified decomposition sections and do not settle
that universal instance.

\subsection{Packets and recovery of the curve's points and open models}

\paragraph{Statement B10 (Proved: packets are the nonzero centralizer fibres).}\label{res:B10}
\passage{section 3; p.28; scan 57; section 4; p.63; scan 127}

For the same $K,X,L$, let $\mathcal S_L(K)$ be the set of continuous
sections $G_K\to E_L$, modulo $P_L$-conjugacy. On its subset
$\mathcal S_L^{\rm dec}(K)$ of decomposition sections the map
\[
 [s]\longmapsto[C_{P_L}(s(G_K^\circ))]_{P_L}
\]
has one fibre for each $x\in X(K)$. For a fixed prolongation $v$ above
$x$, that fibre is precisely the injective image of the set of splittings
$G_K\to D_v$, modulo $I_v$-conjugacy. It is a torsor under
\[
 H^1(G_K,I_v)=H^1(K,\widehat{\mathbf Z}(1))
 \simeq\varprojlim_n K^*/(K^*)^n.
\]
Every $x\in X(K)$ has a nonempty such fibre. This proves the packet
description of cote 149, \S3, manuscript p.~28, scan 57, for function
fields of curves over fields finitely generated over $\mathbf Q$.

\begin{proof}
The map is well defined by conjugacy functoriality, and B7 identifies
its nontrivial value with $I_v$. We first prove the required normalizer
identity in the function field:
\[
 N_{E_L}(I_v)=D_v,\qquad N_{P_L}(I_v)=I_v.
\]
For $g$ in the first normalizer, the inertia groups of $v$ and $gv$
coincide. On every sufficiently small hyperbolic open deleting their
centres they have equal nontrivial images. Stix's Proposition 11 then
identifies the corresponding cusps. Varying the open and the finite
normalization identifies the two valuations, by the construction in B7.
Thus $g\in D_v$. Conversely $D_v$ normalizes its inertia by definition.
Intersecting with $P_L$ gives the second identity, since the residue
field of the prolongation in $\overline L$ is $\overline K$.

Now suppose two splittings $s,t$ into $D_v$ become conjugate by
$g\in P_L$ in $E_L$. By B7, $gI_vg^{-1}=I_v$. The normalizer identity
forces $g\in I_v$, giving injectivity. Distinct underlying points have
nonconjugate inertia groups: a $P_L$-conjugation acts over
$X_{\overline K}$, and hence cannot change their underlying geometric
points. Conversely all prolongations above the same $K$-rational point
are conjugate by $P_L$. At each finite geometric Galois normalization the
Galois group acts transitively on the points above that point; compactness
of $P_L$ gives a compatible conjugating element over the full tower.
This proves the description of the fibres.

For completeness, fix one splitting $s_0$ into $D_v$. Every other one is
uniquely $s_a(\gamma)=a(\gamma)s_0(\gamma)$, where $a$ is a continuous
$I_v$-valued $1$-cocycle for the cyclotomic action. The equation expressing
that $s_a$ is a homomorphism is exactly the cocycle equation, and
conjugation by $I_v$ is exactly addition of a continuous coboundary.
Thus the set is a torsor under continuous $H^1(G_K,I_v)$. The Kummer
comparison in Part C(a) gives the displayed inverse limit. A local
parameter at $x$ gives the initial splitting: the compatible covers
obtained by adjoining its $n$th roots give the splitting of the punctured
henselian neighborhood, as in the opened Stix Proposition 8. Its
inclusion into $E_L$ provides the required decomposition section.
\end{proof}

\paragraph{Statement B11 (Proved: what the packet construction recovers).}\label{res:B11}
\passage{section 3; p.28; scan 57}

From $E_L\to G_K$ with its given cyclotomic character, the nontrivial
centralizers of germs of sections recover $X(\overline K)$ as a $G_K$-set
and hence recover $X(K')$ for every finite extension $K'/K$.
They also recover, for every finite $G_K$-stable subset
$I\subset X(\overline K)$, the quotient extension
$E_{X\setminus I}\to G_K$. These data determine the $K$-isomorphism
class of every open model $X\setminus I$, with compatibility under
inclusions. No assertion that every birational section is a decomposition
section is needed for this reconstruction: the nontrivial-centralizer
subset is already exactly characterized by B7.

\begin{proof}
For a finite $K'/K$ inside $\overline K$, form the set of packets
\[
 \mathscr P_L(K')=
 \bigl\{[C_{P_L}(s(G_{K'}^\circ))]_{P_L}:
            s\text{ a section over }G_{K'},\ C_{P_L}(s(G_{K'}^\circ))\ne1\bigr\}.
\]
This is the set of centralizer conjugacy classes, rather than the set
of individual splittings in its torsor fibres. B10 identifies it with
$X(K')$. Restriction to $G_{K''}$ for $K''\supset K'$ leaves the
centralizer $I_v$ unchanged by B7, so it defines the transition map on
these packet sets and corresponds exactly to $X(K')\subset X(K'')$.
Every geometric point is defined over a finite extension, and every
such point has a decomposition section by B10. Hence their filtered
colimit is exactly $X(\overline K)$.

For $\tau\in G_K$, choose $g\in E_L$ mapping to $\tau$. The map
$[I_v]_{P_L}\mapsto[gI_vg^{-1}]_{P_L}$ is independent of that lift,
since two lifts differ by $P_L$. It sends the packet over $K'$ to the
packet over $\tau(K')$ at the conjugate geometric point. This defines
the required $G_K$-action on the colimit. Its fixed points under
$G_{K'}$ are exactly $X(K')$.

Write $U=X\setminus I$. In $E_L$, form the closed normal subgroup
\[
 R_U=\overline{\langle\!\langle I_v:
       v\text{ lies over }x\in U(\overline K)\rangle\!\rangle}.
\]
It lies in $P_L$. We claim that the quotient map identifies
$E_L/R_U$ with $E_U$. To check this it suffices to check its finite
quotients. A finite Galois extension of $L$ kills these inertia groups
exactly when its normalization over the regular curve $U$ is unramified
at every closed point. In characteristic zero the residue extensions are
separable; a finite normalization unramified at these discrete valuation
rings is finite \emph{\'etale}. Conversely a finite \emph{\'etale} cover
has trivial inertia there. Thus the two profinite quotients have exactly
the same finite quotients, with the same maps from $E_L$. Their canonical
continuous comparison is an isomorphism. This also proves compatibility:
for $V\subset U$ one has $R_V\subset R_U$, producing precisely the
continuous surjection $E_V\twoheadrightarrow E_U$.

To recover the schemes, we use the opened Tamagawa
\cite[Theorem (0.3), with Conjecture (0.1)(ii) as its stated formulation]
{B_Tamagawa1997}. Its exact setting is a field $k$ finitely generated over
$\mathbf Q$ and two smooth, geometrically connected, affine hyperbolic
$k$-curves. The natural map from their $k$-isomorphisms to continuous
isomorphisms of their arithmetic fundamental groups over $G_k$, modulo
inner automorphisms from the target geometric group, is bijective.

Choose a finite $G_K$-stable set $J\supseteq I$ containing at least three
geometric points. Then $V=X\setminus J$ is affine hyperbolic, and the
reconstructed quotient $E_V\to G_K$ determines $V$ by that theorem.
Its unique smooth proper compactification determines $X$. The remaining
labels $J\setminus I$ are its cusp inertia classes, already supplied by
the quotient maps from $E_L$, and adding exactly these points to $V$
gives $U$. For two choices $J_1,J_2$, let $J=J_1\cup J_2$ and first
reconstruct the common hyperbolic refinement $V=X\setminus J$.
Fill the cusp points labelled by $J\setminus J_i$ in its smooth proper
compactification. The arithmetic fundamental group of the resulting
open is the quotient of $E_V$ by the normal closure of those cusp
inertias: the same finite-normalization argument used for $R_U$ proves
this assertion. That quotient is the already reconstructed
$E_{X\setminus J_i}$. The result of filling is itself affine
hyperbolic, so Tamagawa's bijection gives the unique isomorphism to
the independently reconstructed $X\setminus J_i$ inducing this
specified group isomorphism. A curve isomorphism extends uniquely to
the smooth proper compactifications, and its maps on cusp inertia
classes give the specified labels there. Adding the further labelled
points needed for $U$ now gives compatible isomorphisms of the desired
opens. On three choices the isomorphisms agree on their common
hyperbolic refinement by the same uniqueness. Thus both the recovered
isomorphism class and the inclusion maps are independent of the
auxiliary set $J$.
\end{proof}

% END sections/B_birational.tex

\section{D. Fixed points and cyclotomic base fields}

% BEGIN sections/D_root.tex
\subsection{Continuous actions and points of the completed universal cover}
\paragraph{Statement D0a (Proved; Part D(a)).}\label{res:D0a}
\passage{Section 3, manuscript pp. 29–30, scans 59–61}

Let $U/K$ be a hyperbolic curve in characteristic zero, with a geometric base point. The inverse system of all connected pointed finite etale geometric covers defines $\widetilde U$, and their smooth proper normalizations define $X'$. The semilinear automorphisms of $\widetilde U$ that cover the given action on $U_{\overline K}$ form $E_U$. Give this group the topology of its actions on all finite-cover fibres, together with its action on each finite field of coefficients. With this topology, continuous sections $\Gamma_K\to E_U$ are exactly continuous semilinear actions of $\Gamma_K$ on this pointed-cover system compatible with the projection to $U$ and inducing the prescribed action on $\overline K$.

A fixed geometric pro-point at finite distance is a lift of a point of $U(K)$ and its stabilizer is the point section. A fixed pro-point at infinity is a lifted cusp whose inertia subgroup is normalized by the section. These statements give the dictionary requested in cote 149, section 3, manuscript pp.~29--30, scans 59--61.

\begin{proof}
For clarity, a geometric pro-point means a compatible family of geometric points of the finite levels. The fibre over a geometric point of $U$ is a simply transitive profinite $\pi$-set; the fibre over a boundary point in $X'$ is a profinite $\pi$-set with stabilizers $I_c$. Finite etale covers of $U$ with a fixed geometric fibre form its Galois category. An automorphism of the fibre functor is determined by its compatible permutations of the finite fibres. Requiring semilinearity and the given action on $U$ identifies the coefficient automorphism with its image in $\Gamma_K$. This constructs the extension $1\to\pi\to E_U\to\Gamma_K\to1$, including its profinite topology, and gives both directions of the asserted correspondence. It is not necessary to choose a presentation of $\pi$ for this construction.

For the equivalent rational-function description, each function and each cover is defined at a finite level over a finite coefficient field. An open stabilizer of the finitely many coefficients, functions and their conjugates determining that level is precisely a basic neighbourhood in the topology just described. Conversely, fixing a sufficiently refined finite pointed cover fixes every specified finite collection of its rational functions. Thus on this group of automorphisms over $U$ the inverse-limit topology is the topology of pointwise convergence on discrete rational functions and coefficients. Automorphisms without the condition of preserving the map to $U$ are addressed in Part D(e).

A fixed pro-point projects to a $\Gamma_K$-fixed geometric point of $U$, hence a $K$-rational point. Its arithmetic stabilizer maps isomorphically to $\Gamma_K$, because the geometric deck action is free there. Inclusion of the image of a section into that stabilizer is therefore equality. Conversely the decomposition section of a $K$-point fixes the corresponding lift. At a boundary pro-point, the arithmetic stabilizer is $D_c=N_{E_U}(I_c)$ by A2, so fixing it is exactly normalization of its inertia. All assertions are compatible with changing the geometric base point by an inner automorphism.
\end{proof}

\paragraph{Statement D0b (Proved; Part D(b), stronger uniqueness).}\label{res:D0b}
\passage{Section 3, manuscript pp. 29–30, scans 59–61}

Let $F$ be a finite extension of a number-field cyclotomic closure, or a finitely generated extension of such a field. For every hyperbolic $U/F$ and every continuous section $s_F:\Gamma_F\to E_U$, the action on $X'$ has at most one fixed geometric pro-point in all. Consequently, for a number field $K$ and a section $s:\Gamma_K\to E_U$,
\[
 (X')^{s(\Gamma_K^\circ)}=(X')^{s(\Gamma_K)},
 \qquad \#(X')^{s(\Gamma_K^\circ)}\le1.
 \tag{D.1}
\]
The same assertion holds for finitely generated $K/\mathbb Q$. This proves all four requested uniqueness statements and the proper-curve assertion. It identifies the existence part of Conjecture B with the point-or-cusp surjectivity assertion, through D0a; no such surjectivity is assumed. The manuscript passages are scans 59--61, manuscript pp.~29--30.

\begin{proof}
Two distinct pro-points differ at a finite geometric covering level. Replace its subgroup by an open characteristic subgroup $H$ contained in it, and then refine $H$ so that the compactification of its covering has positive genus. Here are details of the latter step. A hyperbolic genus-zero curve has at least three punctures. Over the algebraic closure, carrying three of them to $0,1,\infty$ and pulling back the Fermat cubic gives an etale cover whose compactification has genus one. For positive genus no initial change is needed. A Galois closure followed by an open characteristic refinement gives a cover dominating this one; Riemann--Hurwitz preserves positive genus. Characteristic refinements exist since a finitely generated profinite group has finitely many subgroups of each bounded index. This is also the content of the opened Esnault--Hai Lemma 5.6, stated for a characteristic-zero genus-zero hyperbolic curve \cite{EsnaultHai2008}.

The group $H s_F(\Gamma_F)$ defines a finite etale $V/F$ with the two chosen pro-points projecting to distinct $F$-points $v_1,v_2$ of its compactification $Y$. They may be interior or boundary points. Both induce the same section of $\pi_1(Y)$: their stabilizers contain $s_F(\Gamma_F)$, whose image is a section after the compactification homomorphism. With $v_2$ as origin for $Y\to J_Y$, the difference class $[v_1-v_2]$ therefore has zero Kummer class. The finite multiplication covers on $J_Y$ identify this with divisibility by every integer in $J_Y(F)$. A1 and AE1 make it zero. A divisor of degree zero of the form $v_1-v_2$ can be principal for distinct points only on a genus-zero curve, since its function would have degree one. This contradicts the positive genus of $Y$ and proves uniqueness. The proof applies over every finite extension needed to define a level, and AE1 supplies the same arithmetic input over finitely generated extensions of $F$.

For (D.1), use the base-change identification $\pi_1(U_{K(\mu_\infty)})=E_U\times_{\Gamma_K}\Gamma_K^\circ$. The preceding argument applies to the restricted section. Normality of $\Gamma_K^\circ$ means that $s(\Gamma_K)$ preserves its fixed-point set. A singleton preserved by a group is fixed pointwise. Both sides of (D.1) are empty or the same singleton. This proves the equality also in the absence of a point, without making an existence assertion.
\end{proof}

In particular, two different interior points over a cyclotomic field cannot have the same full section class: choose compatible lifts for the common section and apply D0b. The exact abelianized obstruction and a finite cover separating the smallest example are supplied next. Over number fields, the smallest existence problem left by this argument is an arbitrary section for $\mathbb P^1_{\mathbb Q}\setminus\{0,1,\infty\}$; D0b settles uniqueness for it, while establishing neither an interior point nor a boundary point for that arbitrary section.

% END sections/D_root.tex


% BEGIN sections/D_agent.tex
% Research contribution D: files and identifiers owned by cyclotomic_coverings.
\subsection{Cyclotomic points: abelianization and explicit finite covers}

\paragraph{Statement D1 (Proved).}\label{res:D1}
\passage{section 3; p.30; scan 62}

Let $F$ be a finite extension of $\mathbb Q(\mu_\infty)$, let
$U=X\setminus S$ be a smooth geometrically connected hyperbolic curve over
$F$, and let $x,y\in U(F)$.  We use geometric conjugacy classes of continuous
sections and the natural continuous quotient
$E_U\twoheadrightarrow E_U/\overline{[\pi,\pi]}$.
For $g(X)\geq1$, equality of the images of the two point sections under this
quotient implies $x=y$.  For $g(X)=0$ and $x\ne y$, equality occurs precisely
when every geometric point of $S$ is $F$-rational and
\[
 \frac{h(a)}{h(b)}\in\mu_\infty\qquad(a,b\in S),\qquad
 h(t)=\frac{t-x}{t-y}.
\]
Here we choose any $F$-coordinate on $X\simeq\mathbb P^1$; the displayed
ratios have the evident homogeneous interpretation at infinity and do not
depend on the coordinate.  In particular, on
$U=\mathbb P^1\setminus\{0,1,\infty\}$ equality is equivalent to
$x/y,(1-x)/(1-y)\in\mu_\infty$.  Its distinct ordered pairs are exactly
\[
 (x,y)=\left(\frac{a(1-b)}{a-b},\frac{1-b}{a-b}\right),\quad
 a,b\in\mu_\infty,\quad a,b\ne1,\quad a\ne b.
\]
This answers the abelianized question of scan 62 (manuscript p.~30).

\begin{proof}
We make the comparison with Kummer theory explicit.  Choosing $y$ as origin
of the generalized Abel--Jacobi map gives a morphism $U\to J_S$.  Its map on
geometric abelian fundamental groups is the identification with the full
Tate module $T(J_S)$: at level $n$, pulling back multiplication by $n$ on
$J_S$ gives the abelian covering with monodromy
$\pi\twoheadrightarrow\pi^{\rm ab}/n$.  This assertion follows alternatively
from the degree-zero divisor description of $J_S$: the compact part is
$T(J_X)$ and each puncture contributes its tame inertia generator, with the
single sum relation.  Both constructions give the exact sequence
\[
 0\longrightarrow
 \bigl(\widehat{\mathbb Z}(1)[S]/\widehat{\mathbb Z}(1)\bigr)
 \longrightarrow \pi^{\rm ab}\longrightarrow T(J_X)\longrightarrow0.
\]
The arithmetic comparison sends the difference of point sections to the
Kummer class of $[x-y]\in J_S(F)$.  A Kummer class vanishes precisely when
this element belongs to every $nJ_S(F)$: the finite-level statement is the
exact sequence for $[n]:J_S\to J_S$, and compatible choices follow from the
inverse system of finite nonempty division-point sets.  Thus the comparison
is proved in both directions.

For positive genus its image in $J_X(F)$ is divisible.  The arithmetic
input of Part A says that this image is zero.  A principal divisor $x-y$
with $x\ne y$ would give a degree-one morphism $X\to\mathbb P^1$; such a
morphism between smooth proper curves is an isomorphism, contrary to
$g(X)\geq1$.

In genus zero the existence of $x$ identifies $X$ with $\mathbb P^1$.
The divisor description gives
$J_S=\operatorname{Res}_{S/F}\mathbb G_m/\mathbb G_m$ and
$J_S(F)=\mathcal O(S)^*/F^*$; the equality on rational points follows also
from $H^1(F,\mathbb G_m)=0$.  The class of $x-y$ is represented by the tuple
$(h(a))_{a\in S}$, up to inversion, which is irrelevant for divisibility.
Take a finite Galois splitting field $L/F$ of $S$.  After restriction to $L$,
divisibility and Part A(iv) show that $h(a)/h(b)$ is a root of unity for
all $a,b$.  These roots of unity already belong to $F$.  Write
$h(a)=t\zeta_a$ with $\zeta_a\in\mu_\infty$.  Equivariance implies that
$\lambda(\sigma)=\sigma(t)/t$ is a character
$G=\operatorname{Gal}(L/F)\to\mu_\infty$.
Let $n=|G|$.  Divisibility over $F$ supplies an equivariant tuple $(b_a)$
and $c\in F^*$ with $b_a^n=c h(a)$.  Then $b_a/b_b$ is a root of unity,
so the same argument supplies a character $\psi:G\to\mu_\infty$ with
$\psi^n=\lambda$.  Every character of $G$ has exponent dividing $n$.
Consequently $\lambda=1$, all $h(a)$ belong to $F$, and the inverse
$F$-M\"obius transformation shows $a\in F$ for every cusp.  Conversely,
when $S$ is split and these ratios are roots of unity, every coordinate of
the class in the split torus has roots of every order in $F$, so the class
is divisible.  The parametrization for three punctures follows by solving
$x=ay$, $1-x=b(1-y)$; the listed exclusions are exactly those ensuring
$x,y\notin\{0,1\}$ and $x\ne y$.
\end{proof}

% Further statements are appended below the cyclotomic-base discussion.

\paragraph{Statement D2 (Proved).}\label{res:D2}
\passage{section 3; pp.30-31; scans 61-63}

For every finite extension $F'/F$ and every hyperbolic curve $U/F'$, the
map from $U(F')$ to geometric conjugacy classes of full continuous sections
is injective.  The same holds over every finitely generated extension of
$F$.  These assertions make no assertion of surjectivity.

\begin{proof}
We give the finite-cover comparison, including its descent.  Suppose two
point sections have the same class.  Conjugate the chosen lift of the
second geometric point so that their actual homomorphisms
$s:G_{F'}\to E_U$ agree.  Choose a finite-index characteristic subgroup
$H\subseteq\pi$ whose covering has compactification of positive genus.
Such a subgroup exists: for positive genus take $H=\pi$; in genus zero
choose three punctures over the algebraic closure and pull back the
Fermat cubic cover $w^3+v^3=1$ by a M\"obius coordinate carrying them to
$0,1,\infty$.  Removing the other punctures leaves a connected finite
\'etale covering of $U_{\overline F'}$ of genus one.  Taking its Galois
closure and then the intersection of all subgroups of the resulting
index produces an open characteristic subgroup contained in the original
one.  There are finitely many such subgroups because $\pi$ is finitely
generated.  The induced finite map of compactifications has positive-genus
target, and Riemann--Hurwitz gives positive genus for its source.

The open subgroup $H s(G_{F'})\subseteq E_U$ defines a connected finite
\'etale covering $V\to U$ over $F'$.  It has geometrically connected fibre
because its projection to $G_{F'}$ is surjective and its geometric subgroup
is exactly $H$.  Both geometric lifts chosen above are fixed by the common
$s(G_{F'})$, so descend to $F'$-points $\widetilde x,\widetilde y\in V(F')$.
Their sections in $E_V=H s(G_{F'})$ are the same actual homomorphism.
The morphism from $V$ to its smooth proper compactification carries them
to points with the same abelianized section.  Statement D1 in positive
genus gives $\widetilde x=\widetilde y$, hence $x=y$.

For finitely generated extensions we use the same argument and the AVKF
property established in Part A/E.  The exact external statements are
Tsujimura \cite[Proposition 1.2(ii), Proposition 4.2]{D_Tsujimura2022}:
finite generation preserves AVKF, and geometric points of a hyperbolic
curve over an AVKF field are separated by their germs of point sections.
Hoshi--Mochizuki--Tsujimura \cite[Definition 3.3(vi), Proposition 6.4]{D_HMT}
records the same injectivity for AVKF subfields of $\overline{\mathbb Q}$.
We have supplied the comparison rather than assuming section surjectivity.
\end{proof}

\paragraph{Statement D3 (Proved).}\label{res:D3}
\passage{section 3; p.30; scan 62}

Put $F=\mathbb Q(\mu_\infty)$ and $\xi=\zeta_6$.  For all positive
integers $n,m$ and every $\zeta\in\mu_n$, the finite \'etale covering of
$U_{0,3}$ defined by
\[
 w^n=z,\qquad u^m=1-\zeta w
\]
has isomorphic $G_F$-sets as fibres over $\xi$ and $\xi^{-1}$.
The bound $1\leq n,m\leq12$, with every $\zeta\in\mu_n$, is checked
exactly in \texttt{code/D\_coverings.py}; the argument below proves the
unbounded assertion.  The restriction $\zeta^n=1$ ensures that its extra
branch value is $z=1$, so that this is indeed a covering of the specified
three-punctured line.

\begin{proof}
Every $w$ in either fibre is a root of unity, hence belongs to $F$.
The map $w\mapsto w'=\zeta^{-2}w^{-1}$ bijects the first fibre of $w^n=z$
with the second.  We have the exact identity
\[
 1-\zeta w'=-\frac{1-\zeta w}{\zeta w}.
\]
Choose $\lambda_w\in\mu_\infty\subset F$ with
$\lambda_w^m=-(\zeta w)^{-1}$.  Then
$(w,u)\mapsto(w',\lambda_w u)$ is a $G_F$-equivariant bijection of the
finite fibres.  It commutes with the second-stage $\mu_m$-action.  This
comparison applies to these covers; the cover in D4 supplies a further
finite Galois-set invariant.
\end{proof}

\paragraph{Statement D4 (Proved).}\label{res:D4}
\passage{section 3; p.30; scan 62}

The full point sections at $\zeta_6$ and $\zeta_6^{-1}$ over
$F=\mathbb Q(\mu_\infty)$ are separated by an explicit connected finite
\'etale covering of degree $81$.  In particular, the smallest instance
specified in Part D(c) is decided.

\begin{proof}
Let $E$ be the smooth Fermat cubic $w^3+v^3=1$, with the flex
$(1:-1:0)$ as origin, and let $\varphi:E\to\mathbb P^1$ be $z=w^3$.
It has degree $9$ and ramifies only above $0,1,\infty$.  The formulas
\[
 X=\frac{12}{w+v},\qquad Y=\frac{36(w-v)}{w+v}
\]
identify its smooth projective model with $Y^2=X^3-432$; their inverse is
$w=(36+Y)/(6X)$, $v=(36-Y)/(6X)$.  Fix $a=\zeta_{18}$,
$\omega=a^6$, and $P=(a,a^{-1})\in E(F)$ in Fermat coordinates.
Our program performs the group law in the exact field
$\mathbb Q[a]/(a^6-a^3+1)$ and verifies
\[
 3P=(0,24a^3-12)\ne O,\qquad9P=O.
\]
The nine points above $\xi^{-1}$ are
$Q_{ij}=(\omega^i a^{-1},\omega^j a)$, $0\leq i,j<3$.
For each of them it verifies
\[
 3(Q_{ij}-P)=3P,\qquad9(Q_{ij}-P)=O.
\]
Thus all nine differences have exact order $9$.  Above $\xi$, the
corresponding differences have orders $1,3,\ldots,3$.
These are finite identities checked using rational arithmetic; the output
prints every pair $(i,j)$ and every triple.

We next prove the arithmetic bound that turns these identities into a
Galois-set separation.  The automorphism
$(X,Y)\mapsto(\omega X,Y)$ gives an action of
$\mathbb Z[\omega]$ on $E$, defined over $\mathbb Q(\omega)$.
Its action on $T_3E$ is faithful
\cite[II, Proposition 6.7]{D_MilneEC}.
The $3$-adic Tate module is a free rank-one module over
$\mathbb Z_3[\omega]$: it is free of rank two over $\mathbb Z_3$,
$T^2+T+1$ is irreducible over $\mathbb Q_3$, and
$\mathbb Z_3[\omega]$ is the valuation ring of the ramified quadratic
extension $\mathbb Q_3(\omega)$ (the element $\omega-1$ satisfies the
Eisenstein polynomial $T^2+3T+3$).  The module is finitely generated and
torsion-free over this valuation ring, hence free; comparison of
$\mathbb Z_3$-ranks gives rank one.  For every $n\geq1$ the natural map
\[
 T_3E/3^nT_3E\ \xrightarrow{\sim}\ E[3^n](\overline{\mathbb Q})
\]
is equivariant both for Galois and for the endomorphisms just specified.

The displayed Weierstrass equation has discriminant $-2^{12}3^9$, so
the elliptic curve has good reduction at $7$, and the program counts
$\#E(\mathbb F_7)=9$.  We use the following precise source statements:
Milne \cite[II, Proposition 6.3; V, Propositions 8.1 and 8.3]{D_MilneEC}
gives the rank-two prime-to-characteristic torsion module, unramifiedness
of the Tate representation at good primes away from its prime, and
$\operatorname{tr}(\operatorname{Frob}_p)=p+1-\#E(\mathbb F_p)$;
II, Proposition 6.4 identifies the Tate determinant with the degree of
an elliptic endomorphism.  Hence a Frobenius lift at $7$ acts on $T_3E$
with characteristic polynomial $T^2+T+7$.  As $7$ splits in
$\mathbb Q(\omega)$, a Frobenius lift $g$ in a decomposition group at $7$
restricts to the identity on $\mathbb Q(\omega)$.  It therefore commutes
with the specified endomorphism $\omega$ and, on the rank-one
$\mathbb Z_3[\omega]$-module, is multiplication by
$\alpha=1+3\omega$ or its conjugate.
Complex conjugation $c$ conjugates $\omega$ to $\bar\omega$, so the
commutator $cgc^{-1}g^{-1}$ acts by
$\bar\alpha/\alpha$ (or its inverse).  Each cyclotomic extension of
$\mathbb Q$ is abelian, so this commutator fixes $F$ pointwise.  It
therefore fixes every point defined over $F$.  Every $3$-power torsion point
$R\in E(F)$ belongs to some $E[3^n]$; the equivariant quotient above
therefore gives $(\bar\alpha/\alpha-1)R=0$, with the same conclusion
for the inverse scalar.  Since $\alpha$ is a unit in
$\mathbb Z_3[\omega]$ (its norm is $7$), this implies
$(\alpha-\bar\alpha)R=0$.
Since
\[
 (\alpha-\bar\alpha)(-1-2\omega)
 =(3+6\omega)(-1-2\omega)=9,
\]
we obtain $9R=0$.  The program checks the multiplication matrix
$\left(\begin{smallmatrix}3&-6\\6&-3\end{smallmatrix}\right)$,
its Smith form $(3,9)$ and the kernels modulo $9,27,81$.

Now let $M:E\to E$ be $R\mapsto[3]R+P$ and set
$W=M^{-1}(\varphi^{-1}(U_{0,3}))$.  The composite
$W\xrightarrow{M}\varphi^{-1}(U_{0,3})\xrightarrow{\varphi}U_{0,3}$
is defined over $F$: the curve and $\varphi$ are defined over $\mathbb Q$,
and $P\in E(\mathbb Q(\zeta_{18}))\subset E(F)$ defines the translation.
Over $U_{0,3}$, the map $\varphi$ is given by
$w^3=z$, $v^3=1-z$, with $w,v$ invertible; the two monic Kummer
equations give a finite \'etale algebra of rank $9$.
Multiplication by $3$ is finite \'etale of degree $9$ in characteristic
zero, and translation is an isomorphism, so the displayed composite is
finite \'etale of degree $81$.  Moreover $W$ is a nonempty open
subscheme of the geometrically integral elliptic curve $E$, hence is
geometrically connected.  Its fibre over $\xi$ contains the $F$-point
$R=O$, because $M(O)=P$ and $\varphi(P)=a^3=\xi$.
An $F$-point over $\xi^{-1}$ would satisfy
$3R=Q_{ij}-P$ for some $(i,j)$, forcing $R$ to have exact order $27$,
since $9R=3P\ne O$ and $27R=O$.  This contradicts the proved bound.
The two finite $G_F$-sets therefore have
different numbers of fixed elements.  Conjugate full sections induce
isomorphic $G_F$-sets on every finite \'etale fibre; this one disproves
conjugacy.  This last implication is the fibre-functor map from the
arithmetic fundamental group to the permutation group of the finite
covering, restricted to each actual section.
\end{proof}

\subsection{The proposed extension of the manuscript to cyclotomic bases}

\paragraph{Statement D5 (Refuted, with explicit counterexample).}\label{res:D5}
\passage{section 3; pp.27,30-31; scans 56,62-63}

The extension of Conjecture A to $F=\mathbb Q(\mu_\infty)$ is false,
even for proper hyperbolic curves.  The smooth projective normalization
$C/F$ of
\[
                         t^6+2y^3+4=0
\]
has genus $4$, has no $F$-rational point, and admits a continuous section
of $E_C\twoheadrightarrow G_F$.  The same counterexample persists after
adjoining any finite number of independent transcendental variables to
$F$.  Thus ``tous nos beaux \'enonc\'es'' (scan 62, manuscript p.~30) cannot include
the proposed surjectivity statement.  The point-injectivity and centralizer
statements proved separately remain valid.

\begin{proof}
The morphism $(t,y)\mapsto(t^2:y:1)$ maps the affine equation to
$D:X^3+2Y^3+4Z^3=0$ and extends to a morphism of smooth projective curves.
For $L=\mathbb Q(\mu_N)$ and a place above $2$, the ramification index
$e$ is a power of $2$: the odd-order roots of unity give unramified local
extensions, while the $2$-power roots give ramification degree
$\varphi(2^r)$.  In the integer-normalized valuation of $L$, the three
nonzero summands in the cubic have respective residues $0,e,2e$ modulo
$3$.  These residues are distinct.  Hence the term of least valuation is
unique and their sum cannot vanish.  Zero coordinates do not change the
argument for the remaining terms.  Thus $D(L)=\varnothing$ for every
$N$, and $C(F)=\varnothing$.  This is the case $p=2$, $a=b=c=1$ of
Clark \cite[Theorem 1]{D_Clark2006}, whose full hypothesis is
$p\equiv-1\pmod3$ and $p\nmid abc$; we have reproduced the valuation
proof for the example used here.

The function-field extension $F(t)\subset F(t,y)$ has degree three:
$-(t^6+4)/2$ has a simple zero and so is not a cube even over
$\overline F(t)$.  The six roots of $t^6+4$ are simple.  They are the six
branch values of index three; the pole at infinity has order six, so
infinity has three unramified preimages.  Riemann--Hurwitz gives
$2g-2=3(-2)+6(3-1)=6$, hence $g=4$.  The exact polynomial and valuation
residue checks are in \texttt{code/D\_cyclotomic\_counterexample.py}.

We give the section-existence argument, including the finite lifting
step, to avoid assuming any section conjecture.  First,
$\operatorname{Br}(L)=0$ for every finite extension $L/F$.
A central simple algebra over $L$ descends to a number field $k\subset L$.
Milne \cite[VII, Proposition 7.2]{D_MilneCFT} states that every Brauer
class of a number field is killed by a finite cyclic cyclotomic extension.
That extension of $k$ is contained in $L$, which contains all roots of
unity.  For clarity, the local-global input behind that proposition is
VIII, Theorem 4.2: the map from the Brauer group of a number field to the
direct sum of its local Brauer groups is injective.  III, Theorem 2.1
identifies the nonarchimedean local group with $\mathbb Q/\mathbb Z$ and
restriction with multiplication by the local degree.  Cyclotomic
extensions provide arbitrarily divisible unramified local degrees,
through $\mu_{p^d-1}$ at $p$; a finite compositum kills the finitely many
nonzero local invariants, and adjoining $i$ kills real invariants.

Next $H^2(G_F,A)=0$ for every finite continuous $G_F$-module $A$.
Here is the group-theoretic reduction.  It suffices to treat its
$\ell$-primary parts.  Let $P$ be a pro-$\ell$ Sylow subgroup of $G_F$.
For every open $V$ containing $P$, Kummer theory and the preceding
Brauer vanishing give $H^2(V,\mathbb F_\ell)=0$, since $\mu_\ell\subset F$.
Continuous cochains give
$H^2(P,\mathbb F_\ell)=\varinjlim_{P\subseteq V}H^2(V,\mathbb F_\ell)$:
a locally constant cochain factors through a finite quotient, and its
cocycle or coboundary identities extend to an open neighbourhood $PN$ of
$P$ for a sufficiently small open normal $N$.  A finite $P$-module has a
filtration with trivial $\mathbb F_\ell$ factors: reduce by powers of
$\ell$, and use the nonzero invariants of every finite $\ell$-group on a
nonzero vector space in characteristic $\ell$.  The latter follows by
induction from a central element of order $\ell$, whose operator $g-1$
is nilpotent.  The cohomology long exact sequences now give
$H^2(P,A)=0$.  Restriction to $P$ is injective on the $\ell$-primary
cohomology: vanishing on $P$ gives vanishing on some such $V$, and
corestriction composed with restriction multiplies by $[G_F:V]$, prime
to $\ell$.  The transfer identity is Milne
\cite[II, Proposition 1.30]{D_MilneCFT}; its profinite version follows
from the same finite-quotient cochains.  Thus the claimed vanishing holds.

We now solve a finite embedding problem $G_F\to A_0\leftarrow B_0$
with $B_0\twoheadrightarrow A_0$ surjective.  Replace $A_0$ by the
image of the first map and $B_0$ by its preimage.  Choose a subgroup
$H\subseteq B_0$ minimal among those surjecting onto $A_0$.  Its kernel
$N$ lies in the Frattini subgroup of $H$: a maximal subgroup not
containing $N$ would itself surject onto $A_0$.  Moreover $N$ is
nilpotent.  Indeed, for any Sylow subgroup $P_N$ of the normal subgroup
$N$, the finite Frattini argument gives
$H=N_H(P_N)N$; the Frattini property forces $N_H(P_N)=H$.
All Sylow subgroups of $N$ are therefore normal.  When $N\ne1$, its
centre contains a nonzero elementary abelian subgroup invariant under
$H$, and hence a minimal nonzero $H$-normal elementary abelian subgroup
$D$.  Induction on $|N|$ solves the problem for $H/D\to A_0$.
Pulling back $H\to H/D$ by that solution gives an extension of $G_F$
by the finite abelian module $D$.  Its class in $H^2(G_F,D)$ vanishes,
so it has a continuous splitting, yielding the required weak solution.
The last statement follows directly from a continuous set-theoretic
section and its two-cocycle; such a section exists for a finite-kernel
map of profinite groups by using finitely many clopen cosets.

Every profinite extension of $G_F$ with finite kernel consequently
splits.  In detail, choose an open normal subgroup $V$ of the extension
group meeting its finite kernel trivially.  The extension is the fibre
product of the finite map $E/V\to G_F/p(V)$ with
$G_F\to G_F/p(V)$, so the solved finite embedding problem supplies its
splitting.  Finally the geometric fundamental group of $C$ is finitely
generated and has a countable cofinal sequence
$\pi=N_0\supset N_1\supset\cdots$ of open characteristic subgroups.
Starting with the identity section of $E_C/N_0=G_F$, lift recursively
through $E_C/N_{r+1}\to E_C/N_r$, whose kernels are finite.
The compatible sections give a continuous homomorphism
$G_F\to\varprojlim E_C/N_r=E_C$ splitting the projection.  This is an
actual section-existence proof, although it does not select a canonical
section.

For $F(T_1,\ldots,T_d)$ the arithmetic fundamental-group extension of
the constant curve is the pullback by restriction to $G_F$, so the
section pulls back continuously.  An $F(T_1,\ldots,T_d)$-point would
give a rational map $\mathbb P^d_F\dashrightarrow C$; restricting to
lines and using Riemann--Hurwitz shows that it is constant, and its
constant value is an $F$-point.  This contradicts $C(F)=\varnothing$.
\end{proof}

\paragraph{Statement D6 (Proved).}\label{res:D6}
\passage{section 1; p.10; scan 22; section 3; p.30; scan 62}

The faithfulness assertion of manuscript Th\'eor\`eme 2 (scan 22,
manuscript p.~10) holds for finitely generated field extensions of
$F=K(\mu_\infty)$, including finite extensions of this $F$.
Explicitly, $F$-embeddings $f,g:M\hookrightarrow L$ inducing the same
outer continuous homomorphism $G_L\to G_M$ are equal.

\begin{proof}
As $\mu_n\subset F$, inner conjugation acts trivially on
$H^1(G_M,\mu_n)=\operatorname{Hom}_{\rm cont}(G_M,\mu_n)$.
The two outer maps therefore induce equal Kummer pullbacks for every
$n$, giving $f(a)/g(a)\in\bigcap_n L^{*n}$ for each $a\in M^*$.
Let $F_1$ be the relative algebraic closure of $F$ in $L$; it is finite
over $F$.  On a normal proper model of $L/F_1$, every prime-divisor
valuation of an infinitely divisible element is zero.  Such an element
and its inverse are global regular functions, hence constants in $F_1$.
Every root in $L$ of a constant is algebraic over $F_1$, so belongs to
$F_1$.  The arithmetic input in A(iv) now gives
$\bigcap_nL^{*n}=\mu_\infty$.

Choose an embedding of the countable characteristic-zero field $L$ into
$\mathbb C$.  Root-of-unity ratios give equal complex absolute values
for $f(a)$ and $g(a)$, and the same for $a+1$ and $a+i$.
The resulting equalities of squared absolute values give equal real
parts and equal imaginary parts of $f(a)$ and $g(a)$, respectively.
Here $f(i)=g(i)=i$ after naming the chosen image of $i$; the other image
$-i$ yields the same conclusion.  Thus $f(a)=g(a)$ for all $a$.
This argument also covers $a=-1,-i,0$ directly, since these belong to
$F$.  The comparison used throughout is the contravariant Kummer
pullback on $H^1$, so it respects the precise inner ambiguity in the
manuscript's groupoid formulation.
\end{proof}

\paragraph{Statement D7 (Refuted; exact replacement proved).}\label{res:D7}
\passage{section 1; p.15; scan 32; section 3; p.30; scan 62}

The abelianized point-injectivity assertion of manuscript Th\'eor\`eme 3
(scan 32, manuscript p.~15) fails over $F$ and over its finitely generated
extensions.  Already for $\mathbb G_m/F$, the distinct points $1$ and
$\zeta_6$ have identical full section classes.  For an embedding of an
$F$-variety into a semiabelian variety, its induced Kummer classes agree
exactly when their difference belongs to the intersection of all
multiples of the rational-point group.  For hyperbolic curves the more
precise replacements are D1 and D2.

\begin{proof}
The geometric fundamental group of $\mathbb G_m$ is
$\widehat{\mathbb Z}(1)$, and the difference of the two point sections
is the Kummer class of their ratio.  Each root of unity has compatible
roots of all orders in $F$, so the class of $\zeta_6$ is zero.  This
persists under every extension of $F$.  For a semiabelian variety the
multiplication-by-$n$ exact sequence gives the stated divisibility
criterion at each finite level.  Compatible finite division-point sets
give the full Tate-module criterion, as in the proof of D1.
\end{proof}

\paragraph{Statement D8 (Proved).}\label{res:D8}
\passage{section 1; pp.6-7; scans 13-15; section 3; p.30; scan 62}

Manuscript Th\'eor\`eme 1 (scans 13--15, manuscript pp.~6--7), with its stated
centralizers of open subgroups, holds for every finitely generated
extension $M/F$.  It holds for its geometric absolute Galois group as
well, with the trivial group understood in transcendence degree zero.

\begin{proof}
The field $F$ is an abelian extension of a number field and is Hilbertian;
the exact Hilbertian inputs and the centralizer argument are proved in
Statement B12.  A finite extension preserves this property.  In positive
transcendence degree a rational function field is Hilbertian, and finite
extensions preserve it.  Apply B12 to $M$ and, for its geometric group,
to the positive-transcendence-degree function field over the algebraic
closure of the constant field.  This gives centralizer triviality for
every open subgroup, not only the centre.
\end{proof}

\paragraph{Statement D9 (Proved replacement; Conjecture A refuted).}\label{res:D9}
\passage{section 3; pp.26-27,30; scans 54-56,62}

For the cyclotomic base $F$ and its finitely generated extensions, the
hyperbolic version of the Proposition of scan 54 and the centralizer
``pr\'ecisions'' of scan 56 hold with the same section and conjugacy
conventions as Parts A and C.  The excluded type $(0,2)$ remains a
counterexample to the first assertion of the Proposition.  Conjecture A
itself has the counterexample D5.

\begin{proof}
The finite-extension AVKF property used for packet separation holds by
D2.  The Tate-module vanishing for all finite extensions is the
arithmetic input proved in Part A/E; since the base contains all roots
of unity, a cyclotomically normalized inertia subgroup is centralized.
Apply the already proved inertia-recognition and unique-cusp argument
of Part A over this base and the splitting-class comparison of Part C.
Every comparison is a base extension of the stated finite \'etale
covering or the continuous inclusion of the corresponding open Galois
subgroup, so all its hypotheses persist.  The $(0,2)$ calculation is
literal: for $\mathbb G_m$ the decomposition extension at either cusp
is the full fundamental-group extension.  Its splitting packet
therefore contains every point section, including the pair in D7.
\end{proof}

\subsection{Poincar\'e coverings and the continuity clause}

\paragraph{Statement D10 (Proved).}\label{res:D10}
\passage{section 3; p.30; scan 62}

Let $k$ be an algebraically closed field of characteristic zero and let
$U,V$ be hyperbolic curves over $k$.  Their universal pro-\'etale
coverings are isomorphic as pro-curves over $k$ exactly when there is a
connected curve $W$ with finite \'etale surjections to both $U$ and $V$.
The completion of a universal covering is canonical and functorial for
semilinear isomorphisms.  More precisely, it is the inverse system of
smooth proper compactifications of its finite levels; the set used in
the fixed-point assertions is the set of compatible geometric points
\[
                       X'=\varprojlim_i X_i(k).
\]
Thus a generic scheme point is not being counted as a fixed point in
Conjectures B or B$'$.

\begin{proof}
The transition maps between finite levels are finite, hence affine.
Their scheme limit exists by Stacks, Lemma 32.2.2
\cite[Tag 01YV]{D_Stacks}.  A morphism from this limit to a curve of
finite presentation factors through a finite level.  The precise result
is Stacks, Proposition 32.6.1(3), Tag 01ZB: for an inverse system of
quasi-compact quasi-separated schemes with affine transitions, morphisms
to a locally finitely presented target are the filtered colimit of the
finite-level morphism sets.  The same statement identifies two such
morphisms after a sufficiently fine level when their limit morphisms
coincide.

Let $\alpha:\widetilde U\simeq\widetilde V$ be an isomorphism.  Factor
$p_V\alpha$ through a finite cover $U_i$ of $U$, obtaining $f:U_i\to V$.
Factor $p_{U_i}\alpha^{-1}$ through a finite cover $V_j$ of $V$, obtaining
$g:V_j\to U_i$.  After increasing $j$, the equality of limit morphisms
gives $fg=q_j$, where $q_j:V_j\to V$ is the finite \'etale structure
map.  In particular $f$ and $g$ are nonconstant.

Here is the finiteness step, which is necessary for the asserted
equivalence.  A nonconstant map of smooth curves is quasi-finite and
separated.  Zariski's main theorem, Stacks, Lemma 37.43.3, Tag 05K0,
factors $f$ as an open immersion $U_i\subset T$ followed by a finite map
$T\to V$; replace $T$ by the reduced closure of $U_i$ in $T$.
The composite $V_j\to T$ is proper over $V$: its graph is closed because
$T/V$ is separated, and projection from $V_j\times_V T$ is proper
because $V_j/V$ is finite.  Its image is closed and dense in $T$, so is
all of $T$, while it is contained in $U_i$.  Consequently $T=U_i$ and
$f$ is finite.  Since $g$ is surjective and $fg$ is \'etale, the
derivative of $f$ vanishes nowhere.  A finite map of smooth curves is
flat, as its local modules over the target discrete valuation rings
are torsion-free.  Hence $f$ is \'etale.  The curve $U_i$ is the required
common cover.

Conversely, the pointed finite \'etale covers of a connected finite
\'etale cover $W\to U$ form a cofinal subsystem of the covers of $U$:
for a given cover take the connected component, containing the chosen
geometric lift, of its fibre product with $W$.  Thus both $W\to U$ and
$W\to V$ identify their universal pro-coverings, with the usual deck
ambiguity when the lifts are forgotten.

A finite map between smooth curves extends uniquely between their
smooth proper models.  The exact algebraic statement is the equivalence
between transcendence-degree-one function fields and normal proper
curves with nonconstant morphisms, Stacks, Theorem 53.2.6, Tag 0BY1.
Apply it to the finite-level maps just constructed.  Equality after
refinement ensures compatibility, and the inverse isomorphism gives the
inverse map of completed systems.  This also proves independence of the
initial presentation, using the common cover.  Replacing $k$ by its
image under an automorphism gives the same construction for a semilinear
isomorphism.  On compatible geometric points these maps define exactly
the claimed functorial completion.
\end{proof}

\paragraph{Statement D11 (Proved).}\label{res:D11}
\passage{section 3; pp.29-31; scans 59,62-63}

Let $K$ have characteristic zero and let $\mathcal U$ be a Poincar\'e
pro-curve over $\overline K$.  Give its semilinear automorphism group
the topology whose identity neighbourhoods are the simultaneous
stabilizers of finitely many rational functions.  Every semilinear
$G_K$-action for which each rational function has open stabilizer
preserves, after restriction to an open normal subgroup $G_{K_i}$, a
finite level $U_i$.  This level has a $K_i$-form, and the restricted
action is supplied by a continuous section of its arithmetic
fundamental-group extension.  One may start with any specified finite
level in this assertion.

The families of assertions B$'$ with this continuity clause and B with
$\Gamma^\circ$, each quantified over all finite extensions of $K$ and
all hyperbolic curves there, are equivalent.  At one fixed base $K$,
the precise implications are
\[
 \{\mathrm B_L:\ [L:K]<\infty\} \Longrightarrow\
 \mathrm B'_K\text{ (continuous)}\ \Longrightarrow\ \mathrm B_K.
\]
These are comparisons of the assertions, without assuming their
existence conclusions.  For the cyclotomic base in D5, even the
continuous form has a counterexample.

\begin{proof}
Let $\Omega$ be the union of the function fields of the finite levels,
and select one such field $L_0=\overline K(f_1,\ldots,f_r)$.
Choose equations for its smooth compactification and boundary, with
coefficients in a finite extension $K_0/K$.  Enlarge the finite list of
functions when necessary to realize these equations as the chosen
model.  Intersect the open stabilizers of the $f_j$ and the open
subgroup fixing $K_0$, and then take its normal core in $G_K$.
This gives an open normal $G_{K_i}$.  On $L_0$ its action fixes the
$f_j$ and acts on the coefficients in the prescribed way.  Thus it
is the standard semilinear action on that $K_i$-model, preserving its
boundary.  Equality of function-field actions implies equality of the
morphisms to this smooth separated curve.  The projection
$\mathcal U\to U_i$ is therefore equivariant.

The automorphisms of $\Omega$ fixing $K_i(U_i)$ form $E_{U_i}$: a
finite separable subextension is admitted precisely when its
normalization is \'etale over $U_i$, and adjoining the algebraic
constants gives the geometric universal cover.  This identifies the
actual group, including its projection to $G_{K_i}$, rather than merely
an outer action.  The given semilinear action therefore gives a
homomorphism $s_i:G_{K_i}\to E_{U_i}$ splitting that projection.
The Krull topology on this group is generated by stabilizers of finite
sets of elements of $\Omega$.  The hypothesis makes $s_i$ continuous.
In particular its image is compact.  Conversely a continuous section
has open stabilizers on every element of $\Omega$, since each element
lies in a finite arithmetic covering.  This proves both topological
comparisons directly.

Apply the fixed-point assertion B over $K_i$ to obtain its unique
$G_{K_i}^\circ$-fixed geometric pro-point.  The subgroup
$G_{K_i}^\circ=G_{K_i}\cap G_K^\circ$ is normal in $G_K$, so uniqueness
makes that point fixed by all of $G_K$.  Any $G_K^\circ$-fixed point
is also $G_{K_i}^\circ$-fixed, proving uniqueness in B$'_K$.
The reverse implication at a fixed base is the ordinary
section-to-action construction.  Applying these two implications at
every finite extension proves the displayed family equivalence.
Finally, the section of the proper curve in D5 has no fixed geometric
pro-point, because its image on the curve would be an $F$-rational
point.  Its action is continuous by the construction just proved.
\end{proof}

\paragraph{Statement D12 (Proved for proper curves).}\label{res:D12}
\passage{section 3; pp.30-31; scans 62-63}

Let $U/\overline K$ be proper hyperbolic, choose an embedding
$\overline K\hookrightarrow\mathbb C$, and write
$U(\mathbb C)=\mathbb H/\Lambda$.  Its pro-curve automorphism group is
the Belyaev completion of $\operatorname{Comm}(\Lambda)$, with
$\widehat\Lambda$ a compact open subgroup.  In particular,
\[
 \operatorname{Aut}_{\overline K}(\widetilde U)
       =\widehat\Lambda\operatorname{Comm}(\Lambda),\qquad
 \operatorname{Aut}_{\overline K}(\widetilde U)/\widehat\Lambda
       \simeq\operatorname{Comm}(\Lambda)/\Lambda
\]
as groups in the first formula and as coset sets in the second.
For a non-arithmetic uniformizing group this completion is profinite;
for an arithmetic uniformizing group it is locally compact and
noncompact.

\begin{proof}
The opened source is D\v zambi\'c--Gonz\'alez-Diez
\cite[Theorem 7 and Proposition 3]{D_DzambicGonzalezDiez}.
Its hypothesis is a compact Riemann surface of genus at least two,
and its conclusion identifies the holomorphic automorphisms of the
inverse-limit solenoid with the completion of the commensurator in
the permutation action on cosets of its finite-index subgroups.
Proposition 3 states that $\widehat\Lambda$ is compact open and gives
the two displayed formulas.  The arithmeticity criterion used in its
Corollary 11 is
$[\operatorname{Comm}(\Lambda):\Lambda]=\infty$ precisely in the
arithmetic case; the source records this Margulis criterion explicitly
in its Introduction, p.~3.

Here is the comparison with the algebraic group used here.  Riemann
existence, SGA~1, XII, Theorem 5.1 \cite{D_SGA1}, sends each finite
topological cover to its finite \'etale algebraic cover, and finite
\'etale covers are unchanged by extension between algebraically closed
fields of characteristic zero (SGA~1, XIII, Corollary 2.9).
An algebraic pro-automorphism therefore acts holomorphically on the
inverse-limit solenoid.  Conversely, the source's formula expresses
every holomorphic automorphism as a deck element times a commensurator
element.  The latter gives an isomorphism between two finite covers
$\mathbb H/\Lambda_1$ and $\mathbb H/\Lambda_2$.  It is algebraic by
SGA~1, XII, Corollary 4.5, and
is defined over $\overline K$: the isomorphism scheme between two
proper genus-at-least-two curves is finite, so extending an algebraically
closed base cannot add its geometric points.  These two finite covers
and their isomorphism induce a pro-automorphism by the cofinality map
in D10.  This proves surjectivity of the comparison, while agreement
on finite-level function fields proves injectivity.  Stabilizers of
finite levels give the same compact open subgroups on the two sides.

When the commensurator has finite index over $\Lambda$, completing for
the finite-index-subgroup topology is its usual profinite completion.
When that index is infinite, the quotient by the compact open subgroup
$\widehat\Lambda$ is infinite discrete.  A compact group could not have
such a quotient.  This proves the two topological conclusions from
the stated criterion.
\end{proof}

\paragraph{Statement D13 (Open in the unrestricted cases; proved partial results).}
\label{res:D13}
\passage{section 3; pp.30-31; scans 62-63}

The common-cover and completion assertions of scan 62, and the
continuity descent in Part D(e), are proved in D10--D11.  D12 gives the
automorphism description for proper curves.  An action over a number
field that normalizes no profinite deck group has not been constructed
or excluded here.  The affine commensurator-completion assertion, beyond
the finite-level comparisons of D10, has not been proved here either.
The smallest instance retained for this part is the Poincar\'e curve of
$\mathbb P^1\setminus\{0,1,\infty\}$ over $\overline{\mathbb Q}$,
with its full semilinear $G_{\mathbb Q}$-actions.

To state the missing assertion exactly, fix its usual continuous
tangential action $\rho_0$, put
$A=\operatorname{Aut}_{\overline{\mathbb Q}}(\mathcal U)$, and let
$\mathscr P$ be the profinite deck groups of finite algebraic curve
quotients of $\mathcal U$.  Every abstract semilinear action over
$G_{\mathbb Q}$ is uniquely a function $a:G_{\mathbb Q}\to A$ satisfying
\[
 a_{gh}=a_g\,\rho_0(g)a_h\rho_0(g)^{-1},
 \qquad \rho(g)=a_g\rho_0(g).
\]
The unresolved construction or exclusion asks whether such a function
can satisfy
\[
 \text{for every }P\in\mathscr P\text{ there exists }g\in G_{\mathbb Q}
 \text{ with }a_g\rho_0(g)P\rho_0(g)^{-1}a_g^{-1}\ne P.
\]
These are literal equations in the indicated automorphism group.
Their equivalence with the requested action follows by multiplying
$\rho(g)\rho(h)$ and by the definition of normalization; it does not
assume any continuity or any section conjecture.
For this Poincar\'e curve B$'$ asserts a unique
$G_{\mathbb Q}^\circ$-fixed compatible geometric point for every such
semilinear action.  D11 settles its continuous comparison with B;
the unrestricted assertion above is the outstanding part of D(e).

% END sections/D_agent.tex

\section{E. Other subgroups and local fields}

% BEGIN sections/E_arithmetic_agent.tex
\subsection{Finitely generated fields and the subgroup hypothesis}
\paragraph{Statement AE1 (Proved; Part E(a), arithmetic and centralizer assertions).}\label{res:AE1}
\passage{Section 3, manuscript pp. 27,30, scans 56,61–62}

The manuscript passages are cote 149, section 3, manuscript pp.~27 and 30, scans 56 and 61--62.
Let $K/\mathbb Q$ be finitely generated, $K'/K$ finite, and $A/K'$ an abelian variety. Then $A(K'(\mu_\infty))_{\rm tors}$ is finite. Every field $L$ finitely generated over a finite extension of $\mathbb Q(\mu_\infty)$, together with every finite extension of $L$, has no nonzero divisible abelian-variety points and has zero invariants in every abelian-variety Tate module. Statements A1--A7, in their indicated centralizer and inertia formulations, consequently extend to finitely generated $K/\mathbb Q$.

\emph{Proof of torsion finiteness.} Replace $K$ by $K'$ and let $k$ be its algebraic constant field over $\mathbb Q$. It is a number field and $K/k$ is regular. Choose a smooth geometrically integral variety $S/k$ with function field $K$, shrink it, and extend $A$ to an abelian scheme $\mathcal A/S$. The extension of the group law, inverse, identity, properness and smoothness uses finite presentation and a common open subset on which the corresponding identities and properties hold. Write $k^{\rm cyc}=k(\mu_\infty)$, so $K(\mu_\infty)=Kk^{\rm cyc}$, and $S_{k^{\rm cyc}}$ is geometrically integral.

For every $n$, $\mathcal A[n]$ is finite etale over $S$, since the characteristic is zero. A $Kk^{\rm cyc}$-rational point of $A[n]$ extends uniquely to a section over the normal scheme $S_{k^{\rm cyc}}$: its closure in the finite etale scheme is a finite birational component, and normality identifies that component with the base. Two such sections agreeing at one geometric point agree everywhere, since their equalizer is both open and closed and the base is connected. Choose a closed point $x$ of $S$ with residue field $k_x$, and one point above it after extending constants. Evaluation gives, compatibly with $n$ and the group laws, an injection
\[
 A(Kk^{\rm cyc})_{\rm tors}\hookrightarrow
 \mathcal A_x(k_x(\mu_\infty))_{\rm tors}.
\]
The target is finite by Ribet's Appendix Theorem 1 with the number field $k_x$. This proves the asserted finiteness without using a specialization theorem that requires finite generation of a Mordell--Weil group.

\emph{Proof about divisible points.} A field $L$ as in the statement is $K_0(\mu_\infty)$ for a finitely generated $K_0/\mathbb Q$: descend its finitely many generators and the finitely many coefficients in their defining relations. The same holds after any finite extension of $L$. Descend an abelian variety and a proposed nonzero divisible point $P$ to $K_0$, then spread them to $\mathcal A/S$ and a section $P$ as above. Shrink $S$ so that $P$ does not meet the identity section. For every $n$ the division scheme $[n]^{-1}(P)$ is finite etale over $S$. A division point over $L$ extends over $S_{k^{\rm cyc}}$ by the same normality argument. Evaluation at $x$ would therefore make the nonzero point $P_x$ divisible by every integer over $k_x(\mu_\infty)$, contradicting A1. Hence there is no such $P$. Invariant Tate vectors vanish either by torsion finiteness or by the compatible nonzero divisible torsion points they would provide.

For completeness, the multiplicative assertion of A1 extends too. The image of the cyclotomic character of $K_0$ is open, since its algebraic constant field is a number field. Its kernel fixes $K_0(\mu_\infty)$, so the same cocycle argument as in A1 shows that an infinitely divisible element has a power with zero Kummer class over $K_0$. The kernel of the latter Kummer map is trivial: on a normal proper model over the algebraic constant field, divisibility by every $n$ forces every divisorial valuation to be zero; a rational function with neither zeros nor poles is algebraic over that constant field. The number-field argument in A1 then applies. Splitting a torus over a finite extension proves the same vanishing of homomorphisms to abelian Tate modules.

Finally every finite geometric covering appearing in A4 descends to a finite extension of this finitely generated $K$, and its compactified Jacobian is an abelian variety there. The torsion finiteness just proved supplies exactly the invariant vanishing used in (A.1). The maps, subgroup containments, compactness and boundary calculations in A2--A7 are unchanged. Thus their proofs transfer, with these specified arithmetic input maps, to this generality.\hfill$\square$

\paragraph{Statement AE2 (Proved; H1 with the necessary zero-centralizer qualification).}\label{res:AE2}
\passage{Section 3, manuscript pp. 30, scans 61–62}

This refines H1 and the subgroup question of cote 149, section 3, manuscript p.~30, scans 61--62.
Let $K$ be a number field, $U/K$ a hyperbolic curve, $s:\Gamma_K\to E_U$ a continuous section, and $\Delta\subset\Gamma_K$ a closed subgroup satisfying
\[
 T(A)^{\Delta\cap\Gamma_{K'}}=0
 \quad\text{for every finite }K'/K\text{ and every abelian variety }A/K'.\tag{E.1}
\]
Write $C_\Delta=C_\pi(s(\Delta))$ and
$P_\Delta=\{\ell:\chi_\ell|_\Delta=1\}$. Every nonzero $C_\Delta$ has a unique associated lifted cusp $c$, fixed by $s(\Delta)$, and
\[
 C_\Delta=I_c^{\chi(\Delta)}
 =\prod_{\ell\in P_\Delta}\mathbb Z_\ell(1).
\]
For $P_\Delta\ne\varnothing$, existence of an $s(\Delta)$-fixed cusp is equivalent to $C_\Delta\ne1$ and that fixed cusp is unique. For $P_\Delta=\varnothing$, $C_\Delta=1$ for every section, including cuspidal sections. When $\Delta$ is normal, existence of an $s(\Delta)$-fixed cusp implies that $s$ itself is cuspidal at that unique cusp, including the case $P_\Delta=\varnothing$.

\emph{Proof.} Use the construction in A4 for $z\in C_\Delta$ and $H=N\langle z\rangle$. The stabilizer of $H$ in $\Gamma_K$ is open; write it as $\Gamma_{K'}$. The image of $z$ in the Tate module of the compactified covering is fixed by $\Delta\cap\Gamma_{K'}$, and (E.1) makes it zero. Thus A3 puts $z$ in inertia, and A5 puts all of $C_\Delta$ in one inertia subgroup whenever $C_\Delta\ne1$. A nonidentity $z\in C_\Delta\cap I_c$ and A2 show that $s(\Delta)$ fixes $c$. The action on $I_c$ is $\chi|_\Delta$, so the fixed subgroup of $I_c$ is exactly $C_\Delta$. Computing each factor in the integral domain $\mathbb Z_\ell$ gives the displayed product. If $P_\Delta$ is nonempty, every fixed cusp contributes a nonzero such product, so A5 and A2 give its uniqueness. If $P_\Delta$ is empty, the same calculation makes a nonzero centralizer impossible.

Suppose now that $\Delta$ is normal and put $L=\overline K^{\Delta}$. We verify the additional assertion without obtaining information from a zero centralizer. Every finite extension $L'/L$, every abelian variety over $L'$ and every one of its rational points descend to some finite $K'/K$ with $L'=LK'$. The group $\Delta\cap\Gamma_{K'}$ is normal in $\Gamma_{K'}$. For a point divisible over $L'$, choose compatible division points there. Its Kummer cocycle is zero on this normal subgroup. Comparing the cocycle on $hg$ and $g(g^{-1}hg)$ puts its value at $g$ in $T(A)^{\Delta\cap\Gamma_{K'}}=0$. The point is therefore divisible over $K'$ and is zero by reduction modulo powers of a good finite prime, as proved in A1. Thus all finite extensions of $L$ have no nonzero divisible abelian-variety points.

The continuous base-change homomorphism $\pi_1(U_L)\to E_U$ identifies its source with the inverse image of $\Delta$: finite etale covers descend to finite intermediate fields, and their subgroup identifications pass to this inverse limit. Thus $s|_\Delta$ is an actual section for $U_L/L$. Stix Theorem 14, with this hypothesis verified as in the second proof of A6, now gives uniqueness of an $s(\Delta)$-fixed lifted cusp. Normality means that $s(\Gamma_K)$ permutes the fixed-cusp set; a singleton is fixed. Hence $s$ is cuspidal at that cusp. Alternatively, for $C_\Delta\ne1$, normality preserves $C_\Delta$ and A2 immediately puts $s(\Gamma_K)$ in $D_c$.\hfill$\square$

The original H1 conclusion that nontrivial centralizers require triviality of the entire cyclotomic character is replaced by the primewise product above. There is no inference from $C_\Delta=1$ to absence of a cusp when $P_\Delta$ is empty. For arbitrary nonnormal $\Delta$, (E.1) alone does not supply the normal-subgroup cocycle argument just given; uniqueness of fixed cusps in this zero-centralizer case is not asserted here. The precise remaining assertion would be that such a subgroup can fix at most one lifted cusp even when all $\chi_\ell|_\Delta$ are nontrivial. Our proof establishes the complete assertion needed for the original, normal-subgroup H1 and for every nonzero centralizer.

\paragraph{Statement AE3 (Proved; arithmetic instances of (E.1), Part E(b)).}\label{res:AE3}
\passage{Section 3, manuscript pp. 30, scans 62}

These are the arithmetic instances requested for H1 and cote 149, section 3, manuscript p.~30, scan 62.
For a number field $K$, condition (E.1) holds for the absolute Galois groups of $K(\mu_{\ell^\infty})$, of $K(\mu_\infty)$, and of $K(K^{1/\infty})$. It also holds for the Kummer fields in Murotani--Ozeki Proposition 3.15(2). For $\Delta=\Gamma_{K^{\rm ab}}$, condition (E.1), with all of its quantifiers, holds exactly when $K$ contains no CM subfield.

\emph{Sources and reductions.} The first two cases follow from A1, since $\Gamma_{K(\mu_{\ell^\infty})}\supset\Gamma_{K(\mu_\infty)}$, and intersecting with $\Gamma_{K'}$ preserves that inclusion. For the remaining cases we use the following precise statements from the opened 17 January 2025 manuscript \cite{MurotaniOzeki2025}. Its Proposition 3.15(2) has a characteristic-zero field $k$, a Galois extension $L/k$ containing $\mu_\infty$, open images $\chi_p(\Gamma_k)$ for every prime $p$, and $L$ abelian-variety Kummer-faithful (AVKF); it concludes that every subfield of $L(k^{1/\infty})$ is AVKF. Theorem 3.8, for a number field $K$, equates: absence of a CM subfield of $K$; finiteness of $A(K^{\rm ab})_{\rm tors}$ for every $A/K$; and the AVKF property of $K^{\rm ab}$. AVKF means that every abelian variety over every finite extension has no nonzero divisible rational point.

An AVKF field $M$ has $T_\ell(A)^{\Gamma_{M'}}=0$ for every finite $M'/M$: a nonzero invariant vector would give a nonzero divisible $\ell$-primary torsion subgroup, whose elements are also divisible by integers prime to $\ell$. For an abelian variety defined over $K'$, take $M'=MK'$. This proves (E.1). Apply Proposition 3.15(2) to $k=K$, $L=K(\mu_\infty)$, for which A1 verifies AVKF and the cyclotomic images are open. This gives the asserted full radical field as well as all the stated general instances. No claim of finiteness of the total torsion group over $K(K^{1/\infty})$ is required.

For $K^{\rm ab}$, Theorem 3.8 and the preceding implication prove the forward assertion when $K$ contains no CM subfield. In the opposite case, the proof of that theorem explicitly constructs a finite $K'\subset K^{\rm ab}$ and a nonzero CM abelian variety $A/K'$ all of whose torsion is defined over $K^{\rm ab}$. For a CM subfield $k\subset K$, the inclusion $Kk^{\rm ab}\subset K^{\rm ab}$ follows because base change of an abelian Galois extension remains abelian; base-changing that construction and taking the compositum of its field of definition with $K$ supplies the required $K'$. Thus $T(A)^{\Gamma_{K^{\rm ab}}}=T(A)\ne0$, and (E.1) fails. This also identifies the required finite extension in the failed universal quantifier. The fixed-variety criterion quoted in that source from Zarhin's Theorem 1 concerns $K$-simple factors whose CM endomorphisms are defined over $K$; the quantified assertion (E.1) is the field criterion just proved.\hfill$\square$

\paragraph{Statement AE4 (Refuted original H1; proved counterexample and replacement).}\label{res:AE4}
\passage{Section 3, manuscript pp. 30, scans 61; supplied H1 text}

The original H1 assertion that a nontrivial centralizer requires triviality of the entire cyclotomic character is false. Take $K=\mathbb Q$, $U=\mathbb P^1\setminus\{0,1,\infty\}$, the tangential section at $0$ with tangent vector $1$, and
\[
 \Delta=\operatorname{Gal}(\overline{\mathbb Q}/\mathbb Q(\mu_{3^\infty})).
\]
Then (E.1) holds, $\Delta$ is normal, $\chi|_\Delta$ is nontrivial, and
\[
 C_\pi(s(\Delta))=\mathbb Z_3(1)\subset I_0.
\]
The true statement in place of that H1 assertion is AE2, with its primewise fixed subgroup.

\emph{Proof.} The inclusion $\Gamma_{\mathbb Q(\mu_\infty)}\subset\Delta$ proves (E.1) by A1, after intersection with every $\Gamma_{K'}$. The cyclotomic Galois group over $\mathbb Q$ is the inverse limit of $(\mathbb Z/n\mathbb Z)^*$, hence is $\prod_q\mathbb Z_q^*$ by the Chinese remainder isomorphisms. The image of $\Delta$ consists precisely of tuples whose $3$-component is $1$. Every other prime component remains the full nontrivial $\mathbb Z_q^*$, whereas the $3$-component is trivial. Applying (A.2) therefore gives exactly $\mathbb Z_3(1)$, a nonzero proper subgroup of $I_0$. This verifies all the original hypothesis quantifiers while contradicting its asserted conclusion. The manuscript background is cote 149, section 3, manuscript p.~30, scan 61, together with the supplied H1 text.\hfill$\square$

% END sections/E_arithmetic_agent.tex


% BEGIN sections/E_root.tex
\subsection{The CM punctured elliptic curve and its nilpotent quotients}
\paragraph{Statement E1 (Proved; Part A(f) and the class-two part of E(c)).}\label{res:E1}
\passage{La Longue Marche, section 34, manuscript p. 245, scan 483; Section 3, manuscript pp. 27,30, scans 55–56,61–62}

Put $K=\mathbb Q(i)$, $E:y^2=x^3-x$, $U=E\setminus\{O\}$ and $\Delta=\Gamma_{K^{\rm ab}}$. Use a $K$-rational tangent at $O$ to define $s$. Let $P=\pi_1(U_{\overline K})^{(5)}$, let $\gamma_rP$ denote its closed lower central series, and put $N_r=P/\gamma_{r+1}P$. Then
\[
 N_2^{s(\Delta)}=N_2,\qquad
 N_2^{s(\Gamma_K^\circ)}=\gamma_2P/\gamma_3P\simeq\mathbb Z_5(1).
 \tag{E.2}
\]
The full profinite group has $C(s)=I_O$. Its image in $\pi^{\rm ab}=T(E)$ is zero, and its image in $N_2$ is exactly the centre $\gamma_2P/\gamma_3P$. Thus the first group in (E.2) contains non-inertia elements in this quotient. The comparison is with LM3b D.5 and Longue Marche section 34, scan 483.

The image assertion has the generality asked for in Part A(f): for every
elliptic curve $E_0$ over a number field $K_0$, every rational tangential
section at its removed origin, and every prime $\ell$, the image of its
full $C(s)$ in $T(E_0)$ is zero and its image in the class-two geometric
pro-$\ell$ quotient is the entire central $\mathbb Z_\ell(1)$.
In particular this applies to every CM elliptic curve over a number field
containing its CM field. The additional fixed-group formulas (E.2) retain
the specific curve, subgroup and prime stated above.

\begin{proof}
The topological punctured torus has free group on two generators $x,y$, with its boundary represented by $c=[x,y]$ up to inversion and conjugation. SGA~1 XII.5.1 identifies its pro-$5$ completion with $P$, so $P$ is free pro-$5$ of rank two. The quotient $N_2$ has unique coordinates $x^a y^b c^d$ with $a,b,d\in\mathbb Z_5$. Existence follows by collecting commutators, which are central; uniqueness follows by mapping to the upper triangular unipotent $3$ by $3$ matrices over $\mathbb Z_5$. Hence
\[
 1\longrightarrow Z=\mathbb Z_5c\longrightarrow N_2
 \longrightarrow V=T_5(E)\longrightarrow1,
 \qquad Z=\bigwedge^2V=\mathbb Z_5(1).
\]
The determinant action on $Z$ follows by computing the commutator of two lifted vectors; the last identification is the Weil pairing.

The endomorphism $[i]:(x,y)\mapsto(-x,iy)$ is defined over $K$ and commutes with Galois. Its action $J$ on $V$ has $J^2=-1$. The two roots of $X^2+1$ in $\mathbb Z_5$ differ by a unit. Their idempotents decompose $V$ into two free rank-one eigenspaces. Both have rank one because on the complex homology lattice multiplication by $i$ has characteristic polynomial $X^2+1$, and comparison preserves that integral action. Each Galois action on one rank-one summand is a character, hence factors through $\Gamma_K^{\rm ab}$. Therefore $\Delta$ acts trivially on $V$ and on $Z$.

We must also compute the possible shear from $V$ to $Z$ in the actual tangential action. The parameter $t=-x/y$ at $O$ satisfies $[-1]^*t=-t$ exactly. On the local pro-$5$ covers $t_n^{5^n}=t$ the maps $t_n\mapsto-t_n$ are compatible, square to the identity and commute with every coefficient automorphism. They define the Galois-invariant path between the two opposite tangential base points. Composing this path with $[-1]$ gives a genuine involution $\alpha$ of $P$ commuting with $s(\Gamma_K)$. This construction uses the odd prime $5$ and asserts no corresponding order-two path for the full profinite group.

Since $2$ is invertible in $\mathbb Z_5$, the class-two group law can be written in integral Campbell--Hausdorff coordinates on a Lie module $L$, with product $u*v=u+v+\tfrac12[u,v]$. The involution is $-1$ on $L/Z=V$ and $+1$ on $Z$. The idempotents $(1\pm\alpha)/2$ give $L=L^-\oplus Z$ and the projection $L^-\to V$ is an isomorphism. Every Galois element preserves this splitting. Its action is therefore exactly its action on $V$ together with the determinant action on $Z$, with no additional shear. This proves $N_2^{s(\Delta)}=N_2$.

For $\Gamma_K^\circ$, A1 gives $V^{\Gamma_K^\circ}=0$, while $Z$ is fixed because its character is $\chi_5$. The same splitting proves the second equality in (E.2). In the full group A6 gives $C(s)=I_O$. The boundary $[x,y]$ is zero in abelianization and generates $Z$ in $N_2$, proving both asserted images. A different rational tangent changes the action by a cocycle with values in $I_O$. Its image in $N_2$ is central, so all the class-two conclusions are independent of this change.

For the general image assertion, A6 again identifies the full centralizer
with origin inertia. The punctured-torus presentation and the collection
map into the Heisenberg group apply over $\mathbb Z_\ell$ for every
$\ell$. They identify its boundary commutator with a generator of the
central factor, and its abelian image with zero. These are the two actual
quotient maps used in the assertion; the general argument needs no
Campbell--Hausdorff coordinate change.
\end{proof}

\paragraph{Statement E2 (Proved; audit of LM3b D.5).}\label{res:E2}
\passage{La Longue Marche, section 34, manuscript p. 245, scan 483; LM3b;D.5;not a manuscript scan}

In the class-two arithmetic pro-$5$ quotient for this curve, the cyclic subgroup generated by a lift $x$ of a primitive vector in a CM eigenspace of $V$ is maximal procyclic, has arithmetic normalizer mapping onto $\Gamma_K$, and is not a cusp inertia subgroup in that quotient.

\begin{proof}
Choose a second lift $y$ completing that vector to a basis of $V$ and write $c=[x,y]$. Every lift of $\sigma\in\Gamma_K$ acts on $x$ by
\[
 x\longmapsto x^{a_\sigma}c^{d_\sigma},\qquad
 a_\sigma\in\mathbb Z_5^*,\ d_\sigma\in\mathbb Z_5.
\]
Replacing that arithmetic lift by its product with $y^{d_\sigma/a_\sigma}$ removes the central factor, because
\[
 y^r x^a c^d y^{-r}=x^a c^{d-ar}.
\]
Thus the normalizer maps onto $\Gamma_K$. The primitive nonzero image of $x$ in $V$ proves maximality: a proper overgroup that is procyclic would express $x$ as a proper $5$-power in that group, forcing its image in $V$ to be divisible by $5$. All cusp inertia here lies in $Z$, while this subgroup has nonzero image in $V$. These conclusions use the explicit quotient $\pi\twoheadrightarrow N_2$; they do not assert a lift of its normalizer to the full arithmetic fundamental group.

The program checks every unit $a$ and every $d$ modulo $5,25,125$, respectively $20$, $500$, $12500$ conjugations. It also computes the two distinct roots of $X^2+1$ at those moduli, verifying the eigenspace denominators used in E1.
\end{proof}

\paragraph{Statement E3 (Open in the full group; finite classes proved).}\label{res:E3}
\passage{La Longue Marche, section 34, manuscript p. 245, scan 483; LM3b;D.4–D.5;not a manuscript scan}

For the actual $s(\Delta)$-action above, the entire last central layer $\gamma_rP/\gamma_{r+1}P$ is fixed in $N_r$ for every $r\ge1$. For every $r\ge3$, it contains a nonidentity fixed element outside every conjugate of the image of cusp inertia. The calculations through $r=10$ give the degree ranks
\[
 2,1,2,3,6,9,18,30,56,99.
\]
They do not determine whether the full group $P^{s(\Delta)}$ contains a non-inertia element. The same unresolved issue remains for $\pi^{s(\Delta)}$ and for the character-free maximal-procyclic assertion of LM3b D.4.

\begin{proof}[Proof and exact remaining obligation]
An automorphism inducing the identity on $P/\gamma_2P$ induces the identity on every $\gamma_rP/\gamma_{r+1}P$: write generators of the latter as length-$r$ iterated commutators, replace each generator by itself times an element of $\gamma_2P$, and collect. Each error has length at least $r+1$. These commutators topologically generate the layer by the definition of the closed lower central series and the commutator identities. E1 shows that each element of $s(\Delta)$ has this property. Since the last layer is central in $N_r$, its elements are actual fixed group elements, not merely fixed vectors in a different representation.

Take $w_r=[x,[x,\ldots,[x,y]\ldots]]$ of length $r$. Under the Magnus substitution $x\mapsto1+X$, $y\mapsto1+Y$, its lowest term is $\operatorname{ad}_X^{r-1}(Y)$, whose coefficient of $X^{r-1}Y$ is $1$. Mapping into the finite group of units with constant term $1$ in
\[
 (\mathbb Z/5)[X,Y]_{\rm noncomm}/(\text{all words of degree }>r)
\]
makes every commutator of length $r+1$ trivial and keeps this coefficient nonzero. The universal property of the free pro-$5$ group gives that continuous finite quotient. This proves $w_r\ne1$ in $N_r$. Every power $c^a$ has degree-two image $a c$ in the torsion-free centre of $N_2$. Consequently the image of any conjugate of the cusp inertia intersects the last layer of $N_r$, $r\ge3$, trivially. This proves the asserted non-inertia witness for every finite class.

The program enumerates the Lyndon words in $X,Y$ through degree ten and expands their standard bracketings in the associative word algebra over $\mathbb Z$. The leading word of each expansion has coefficient one and all other terms are larger. This gives an exact triangular independence certificate. Collection using antisymmetry and the Jacobi identity expresses all Lie words in these bracketings; it is obtained inductively by factoring a Lyndon word at its longest Lyndon suffix. The counted ranks agree with $r^{-1}\sum_{d\mid r}\mu(d)2^{r/d}$. The explicit witness expansions and all ranks are in the output. The fixed-layer conclusion follows from the preceding commutator proof, rather than from treating arbitrary matrices as Galois images.

For the passage to the actual pro-$5$ graded layer, make the same
truncated-unit substitution over $\mathbb Z/5^m$ for every $m$.
The triangular leading coefficients remain one at each such level.
Consequently any $\mathbb Z_5$-linear relation between the displayed
brackets has every coefficient divisible by $5^m$ for every $m$, and
all coefficients vanish. Together with the collection argument, this
proves the stated $\mathbb Z_5$-ranks in the geometric group itself.

For the precise missing passage, the equivariant quotient $q_r:P\to N_r$ gives
\[
 q_r(P^{s(\Delta)})\subseteq N_r^{s(\Delta)},\qquad
 P^{s(\Delta)}=\varprojlim_r N_r^{s(\Delta)}.
 \tag{E.3}
\]
The second equality uses $P=\varprojlim_rN_r$: every finite $5$-group is nilpotent, so every finite quotient of $P$ kills some $\gamma_rP$, and compactness proves both injectivity and surjectivity of this inverse-limit identification. The explicit $w_r$ maps to $1$ in $N_{r-1}$. Thus these particular fixed elements do not form a nontrivial compatible element of (E.3). Already at the transition $N_3\to N_2$, lifting a specified noncentral fixed $v\in N_2$ requires vanishing of the continuous cocycle
\[
 \delta\longmapsto s(\delta)(\widetilde v)\widetilde v^{-1}
 \ \in\gamma_3P/\gamma_4P.
\]
It is independent of the chosen lift, because this kernel is fixed pointwise. Its vanishing for all $\delta\in\Delta$ is exactly the missing lift condition at this first transition; success at it still leaves the compatible conditions at every higher transition. This paper computes neither that arithmetic cocycle nor a compatible non-inertia tower. In the full profinite group there is the additional requirement of compatible lifts through its finite quotients beyond the pro-$5$ quotient. These are the explicit remaining obligations; the quotient counterexample E2 does not dispose of them.
\end{proof}

% END sections/E_root.tex


% BEGIN sections/E_local_agent.tex
\subsection{The local-field comparison: proved quotient calculations}

\paragraph{Statement D14 (Proved).}\label{res:D14}
\passage{section 3; pp.27,30; scans 55-56,62; local-field extension requested in E(d)}

Let $K/\mathbb Q_p$ be finite, put
$\Delta=G_{K(\mu_\infty)}$, and write $\pi^{(p')}$ for the maximal
prime-to-$p$ geometric quotient.  Suppose the hyperbolic curve
$U=X\setminus S$ extends to
a smooth proper curve $\mathscr X/\mathcal O_K$ with boundary
$\mathscr S$ finite \'etale over $\mathcal O_K$.
For a point $x\in(\mathscr X\setminus\mathscr S)(\mathcal O_K)$,
\[
       C_{\pi^{(p')}}\bigl(s_x(\Delta)\bigr)=\pi^{(p')}.
\]
For a rational boundary section $c$ choose an integral nonzero tangent
vector $v$, whose reduction is nonzero, and the associated tangential
section $s_v$.  It satisfies the same equality.  For the cuspidal
section $s_a$ obtained from $s_v$ by a Kummer class
$a\in\widehat{K^*}$, let
$b=(b_\ell)_\ell\in\widehat{\mathbb Z}$ be its valuation.  Then
\[
 C_{\pi^{(p')}}\bigl(s_a(\Delta)\bigr)=
 \begin{cases}
       \pi^{(p')},& b_\ell=0\text{ for every }\ell\ne p,\\
       I_c^{(p')},& b_\ell\ne0\text{ for at least one }\ell\ne p.
 \end{cases}                                                   \tag{E.4}
\]
These are fixed groups in the quotient, computed using the induced
actual section action.

\begin{proof}
All finite unramified extensions of $K$ are contained in
$K(\mu_\infty)$: roots of unity of order $q^n-1$, where $q$ is the
residue-field cardinality, give the unramified extension of degree $n$.
Thus $\Delta\subset I_K$.

The specialization input is SGA~1, XIII, Corollary 2.9 and the
construction in 2.10 \cite{D_SGA1}: for a smooth proper morphism with
a relative normal-crossings boundary, the categories of torsors of
order prime to the residue characteristic specialize equivalently.
In the proper case the resulting isomorphism on maximal prime-to-$p$
fundamental groups is also Stacks, Theorem 58.30.3, Tag 0C0R.
Use the section $x$ to identify the fibre functors on the generic and
special fibres.  The resulting based action factors through the
Galois group of the residue field.  Its restriction to $I_K$, hence to
$\Delta$, is the identity automorphism, proving the first equality.
For the integral tangent vector the same comparison uses the tame
local fibre functor: in a parameter with unit tangent, its finite
prime-to-$p$ coverings are obtained by extracting roots of that
parameter of order prime to $p$.  The roots of unity needed for these
coverings are unramified constants.  Specialization therefore
identifies this pointed fibre functor as well.  The local inertia
comparison is explicitly SGA~1, XIII, Lemma 2.11.  It gives the stated
identity action for $s_v$.

Let $\kappa_a:G_K\to\widehat{\mathbb Z}(1)$ be the Kummer cocycle of
the twist.  Since $s_v$ acts trivially on $\pi^{(p')}$ over $\Delta$
and $\chi|_\Delta=1$, the action of $s_a(\sigma)$ for $\sigma\in\Delta$
is the inner automorphism by $\kappa_a^{(p')}(\sigma)\in I_c^{(p')}$.
Every unit of $\mathcal O_K$ has an $n$-th root in $K^{\rm nr}$ when
$n$ is prime to $p$: take a root in the algebraic residue field and
use Hensel's lemma, whose derivative is a unit.  The compatible
division-point sets give the same assertion for the prime-to-$p$
Kummer class of a completed unit.  Consequently on $\Delta$
\[
       \kappa_{a,\ell}=b_\ell\kappa_{\varpi,\ell}\qquad(\ell\ne p),
\]
where $\varpi$ is a uniformizer.
Each $\kappa_{\varpi,\ell}|_\Delta$ is nonzero.  Indeed, the
ramification indices of $K(\mu_\infty)/K$ have bounded
$\ell$-adic valuation for $\ell\ne p$, whereas the value of an
$\ell^n$-th root of $\varpi$ would be $1/\ell^n$ in the
$K$-normalized valuation.  Vanishing of that Kummer class would put
all these roots in $K(\mu_\infty)$, contradicting the bounded
denominators.  Its image is a nonzero closed subgroup of
$\mathbb Z_\ell(1)$, so multiplication by a nonzero $b_\ell$ remains
nonzero.

We record why the geometric centralizer calculation applies to this
quotient.  Finite groups of order prime to $p$ form a class closed
under subgroups, quotients and extensions.  Nakamura
\cite[Lemma (2.1.2)]{Nakamura1994} states for every such full class that
a nonidentity element of a free-generator subgroup has precisely that
subgroup as centralizer.  For at least two punctures a chosen peripheral
loop is part of a free basis.  For one puncture, pass to the compact
curve's mod-$\ell$ homology cover, $\ell\ne p$, as in the last
paragraph of Nakamura's Lemma (2.1.4); it has at least two punctures.
Equivalently, the boundary-separation proof of A2 works over
$\prod_{\ell\ne p}\mathbb Z_\ell$: finite normal covers separate two
cusps and can have at least three boundary points, whose distinct
summand lines in the abelianization intersect trivially.  Nonzero
inertia elements have nonzero positive integral powers in the open
cover group.  Thus distinct inertia groups in the prime-to-$p$
universal cover meet trivially, and the centralizer of every
nonidentity element of $I_c^{(p')}$ is $I_c^{(p')}$.
The image of $\kappa_a^{(p')}$ is either zero or contains such an
element.  This proves both alternatives of (E.4).
\end{proof}

\paragraph{Statement D15 (Open; full and bad-reduction
computations unfinished).}\label{res:D15}
\passage{section 3; pp.27,30; scans 55-56,62; local-field extension requested in E(d)}

Part E(d) has the good-marked-reduction, prime-to-$p$ calculations D14.
For instance, over $\mathbb Q_5$ on $U_{0,3}$, the section of the
integral point $2$ and the integral tangential section at $0$ have all
of the free prime-to-$5$ group of rank two fixed.  The tangent twist by
$a=5$ has exactly $I_0^{(5')}$ fixed in that quotient.  Hence the
prime-to-$p$ answer already depends on the specified cuspidal section.

The full group is compared with the computed quotient by
$q:\pi\twoheadrightarrow\pi^{(p')}$, with kernel $N$.
There is always the inclusion
\[
       q\bigl(\pi^{s(\Delta)}\bigr)
                    \subseteq(\pi^{(p')})^{s(\Delta)}.          \tag{E.5}
\]
For a fixed quotient element $\bar h$ and a lift $h\in\pi$, the
elements
\[
           d_\sigma=h^{-1}s(\sigma)h s(\sigma)^{-1}\in N
                 \qquad(\sigma\in\Delta)
\]
form a continuous nonabelian $1$-cocycle for the action of $\Delta$ on
$N$.  The element $\bar h$ belongs to the left side of (E.5) exactly
when this cocycle is a coboundary.  Indeed, replacing $h$ by $hn$ makes
it fixed exactly when $d_\sigma=n\,\sigma(n)^{-1}$ for every $\sigma$.
This proves the precise obstruction to passing from D14 to a full-group
calculation.

The smallest retained full-group instance is $U_{0,3}/\mathbb Q_5$
with $s=s_2$, or with the integral tangent at $0$.  The classes just
displayed have not been computed for all quotient elements there.
A retained bad-marked-reduction instance is
$\mathbb P^1\setminus\{0,1,5,\infty\}/\mathbb Q_5$ with the point $2$;
the two boundary sections $0$ and $5$ meet in the special fibre, so
D14's hypothesis fails.  The full and prime-to-$5$ fixed groups for
this instance have not been computed here.

For the requested literature comparison, Mochizuki
\cite[Lemma 1.3.9]{D_Mochizuki2004} assumes two hyperbolic curves over
finite extensions of $p_i$-adic local fields and an isomorphism of
their full arithmetic fundamental groups.  It proves equality of
their types and preservation of cusp inertia subgroups, both in the
full geometric groups and in every pro-$\ell$ quotient.  Its comparison
map is that arithmetic isomorphism.  D14 instead uses the action of a
specified section restricted to $\Delta$; (E.5) gives the available
map between its full and prime-to-$p$ fixed groups.  No full-group
fixed-group formula is inferred from Lemma 1.3.9.  This identifies the
remaining local calculation rather than suppressing it.

% END sections/E_local_agent.tex


% BEGIN sections/E4_prior_hunches.tex
\paragraph{Statement E4 (Open for the degeneration conclusion; proved
normalization and branch-sign assertions).}\label{res:E4}
\passage{Supplied war-game excerpts of 12 November 2026, entries R4 and H2.
Their references to scans 101--103 have not been independently collated here.}
The normalization assertion H2(a) in the supplied
\texttt{R149\_WARGAME\_EXCERPTS\_20261112.md} is proved in its precise
arithmetic-normalizer formulation: for a hyperbolic curve $C$ over a
number field, a nontrivial closed procyclic subgroup
$I\subset\pi_1(C_{\overline k})$ is essentially cyclotomically
normalized exactly when it is contained in cusp inertia. Thus no such
$I$ exists for a proper hyperbolic curve. Here ``essentially'' means
that a closed subgroup of $N_{\pi_1(C)}(I)$ maps onto an open subgroup
of $G_k$ and acts on $I$ by the cyclotomic character. The same assertion
holds over fields finitely generated over $\mathbf Q$ by AE1.
At a split node with local equation $xy=t$, exchanging the branches
acts on local geometric annulus inertia by inversion. The full
degeneration assertion H2(b), and hence its proposed completion of R4,
is not established in this work and is not used as an input.

\begin{proof}
Stix~\cite[Theorem 20 and \S5.2]{B_Stix2008} proves the stated
equivalence for a characteristic-zero field satisfying
$\operatorname{Hom}_{G_k}(T(S),T(A))=0$ for every torus $S$ and
abelian variety $A$. Part A verifies that hypothesis over number
fields. For the stated finitely generated extension, split $S$ over a
finite extension $k'/k$. A $G_{k'}$-equivariant map from
$T(S)=\widehat{\mathbf Z}(1)^r$ to $T(A)$ has image fixed by
$G_{k'}^\circ$; AE1 makes that image zero. Stix's theorem therefore
applies. In the proper case the set of cusp inertia subgroups is
empty; in the affine case they give exactly the asserted subgroups.
This verifies H2(a) through its actual normalizer map, without choosing
an unproved splitting of that normalizer.

For the branch calculation choose compatible $n$th roots $t_n$ of
$t$ over the geometric generic field, and use the Kummer cover
$z^n=x$. The exchange $x\leftrightarrow y=t/x$ lifts to
$z\mapsto t_n/z$. Conjugating the deck transformation
$z\mapsto\zeta_n z$ by this lift gives
$z\mapsto\zeta_n^{-1}z$. Passing over all $n$ proves the inversion
map on $\widehat{\mathbf Z}(1)$. An arithmetic action additionally
acts on $\zeta_n$ by its cyclotomic character. Thus, when it exchanges
the two branches, its local formula contains that inversion sign as
well. After a finite extension fixing a node and its two branches,
the permutation and branch-exchange issue is removed. This is a
calculation in the local group; it does not assert that the required
global inertia embedding or its arithmetic equivariance has already
been established.
\end{proof}

For a precise record of the remaining source obligation, the opened
Hoshi--Mochizuki manuscript
\cite[Definition 2.10(iv) and Lemma 2.11(ii),(iii)]{B_HoshiMochizuki2011}
gives a local log-group injection whose map on the two node coordinates
is $(a,b)\mapsto a+b$. Write
$I_e^{\log}:=\ker[T^2\xrightarrow{+}T]=\{(a,-a)\}$ for its
geometric node kernel, embedded in the geometric group $\Pi_{G_s}$ of
the associated semigraph. To finish H2(b) in the stated generality one must
supply the log-specialization comparison to the \emph{full} geometric
generic-fibre group and check the displayed maps in the chain
\[
 \widehat{\mathbf Z}(1)
 \xrightarrow{a\mapsto(a,-a)} I_e^{\log}
 \lhook\joinrel\longrightarrow\Pi_{G_s}
 \xrightarrow{\ \mathrm{required\ comparison}\ }
 \pi_1(C_{\overline\eta}),
\]
where the final arrow denotes the required inverse specialization
comparison, not a map constructed in this work. Its obligations include
injectivity, compatibility with the residue-field action,
and the action obtained after \emph{every} splitting of the trait
Galois extension. One must also control the finite permutation of nodes
and branches before asserting individual invariance. The cited local
lemma alone does not provide these entire comparisons. The opened
characteristic-zero Stix theorem also does not by itself prove a
full-group assertion over a finite residue field, where a prime-to-$p$
comparison leaves the $p$-primary part outside its scope.

Accordingly R4(i)'s weight argument and R4(ii)'s proposed passage from
Tate modules to polarized abelian varieties are recorded as unverified
arguments in this deliverable, rather than used as proved inputs. The
smallest retained H2 test is a genus-two degeneration whose special
fibre is two isomorphic CM elliptic curves meeting at one separating
node. Under a surface marking the separating cycle is a commutator
$[a_1,b_1]$, whose image under the actual abelianization map
$\pi\to\pi^{\mathrm{ab}}$ is zero. Thus the abelianized calculation
does not detect that candidate cycle; its global cyclotomic
normalization is exactly the additional assertion requiring the
comparisons above. This status refers to the attached R4 and H2
entries; no independent reading of scans 101--103 is claimed here.

% END sections/E4_prior_hunches.tex

\nocite{Cote149,LongueMarche,GrothendieckFaltings1983}
\bibliographystyle{alpha}
\bibliography{CYCLOTOMIC_KERNEL_references}
\end{document}
